% Smartphones et cybercafés (avantages et inconvénients des médias) — Allemand 1ère A — Cours signé OnBuch+
% Programme officiel MINESEC. Compiler avec : tectonic ALA15-smartphones-et-cybercafes-avantages-et-inconvenients-des-med.tex
\newcommand{\DOCMATIERE}{Allemand}
\newcommand{\DOCNIVEAU}{1ère A}
\newcommand{\DOCMODULE}{Chapitre 4 : Média et communication}
\newcommand{\DOCLECON}{ALA15}
\newcommand{\DOCTITRE}{Smartphones et cybercafés (avantages et inconvénients des médias)}
% OnBuch+ — préambule « cours signés » ENRICHI (compilé avec tectonic / XeLaTeX)
% Système visuel « Pop » d'OnBuch+ : fond crème (jamais blanc), bordures encre
% épaisses, ombres « dures » décalées, titres Archivo Black en capitales.
% Version MATHÉMATIQUES : graphes (pgfplots/tikz), boîtes pédagogiques
% (définition, propriété, méthode, attention, le savais-tu, exercice résolu,
% exercice, TP, bilan), page de garde et signature OnBuch+ ancrée en bas.
%
% Macros attendues dans main.tex AVANT \input{preamble} :
%   \DOCMATIERE  \DOCNIVEAU  \DOCMODULE  \DOCLECON (numéro)  \DOCTITRE
\documentclass[11pt]{article}

\usepackage{fontspec}
\usepackage[ngerman,french]{babel} % français = langue principale ; l'allemand via \all{..} / textebox
\usepackage[a4paper,margin=2cm,top=3.4cm,bottom=2.8cm,headheight=1.4cm,footskip=1.3cm]{geometry}
\usepackage{amsmath,amssymb,amsfonts}
\usepackage{mathtools} % \xmapsto, \coloneqq... souvent utilisés par les modèles
\usepackage{cancel} % simplifications barrées (bilans de forces)
% \abs{z} (module, valeur absolue) et \norm{u} (norme), utilisés par les modèles
\providecommand{\abs}{}\let\abs\relax\DeclarePairedDelimiter{\abs}{\lvert}{\rvert}
\providecommand{\norm}{}\let\norm\relax\DeclarePairedDelimiter{\norm}{\lVert}{\rVert}
\usepackage[table]{xcolor} % [table] : fournit aussi \rowcolors (alternance de lignes)
\usepackage{pagecolor}
\usepackage{tikz}
\usetikzlibrary{arrows.meta,positioning,calc,shapes.geometric,shapes.misc,shapes.arrows,decorations.pathmorphing,decorations.pathreplacing,decorations.markings,patterns,shadows,mindmap,trees,fit,backgrounds,matrix,intersections,angles,quotes}
\usepackage{pgfplots}
\usepackage[version=4]{mhchem} % équations nucléaires \ce{^{235}_{92}U}
\usepackage[european,siunitx]{circuitikz} % schémas électriques
\pgfplotsset{compat=1.18}
\usepgfplotslibrary{fillbetween,groupplots} % aires entre courbes (calcul intégral)
\usepackage{enumitem}
\usepackage{fancyhdr}
% [most] charge toutes les bibliothèques tcolorbox (listings, minted, raster,
% documentation...) et gonfle inutilement la mémoire TeX sur un document long
% (~300 boîtes) : on ne charge que ce qui est réellement utilisé.
\usepackage[skins,breakable]{tcolorbox}
\usepackage{graphicx}
\usepackage{colortbl}
\usepackage{booktabs}
\usepackage{tabularx}
\usepackage{diagbox} % case d'en-tête diagonale des tableaux à double entrée
% Colonnes extensibles centrée / gauche / droite (C, L, R), souvent utilisées par les modèles
\newcolumntype{C}{>{\centering\arraybackslash}X}
\newcolumntype{L}{>{\raggedright\arraybackslash}X}
\newcolumntype{R}{>{\raggedleft\arraybackslash}X}
\usepackage{array}
\usepackage{multicol}
\usepackage{multirow}
\usepackage{siunitx}
% parse-numbers=false : accepte aussi \SI{2,5 \times 10^{-3}}{...}
\sisetup{locale=FR,output-decimal-marker={,},group-minimum-digits=4,parse-numbers=false}
\DeclareSIUnit\ohm{\ensuremath{\symup{\Omega}}} % U+2126 absent de la police
\DeclareSIUnit\division{div} % réglages d'oscilloscope (ms/div, V/div)
\DeclareSIUnit\tr{tr} % tours (vitesse de rotation, ex. \SI{1500}{\tr\per\minute})
\DeclareSIUnit\molar{M} % concentration molaire (ex. \SI{3}{\molar}, \SI{1}{\micro\molar})
\DeclareSIUnit\an{an} % année (le modèle confond parfois avec \year, un compteur TeX) — ex. \SI{5}{\kilo\meter\per\an}
\DeclareSIUnit\FCFA{FCFA} % franc CFA (ex. \SI{500}{\FCFA\per\kilo\gram})
\DeclareSIUnit\degreeBrix{\textdegree Brix} % teneur en sucre d'un sirop/jus (réfractométrie)
\DeclareSIUnit\month{mois} % le modèle utilise parfois \month, un compteur TeX, comme unité de durée
\usepackage{qrcode}
% \degree : commande fréquemment utilisée par les modèles pour un angle ou une
% température en dehors de \SI{}{} (non fournie par les paquets chargés).
\providecommand{\degree}{^{\circ}}
% \qed (fin de démonstration), utilisé par les modèles sans amsthm
\providecommand{\qed}{\hfill$\square$}
% \barre : double barre des valeurs interdites dans un tableau de variations (tkz-tab)
\providecommand{\barre}{\ensuremath{\|}}
% environnement proof (amsthm non chargé) utilisé par les modèles
\ifdefined\proof\else\newenvironment{proof}[1][Démonstration]{\par\noindent\textit{#1.}\ }{\qed\par}\fi
\usepackage{lastpage}
\usepackage{hyperref}
\hypersetup{hidelinks}

% ============================================================
% Palette « Pop » OnBuch+ — mêmes hex que la classe P (plus_ui.dart)
% ============================================================
\definecolor{poporange}{HTML}{F59321}
\definecolor{popdark}{HTML}{E07A0C}
\definecolor{popcream}{HTML}{FFF6E8}
\definecolor{popink}{HTML}{1C1714}
\definecolor{poppurple}{HTML}{7A5AE0}
\definecolor{popgreen}{HTML}{2E9E6A}
\definecolor{poppink}{HTML}{E0457B}
\definecolor{popblue}{HTML}{2F7DE1}
\definecolor{popgold}{HTML}{C98A12}
\definecolor{popmuted}{HTML}{8A7F76}
\definecolor{popline}{HTML}{EADFD2}
\definecolor{popgray}{HTML}{8A7F76}
\colorlet{popgrey}{popgray}
% Alias tolérants (le modèle nomme parfois les couleurs différemment)
\colorlet{popwhite}{white}
\colorlet{popblack}{popink}
\colorlet{popred}{poppink}
\colorlet{popyellow}{popgold}
% Teintes claires (fonds de tableaux/encadrés) — le modèle nomme parfois
% les couleurs avec un suffixe L (ex. popblueL) sans les avoir définies.
\colorlet{poporangeL}{poporange!20}
\colorlet{popdarkL}{popdark!20}
\colorlet{poppurpleL}{poppurple!20}
\colorlet{popgreenL}{popgreen!20}
\colorlet{poppinkL}{poppink!20}
\colorlet{popblueL}{popblue!20}
\colorlet{popgoldL}{popgold!20}
\colorlet{popredL}{popred!20}
\colorlet{popgrayL}{popgray!20}
% Teintes claires pour fonds de tableaux / zones
\colorlet{popblueL}{popblue!10!white}
\colorlet{poporangeL}{poporange!14!white}
\colorlet{popgreenL}{popgreen!12!white}
\colorlet{poppurpleL}{poppurple!12!white}
\colorlet{poppinkL}{poppink!12!white}
\colorlet{popgoldL}{popgold!14!white}

\pagecolor{popcream}

% ============================================================
% Typographie
% ============================================================
\defaultfontfeatures{Ligatures=TeX}
\setmainfont{PlusJakartaSans-Medium.ttf}[Path=../../../fonts/,BoldFont=PlusJakartaSans-Bold.ttf]
% Maths en sans-serif (Fira Math) avec chiffres et lettres droites de Plus
% Jakarta, pour que les formules se fondent dans le texte.
\usepackage[math-style=ISO,bold-style=ISO]{unicode-math}
\setmathfont{FiraMath-Regular.otf}[Path=../../../fonts/]
\setmathfont{PlusJakartaSans-Medium.ttf}[Path=../../../fonts/,range={up/num,up/{Latin,latin},\mathup}]
\usepackage{icomma} % virgule décimale sans espace en mode math
% Caractères mathématiques Unicode absents des polices de texte (Plus Jakarta,
% Archivo Black) : rendus par leur équivalent mathématique, en texte comme en
% formule — sinon ils s'affichent en carré vide (ex. « Arithmétique dans ℤ »).
\usepackage{newunicodechar}
\newunicodechar{ℝ}{\ensuremath{\mathbb{R}}}
\newunicodechar{ℤ}{\ensuremath{\mathbb{Z}}}
\newunicodechar{ℂ}{\ensuremath{\mathbb{C}}}
\newunicodechar{ℕ}{\ensuremath{\mathbb{N}}}
\newunicodechar{ℚ}{\ensuremath{\mathbb{Q}}}
\newunicodechar{𝔻}{\ensuremath{\mathbb{D}}}
\newunicodechar{α}{\ensuremath{\alpha}}
\newunicodechar{β}{\ensuremath{\beta}}
\newunicodechar{γ}{\ensuremath{\gamma}}
\newunicodechar{δ}{\ensuremath{\delta}}
\newunicodechar{ε}{\ensuremath{\varepsilon}}
\newunicodechar{θ}{\ensuremath{\theta}}
\newunicodechar{λ}{\ensuremath{\lambda}}
\newunicodechar{μ}{\ensuremath{\mu}}
\newunicodechar{σ}{\ensuremath{\sigma}}
\newunicodechar{τ}{\ensuremath{\tau}}
\newunicodechar{φ}{\ensuremath{\varphi}}
\newunicodechar{ψ}{\ensuremath{\psi}}
\newunicodechar{ω}{\ensuremath{\omega}}
\newunicodechar{Σ}{\ensuremath{\Sigma}}
% majuscules produites par \MakeUppercase dans les titres (α-aminés -> Α)
\newunicodechar{Α}{\ensuremath{\alpha}}
\newunicodechar{Β}{\ensuremath{\beta}}
\newunicodechar{Γ}{\ensuremath{\Gamma}}
\newunicodechar{Δ}{\ensuremath{\Delta}}
\newunicodechar{Ε}{\ensuremath{\varepsilon}}
\newunicodechar{Θ}{\ensuremath{\Theta}}
\newunicodechar{Λ}{\ensuremath{\Lambda}}
\newunicodechar{Μ}{\ensuremath{\mu}}
\newunicodechar{Τ}{\ensuremath{\tau}}
\newunicodechar{Φ}{\ensuremath{\Phi}}
\newunicodechar{Ψ}{\ensuremath{\Psi}}
\newunicodechar{Ω}{\ensuremath{\Omega}}
\newunicodechar{↦}{\ensuremath{\mapsto}}
\newunicodechar{∈}{\ensuremath{\in}}
\newunicodechar{∉}{\ensuremath{\notin}}
\newunicodechar{∧}{\ensuremath{\wedge}}
\newunicodechar{⇔}{\ensuremath{\Leftrightarrow}}
\newunicodechar{⟺}{\ensuremath{\Longleftrightarrow}}
\newunicodechar{⇒}{\ensuremath{\Rightarrow}}
\newunicodechar{⟹}{\ensuremath{\Longrightarrow}}
\newunicodechar{∘}{\ensuremath{\circ}}
\newunicodechar{∪}{\ensuremath{\cup}}
\newunicodechar{∩}{\ensuremath{\cap}}
\newunicodechar{≡}{\ensuremath{\equiv}}
\newunicodechar{⊂}{\ensuremath{\subset}}
\newunicodechar{⊥}{\ensuremath{\perp}}
\newunicodechar{∃}{\ensuremath{\exists}}
\newunicodechar{∀}{\ensuremath{\forall}}
\newunicodechar{∛}{\ensuremath{\sqrt[3]{\,}}}
\newunicodechar{∜}{\ensuremath{\sqrt[4]{\,}}}
\newunicodechar{⃗}{\ensuremath{{}^{\to}}}
\newunicodechar{ⁿ}{\ifmmode^{n}\else\textsuperscript{\ensuremath{n}}\fi}
\newunicodechar{ˣ}{\ifmmode^{x}\else\textsuperscript{\ensuremath{x}}\fi}
\newunicodechar{ᵇ}{\ifmmode^{b}\else\textsuperscript{\ensuremath{b}}\fi}
\newunicodechar{ʳ}{\ifmmode^{r}\else\textsuperscript{\ensuremath{r}}\fi}
\newunicodechar{ᵘ}{\ifmmode^{u}\else\textsuperscript{\ensuremath{u}}\fi}
\newunicodechar{ʸ}{\ifmmode^{y}\else\textsuperscript{\ensuremath{y}}\fi}
\newunicodechar{ᵃ}{\ifmmode^{a}\else\textsuperscript{\ensuremath{a}}\fi}
\newunicodechar{ᵉ}{\ifmmode^{e}\else\textsuperscript{\ensuremath{e}}\fi}
\newunicodechar{ᵏ}{\ifmmode^{k}\else\textsuperscript{\ensuremath{k}}\fi}
\newunicodechar{ᵖ}{\ifmmode^{p}\else\textsuperscript{\ensuremath{p}}\fi}
\newunicodechar{ᴺ}{\ifmmode^{N}\else\textsuperscript{\ensuremath{N}}\fi}
\newunicodechar{ᶜ}{\ifmmode^{c}\else\textsuperscript{\ensuremath{c}}\fi}
\newunicodechar{ʰ}{\ifmmode^{h}\else\textsuperscript{\ensuremath{h}}\fi}
\newunicodechar{ᵠ}{\ifmmode^{q}\else\textsuperscript{\ensuremath{q}}\fi}
\newunicodechar{ₙ}{\ifmmode_{n}\else\textsubscript{\ensuremath{n}}\fi}
\newunicodechar{ₐ}{\ifmmode_{a}\else\textsubscript{\ensuremath{a}}\fi}
\newunicodechar{ᵢ}{\ifmmode_{i}\else\textsubscript{\ensuremath{i}}\fi}
\newunicodechar{ᵣ}{\ifmmode_{r}\else\textsubscript{\ensuremath{r}}\fi}
\newunicodechar{ₖ}{\ifmmode_{k}\else\textsubscript{\ensuremath{k}}\fi}
\newunicodechar{ᵦ}{\ifmmode_{\beta}\else\textsubscript{\ensuremath{\beta}}\fi}
\newunicodechar{ₓ}{\ifmmode_{x}\else\textsubscript{\ensuremath{x}}\fi}
\newfontfamily\popdisplay{ArchivoBlack-Regular.ttf}[Path=../../../fonts/]
\newfontfamily\poplogo{SpaceGrotesk-Bold.ttf}[Path=../../../fonts/]
\newfontfamily\popsemi{PlusJakartaSans-SemiBold.ttf}[Path=../../../fonts/]
\newfontfamily\popxbold{PlusJakartaSans-ExtraBold.ttf}[Path=../../../fonts/]
\newfontfamily\popxboldls{PlusJakartaSans-ExtraBold.ttf}[Path=../../../fonts/,LetterSpace=3.0]

\setlist[enumerate]{leftmargin=1.7em,itemsep=5pt,topsep=4pt}
\setlist[itemize]{leftmargin=1.7em,itemsep=4pt,topsep=4pt,label={\color{poporange}\textbullet}}
\setlength{\parindent}{0pt}
\setlength{\parskip}{5pt}
\renewcommand{\arraystretch}{1.25}

% pgfplots : style Pop par défaut
\pgfplotsset{
  every axis/.append style={axis line style={popink,line width=1pt},tick style={popink},
    grid style={popline,line width=0.5pt},label style={font=\small},tick label style={font=\footnotesize}},
  popcurve/.style={poporange,line width=1.8pt,no markers,smooth},
}
% Même style disponible dans un \draw TikZ simple (hors pgfplots)
\tikzset{popcurve/.style={poporange,line width=1.8pt}}
\tikzset{popgrid/.style={popline,very thin}}
\colorlet{popcurve}{poporange} % le nom du style est parfois utilisé comme couleur

% ============================================================
% Pill / squiggle
% ============================================================
\newcommand{\poppill}[3]{%
  \tikz[baseline=-0.55ex]\node[draw=popink,line width=1.2pt,rounded corners=2.6pt,%
    fill=#2,inner xsep=6.5pt,inner ysep=3pt]{%
    \popxboldls\fontsize{7}{7}\selectfont\color{#3}\MakeUppercase{#1}};%
}
\newcommand{\popsquiggle}[1]{%
  \tikz[baseline=-0.4ex]{
    \draw[#1,line width=2.6pt,line cap=round]
      plot [smooth] coordinates {(0,0.09) (0.34,0.01) (0.68,0.21) (1.02,0.01) (1.36,0.10)};
  }%
}
% Mot-clé surligné (marqueur orange sous le texte)
\newcommand{\cle}[1]{{\popxbold\color{popdark}#1}}

% ============================================================
% En-tête / pied
% ============================================================
\pagestyle{fancy}
\fancyhf{}
\renewcommand{\headrulewidth}{0pt}
\renewcommand{\footrulewidth}{0pt}
\lhead{\raisebox{-0.15ex}{\poplogo\fontsize{13}{13}\selectfont\color{popink}OnBuch\color{poporange}+}}
\rhead{\poppill{\DOCMATIERE\ \textbullet\ \DOCNIVEAU\ \textbullet\ Leçon \DOCLECON}{white}{popink}}
\lfoot{\popxbold\fontsize{7.6}{7.6}\selectfont\color{popmuted}\MakeUppercase{Cours signé OnBuch+}\ \textbullet\ {\popsemi\DOCTITRE}}
\rfoot{\tikz[baseline=-0.6ex]\node[rounded corners=8.5pt,fill=popink,minimum height=17pt,minimum width=17pt,inner xsep=5pt,inner sep=0]{\popxbold\fontsize{8}{8}\selectfont\color{popcream}\thepage\,/\,\pageref*{LastPage}};}

% ============================================================
% Titres
% ============================================================
\newcommand{\chaptitre}[2]{%
  \begin{tcolorbox}[enhanced,colback=poporange,colframe=popink,boxrule=2.6pt,arc=11pt,
      drop shadow=popink,left=15pt,right=15pt,top=12pt,bottom=11pt,before skip=3pt,after skip=18pt]
    \poppill{#2}{popink}{popcream}\par\vspace{9pt}
    {\raggedright\hyphenpenalty=10000\exhyphenpenalty=10000\popdisplay\fontsize{22}{24}\selectfont\color{popcream}\MakeUppercase{#1}\par}
  \end{tcolorbox}
}
% Section numérotée : pastille orange avec le numéro + titre Archivo
\newcounter{popsec}
\newcommand{\coursec}[1]{%
  \stepcounter{popsec}\setcounter{popsub}{0}%
  \par\vspace{16pt}\noindent
  \begin{minipage}{\linewidth}
  \tikz[baseline=-0.7ex]\node[draw=popink,line width=1.6pt,fill=poporange,rounded corners=4pt,minimum size=22pt,inner sep=0]{\popdisplay\fontsize{12}{12}\selectfont\color{popcream}\thepopsec};%
  \hspace{8pt}{\popdisplay\fontsize{15}{17}\selectfont\color{popink}\MakeUppercase{#1}}\par
  \vspace{4pt}\noindent{\color{poporange}\rule{36pt}{3.6pt}}
  \end{minipage}\par\nopagebreak\vspace{8pt}
}
\newcounter{popsub}
\newcommand{\courssub}[1]{%
  \stepcounter{popsub}%
  \par\vspace{9pt}\noindent{\popxbold\fontsize{11.5}{13}\selectfont\color{popink}{\color{poporange}\thepopsec.\thepopsub}\hspace{6pt}#1}\par\nopagebreak\vspace{4pt}
}

% ============================================================
% Boîtes pédagogiques — même recette : fond blanc, bordure épaisse colorée,
% ombre dure encre, badge plein. Titre optionnel [#1].
% ============================================================
\tcbset{popbase/.style={breakable,enhanced,colback=white,boxrule=2.3pt,arc=9pt,
  shadow={3pt}{-3pt}{0pt}{popink},left=14pt,right=14pt,top=11pt,bottom=11pt,before skip=11pt,after skip=15pt,
  fontupper=\popsemi\fontsize{10.6}{14.5}\selectfont\color{popink},
  fontlower=\popsemi\fontsize{10.6}{14.5}\selectfont\color{popink}}}
% Coche dessinée (le glyphe ✓ n'existe pas dans les polices du document)
\newcommand{\popcheck}[1]{\tikz[baseline=-0.5ex]\draw[#1,line width=1.6pt,line cap=round,line join=round](-0.13,0.0)--(-0.04,-0.09)--(0.13,0.1);}
\providecommand{\coche}{\popcheck{popgreen}}
\providecommand{\resultat}[1]{\par\smallskip\noindent\fcolorbox{popgreen}{popgreenL}{\parbox{\dimexpr\linewidth-2\fboxsep-2\fboxrule\relax}{\textbf{Résultat.} #1}}\par\smallskip}
\newcommand{\popboxhead}[3]{\poppill{#1}{#2}{white}\ifx\relax#3\relax\else\hspace{6pt}{\popxbold\fontsize{10.6}{12}\selectfont\color{popink}#3}\fi\par\vspace{7pt}}

\newtcolorbox{aretenir}[1][]{popbase,colframe=popgreen,before upper={\popboxhead{À retenir}{popgreen}{#1}}}
% Texte en allemand (document de lecture, dialogue, extrait) : cadre
% d'étude. Un titre optionnel (source, auteur) s'affiche dans l'en-tête.
\newtcolorbox{textebox}[1][]{popbase,colframe=popblue,before upper={\popboxhead{Text}{popblue}{#1}\begin{otherlanguage}{ngerman}},after upper={\end{otherlanguage}}}
% Marques ✓ / ✗ (forme correcte / fautive) souvent utilisées par les modèles
\providecommand{\cross}{\textcolor{poppink}{\ensuremath{\times}}}
\providecommand{\xmark}{\cross}\providecommand{\crossmark}{\cross}\providecommand{\cmark}{\ensuremath{\checkmark}}
% Ligne de réponse / justification à compléter (inventée par les modèles)
\providecommand{\justification}[1]{\par\noindent\textit{Justification~:} \dotfill\par}
% \textipa (package tipa non chargé) : la police ne couvre pas l'API -> simple passage
\providecommand{\textipa}[1]{#1}
% Soulignement (ulem non chargé) : \uline, \ul
\providecommand{\uline}[1]{\underline{#1}}\providecommand{\ul}[1]{\underline{#1}}
% \alert (beamer) inventée par les modèles : texte mis en évidence
\providecommand{\alert}[1]{\textbf{\textcolor{poppink}{#1}}}
% variantes ASCII d'ordinaux écrites en macro
\providecommand{\iere}{\textsuperscript{re}}\providecommand{\ieme}{\textsuperscript{e}}\providecommand{\ere}{\textsuperscript{re}}\providecommand{\eme}{\textsuperscript{e}}\providecommand{\er}{\textsuperscript{er}}\providecommand{\ie}{\textsuperscript{e}}
% Ordinaux français écrits en macro par les modèles : 1\ière, 3\ième
\newcommand{\ière}{\textsuperscript{re}}\newcommand{\ième}{\textsuperscript{e}}
% \exemple (inventée par les modèles) : étiquette « Exemple : »
\providecommand{\exemple}{\textbf{Exemple~:}\ }
% \square utilisable aussi hors mode math (cases à cocher)
\AtBeginDocument{\let\squareMath\square\renewcommand{\square}{\ensuremath{\squareMath}}}
% Mot ou phrase en allemand à l'intérieur d'un paragraphe français
\newcommand{\all}[1]{\ifmmode\text{\textit{#1}}\else\begingroup\selectlanguage{ngerman}\itshape #1\endgroup\fi}
\let\esp\all % alias toléré
\newtcolorbox{exemplebox}[1][]{popbase,colframe=poporange,before upper={\popboxhead{Exemple}{poporange}{#1}}}
\newtcolorbox{definition}[1][]{popbase,colframe=popblue,before upper={\popboxhead{Définition}{popblue}{#1}}}
\newtcolorbox{propriete}[1][]{popbase,colframe=poppurple,before upper={\popboxhead{Propriété}{poppurple}{#1}}}
\newtcolorbox{methode}[1][]{popbase,colframe=popgold,before upper={\popboxhead{Méthode}{popgold}{#1}}}
\newtcolorbox{attention}[1][]{popbase,colframe=poppink,before upper={\popboxhead{Attention}{poppink}{#1}}}
\newtcolorbox{experience}[1][]{popbase,colframe=popblue,colback=popblueL,before upper={\popboxhead{Expérience}{popblue}{#1}}}
% « Le savais-tu ? » : carte violette pleine (contexte, culture, Cameroun)
\newtcolorbox{savaistu}[1][]{popbase,colback=poppurple,colframe=popink,
  fontupper=\popsemi\fontsize{10.4}{14.2}\selectfont\color{white},
  before upper={\poppill{Le savais-tu\,?}{white}{poppurple}\ifx\relax#1\relax\else\hspace{6pt}{\popxbold\color{white}#1}\fi\par\vspace{7pt}}}
% Exercice résolu : énoncé en haut, \tcblower puis la solution sur fond crème
\newcounter{popexo}\newcounter{popexr}
\newtcolorbox{exoresolu}[1][]{popbase,colframe=poporange,colbacklower=popcream,
  segmentation style={popink,line width=1.2pt,dashed},
  before upper={\stepcounter{popexr}\popboxhead{Exercice résolu \thepopexr}{poporange}{#1}},
  before lower={\poppill{Solution}{popink}{popcream}\par\vspace{6pt}}}
% Exercice (non résolu en place) — la correction est dans la partie Corrigés
\newtcolorbox{exercice}[1][]{popbase,colframe=popink,
  before upper={\stepcounter{popexo}\popboxhead{Exercice \thepopexo}{popink}{#1}}}
% Corrigé
\newtcolorbox{corrige}[1][]{popbase,colframe=popgreen,colback=popgreenL,
  before upper={\popboxhead{Corrigé}{popgreen}{#1}}}
% Objectifs / prérequis / bilan
\newtcolorbox{objectifs}{popbase,colframe=popink,colback=poporangeL,
  before upper={\popboxhead{Objectifs de la leçon}{popink}{}}}
\newtcolorbox{prerequis}{popbase,colframe=popink,colback=popblueL,
  before upper={\popboxhead{Prérequis}{popblue}{}}}
\newtcolorbox{bilan}{popbase,colframe=popink,colback=white,boxrule=2.8pt,
  before upper={\popboxhead{Fiche bilan}{poporange}{L'essentiel à connaître pour l'examen}}}
% Figure encadrée avec légende
\newtcolorbox{popfigure}[1][]{enhanced,breakable=false,colback=white,colframe=popline,boxrule=1.4pt,arc=8pt,
  left=10pt,right=10pt,top=10pt,bottom=8pt,before skip=10pt,after skip=12pt,halign=center,
  lower separated=false,halign lower=center,
  fontlower=\popsemi\fontsize{9}{11.5}\selectfont\color{popmuted},
  #1}
\newcommand{\legende}[1]{\par\vspace{6pt}{\popsemi\fontsize{9}{11.5}\selectfont\color{popmuted}#1}}

% ============================================================
% Signature OnBuch+
% ============================================================
\newcommand{\poplogoinline}[1]{{\poplogo\fontsize{#1}{#1}\selectfont\color{popink}OnBuch\color{poporange}+}}
\newcommand{\signatureonbuch}{%
  \begin{tcolorbox}[enhanced,colback=popink,colframe=popink,boxrule=2.2pt,arc=10pt,
      drop shadow=poporange,left=14pt,right=14pt,top=10pt,bottom=10pt,before skip=0pt,after skip=0pt]
    \tikz[baseline=-0.7ex]\node[circle,fill=popgreen,draw=popcream,line width=1.3pt,minimum size=22pt,inner sep=0]{\popcheck{white}};%
    \hspace{9pt}\begin{minipage}[c]{0.62\linewidth}
      {\popxbold\fontsize{12}{13}\selectfont\color{popcream}Signé\ \textbullet\ {\poplogo OnBuch\color{poporange}+}}\par\vspace{2pt}
      {\raggedright\popsemi\fontsize{8}{10.5}\selectfont\color{popline}\MakeUppercase{Équipe pédagogique OnBuch+}\\\MakeUppercase{Conforme au programme officiel MINESEC}\par}
    \end{minipage}\hfill
    {\poplogo\fontsize{22}{22}\selectfont\color{popcream}OnBuch\color{poporange}+}
  \end{tcolorbox}%
}

% ============================================================
% Page de garde
% #1 titre · #2 matière · #3 niveau · #4 module · #5 n° leçon · #6 durée ·
% #7 description / objectifs (texte ou itemize)
% ============================================================
\newcommand{\pagegarde}[7]{%
  \thispagestyle{empty}
  \newgeometry{a4paper,margin=1.7cm,top=1.6cm,bottom=1.4cm}%
  \begin{tikzpicture}[remember picture,overlay]
    \fill[poporange!18] ([xshift=-3.2cm,yshift=-3.6cm]current page.north east) circle (4.6cm);
    \fill[poppurple!14] ([xshift=2.4cm,yshift=7.6cm]current page.south west) circle (3.8cm);
  \end{tikzpicture}%
  {\poplogo\fontsize{30}{30}\selectfont\color{popink}OnBuch\color{poporange}+}\hfill
  \raisebox{0.6ex}{\poppill{Cours signé \textbullet\ Premium}{popink}{popcream}}\par
  \vspace{1pt}\popsquiggle{poppurple}\par
  \vspace{8pt}
  \begin{tcolorbox}[enhanced,colback=poporange,colframe=popink,boxrule=2.8pt,arc=13pt,
      drop shadow=popink,left=18pt,right=18pt,top=12pt,bottom=13pt,after skip=10pt]
    \poppill{#2 \textbullet\ #3}{popink}{popcream}\hspace{5pt}\poppill{#4}{white}{popink}\par\vspace{8pt}
    {\popdisplay\fontsize{14}{14}\selectfont\color{popink}LEÇON #5}\par\vspace{4pt}
    {\raggedright\hyphenpenalty=10000\exhyphenpenalty=10000\popdisplay\fontsize{25}{27}\selectfont\color{popcream}\MakeUppercase{#1}\par}
    \vspace{8pt}
    {\popxbold\fontsize{9.5}{9.5}\selectfont\color{popink}\MakeUppercase{Durée conseillée : #6}}
  \end{tcolorbox}
  \begin{tcolorbox}[enhanced,colback=white,colframe=popink,boxrule=2.2pt,arc=10pt,
      drop shadow=popink,left=16pt,right=16pt,top=10pt,bottom=10pt,after skip=8pt]
    \poppill{Dans cette leçon}{poppurple}{white}\par\vspace{5pt}
    \popsemi\fontsize{9.3}{12.6}\selectfont\color{popink}#7
  \end{tcolorbox}
  \noindent\begin{minipage}[t]{0.487\linewidth}
    \begin{tcolorbox}[enhanced,colback=popgreen,colframe=popink,boxrule=2.2pt,arc=10pt,
        drop shadow=popink,left=11pt,right=11pt,top=10pt,bottom=10pt,height=3.1cm,valign=center]
      \tcbox[colback=white,colframe=popink,boxrule=1.3pt,arc=5pt,boxsep=0pt,nobeforeafter,
        left=3pt,right=3pt,top=3pt,bottom=3pt,valign=center]{\qrcode[height=1.9cm]{https://play.google.com/store/apps/details?id=com.luvvix.onbuch&hl=fr}}%
      \hspace{8pt}\begin{minipage}[c]{0.5\linewidth}
        {\popdisplay\fontsize{11}{12.5}\selectfont\color{white}\MakeUppercase{Télécharge OnBuch}}\par\vspace{4pt}
        {\raggedright\popsemi\fontsize{8.4}{11}\selectfont\color{white}Gratuit \textbullet\ Google Play\\Tuteur IA, annales, concours, résultats\par}
      \end{minipage}
    \end{tcolorbox}
  \end{minipage}\hfill
  \begin{minipage}[t]{0.487\linewidth}
    \begin{tcolorbox}[enhanced,colback=poppurple,colframe=popink,boxrule=2.2pt,arc=10pt,
        drop shadow=popink,left=11pt,right=11pt,top=10pt,bottom=10pt,height=3.1cm,valign=center]
      \tikz[baseline=-0.7ex]\node[circle,fill=white,draw=popink,line width=1.3pt,minimum size=1.6cm,inner sep=0]{\popdisplay\fontsize{24}{24}\selectfont\color{poppurple}+};%
      \hspace{8pt}\begin{minipage}[c]{0.6\linewidth}
        {\popdisplay\fontsize{11}{12.5}\selectfont\color{white}\MakeUppercase{Passe à OnBuch+}}\par\vspace{4pt}
        {\raggedright\popsemi\fontsize{8.4}{11}\selectfont\color{white}Tous les cours signés, Tuteur IA illimité, fiches PDF — onglet OnBuch+ de l'app\par}
      \end{minipage}
    \end{tcolorbox}
  \end{minipage}
  \vspace{8pt}
  \popcontenu
  \vfill
  \signatureonbuch
  \restoregeometry
  \newpage
}

% Rangée « Ce cours contient » (page de garde)
\newcommand{\poptile}[3]{% #1 couleur #2 gros texte #3 légende
  \begin{tcolorbox}[enhanced,colback=#1,colframe=popink,boxrule=2pt,arc=9pt,shadow={2.5pt}{-2.5pt}{0pt}{popink},
      left=6pt,right=6pt,top=9pt,bottom=9pt,halign=center,valign=center,height=1.75cm,width=\linewidth,nobeforeafter]
    {\popdisplay\fontsize{12.5}{13}\selectfont\color{white}\MakeUppercase{#2}}\par\vspace{3pt}
    {\centering\popsemi\fontsize{7.8}{9.5}\selectfont\color{white}#3\par}
  \end{tcolorbox}}
\newcommand{\popcontenu}{%
  \par\noindent{\popxboldls\fontsize{7.5}{7.5}\selectfont\color{popmuted}CE COURS SIGNÉ CONTIENT}\par\vspace{7pt}
  \noindent\begin{minipage}[t]{0.235\linewidth}\poptile{popblue}{Cours}{complet, illustré, pas à pas}\end{minipage}\hfill
  \begin{minipage}[t]{0.235\linewidth}\poptile{poporange}{Méthodes}{et exercices résolus}\end{minipage}\hfill
  \begin{minipage}[t]{0.235\linewidth}\poptile{poppink}{Exercices}{type examen + corrigés}\end{minipage}\hfill
  \begin{minipage}[t]{0.235\linewidth}\poptile{popgreen}{Fiche bilan}{l'essentiel à retenir}\end{minipage}\par}

% Sommaire
\newtcolorbox{sommaire}{enhanced,colback=white,colframe=popink,boxrule=2.2pt,arc=10pt,
  drop shadow=popink,left=18pt,right=18pt,top=13pt,bottom=13pt,before skip=6pt,after skip=20pt,
  before upper={\poppill{Sommaire}{poppurple}{white}\par\vspace{9pt}\popsemi\fontsize{10.6}{14.5}\selectfont\color{popink}}}

% Signature de fin de cours — poussée tout en bas de la dernière page
\newcommand{\findecours}{%
  \par\vfill\nopagebreak
  \begin{center}\popsquiggle{poporange}\end{center}\vspace{4pt}
  \signatureonbuch
}
\AtBeginDocument{\def\llbracket{\mathopen{[\mkern-3.5mu[}}\def\rrbracket{\mathclose{]\mkern-3.5mu]}}} % intervalles d'entiers (glyphe ⟦ absent de la police)
% Glyphes absents de Fira Math : redéfinis à partir de glyphes disponibles
\AtBeginDocument{%
  \def\setminus{\mathbin{\backslash}}\let\smallsetminus\setminus
  \def\vdots{\mathord{\vbox{\baselineskip3.2pt\lineskiplimit0pt\kern4pt\hbox{.}\hbox{.}\hbox{.}}}}%
  \def\ddots{\mathinner{\mkern1mu\raise6.5pt\hbox{.}\mkern2mu\raise3.6pt\hbox{.}\mkern2mu\raise0.7pt\hbox{.}\mkern1mu}}%
  \def\triangle{\mathord{\scalebox{0.9}{$\symup{\Delta}$}}}%
  \def\nabla{\mathord{\raisebox{\depth}{\scalebox{1}[-1]{$\symup{\Delta}$}}}}%
  \def\checkmark{\popcheck{popdark}}%
  \def\star{\mathbin{\ast}}\def\bigstar{\mathord{\ast}}%
  \def\diamond{\mathbin{\scalebox{0.6}{$\Diamond$}}}%
  \def\bigcup{\mathop{\mathchoice{\raisebox{-0.3em}{\scalebox{1.7}{$\cup$}}}{\scalebox{1.25}{$\cup$}}{\cup}{\cup}}\displaylimits}%
  \def\bigcap{\mathop{\mathchoice{\raisebox{-0.3em}{\scalebox{1.7}{$\cap$}}}{\scalebox{1.25}{$\cap$}}{\cap}{\cap}}\displaylimits}%
  \def\lesseqqgtr{\mathrel{\lessgtr}}\def\bigoplus{\mathop{\oplus}\displaylimits}%
}
\providecommand{\ssi}{si et seulement si} % abréviation fréquente
\AtBeginDocument{\providecommand{\wideparen}{\overparen}} % arc de cercle

\begin{document}

\pagegarde{Smartphones et cybercafés (avantages et inconvénients des médias)}{Allemand}{1ère A}{Chapitre 4 : Média et communication}{ALA15}{2 h}{Cette leçon propose la lecture et l'exploitation d'un texte sur les smartphones et les cybercafés, leurs avantages et leurs inconvénients. Elle développe le vocabulaire thématique des médias et entraîne la comparaison, les subordonnées en « dass » et le parfait des verbes forts.
\vspace{4pt}
\begin{multicols}{2}\small
\begin{itemize}
\item À la fin de cette leçon, je sais comprendre globalement puis en détail un texte allemand sur les smartphones et les cybercafés.
\item À la fin de cette leçon, je sais employer correctement le vocabulaire thématique des médias (articles, pluriels, familles de mots).
\item À la fin de cette leçon, je sais former et utiliser la comparaison (Komparativ, Superlativ, als/wie) pour comparer deux médias.
\item À la fin de cette leçon, je sais construire des subordonnées en « dass » avec le verbe à la fin pour exprimer une opinion.
\item À la fin de cette leçon, je sais conjuguer les verbes forts au parfait (Perfekt) avec haben ou sein.
\item À la fin de cette leçon, je sais répondre à des questions de compréhension en phrases complètes.
\end{itemize}\end{multicols}}

\begin{sommaire}
\begin{multicols}{2}
\begin{enumerate}
\item Texte support et lecture globale
\item Compréhension détaillée : questions et exploitation
\item Lexique thématique : les médias et la communication
\item Grammaire 1 : la comparaison (Komparativ und Superlativ)
\item Grammaire 2 : les subordonnées en « dass » et le parfait des verbes forts
\item Méthode du commentaire, commentaire modèle et production guidée
\item Activité d'intégration
\item Méthodes et exercices résolus
\item Exercices
\item Corrigés des exercices
\item Fiche bilan
\end{enumerate}
\end{multicols}
\end{sommaire}

\begin{objectifs}
\begin{itemize}
\item À la fin de cette leçon, je sais comprendre globalement puis en détail un texte allemand sur les smartphones et les cybercafés.
\item À la fin de cette leçon, je sais employer correctement le vocabulaire thématique des médias (articles, pluriels, familles de mots).
\item À la fin de cette leçon, je sais former et utiliser la comparaison (Komparativ, Superlativ, als/wie) pour comparer deux médias.
\item À la fin de cette leçon, je sais construire des subordonnées en « dass » avec le verbe à la fin pour exprimer une opinion.
\item À la fin de cette leçon, je sais conjuguer les verbes forts au parfait (Perfekt) avec haben ou sein.
\item À la fin de cette leçon, je sais répondre à des questions de compréhension en phrases complètes.
\item À la fin de cette leçon, je sais résumer un texte en quelques phrases.
\item À la fin de cette leçon, je sais produire un court paragraphe argumenté sur les avantages et les inconvénients d'un média.
\end{itemize}
\end{objectifs}

\begin{prerequis}
\begin{itemize}
\item Le présent des verbes réguliers et irréguliers (Seconde)
\item Les articles définis et indéfinis au nominatif et à l'accusatif (Seconde)
\item La place du verbe en position 2 dans la phrase déclarative (Seconde)
\item Le parfait des verbes faibles (début de 1ère)
\item Le vocabulaire de base de la vie quotidienne et de l'école (Seconde)
\end{itemize}
\end{prerequis}

\newpage

\coursec{Texte support et lecture globale}

\courssub{Situation d'accroche et problématique}

À Yaoundé comme à Douala, presque chaque élève possède un \all{Smartphone} et beaucoup se rendent encore dans les cybercafés pour imprimer leurs devoirs ou se connecter quand les données mobiles sont chères. Le matin, avant les cours, on voit des élèves de Première A consulter leurs messages WhatsApp ; le soir, après les études, les cybercafés du quartier se remplissent de lycéens qui tapent un exposé ou impriment un document. En Allemagne aussi, les jeunes passent des heures sur leur téléphone, mais les débats sur les dangers des réseaux sociaux y sont très vifs : les journaux scolaires publient régulièrement des articles sur le sujet, et les élèves apprennent à argumenter, à comparer, à donner un avis nuancé.

Comment parler en allemand des avantages et des inconvénients de ces médias ? Comment comparer le \all{Smartphone} et le cybercafé et exprimer une opinion argumentée ? C'est exactement ce que nous allons apprendre dans cette leçon. Elle nous donne le texte, les mots et les outils grammaticaux pour le faire : la comparaison, les subordonnées en \all{dass} et le parfait des verbes forts. Tout commence par la lecture attentive d'un texte court, écrit dans le style d'un article de journal scolaire allemand.

\courssub{Préparation à la lecture : la méthode en quatre étapes}

Avant même de lire le texte, il faut savoir \emph{comment} lire un texte allemand. Beaucoup d'élèves se précipitent sur le dictionnaire ou le glossaire dès le premier mot inconnu : c'est une erreur. La lecture d'un texte en langue étrangère est une démarche active, qui se fait en plusieurs passages.

\begin{methode}[Lire un texte allemand en quatre étapes]
\begin{enumerate}
\item \textbf{Lire le titre et anticiper.} Le titre donne presque toujours le thème. Ici : \all{Das Smartphone oder das Internetcafé?} — le mot \all{oder} (« ou ») annonce une \emph{comparaison} entre deux médias. On peut déjà prévoir que le texte parlera des avantages et des inconvénients de chacun.
\item \textbf{Première lecture rapide pour le sens global.} On lit le texte entier sans s'arrêter, sans traduire, juste pour répondre à la question : « De quoi parle ce texte ? »
\item \textbf{Deuxième lecture avec le glossaire.} Cette fois, on s'aide du glossaire pour comprendre les détails : qui parle ? à qui ? quelle est l'opinion de l'auteur ?
\item \textbf{Souligner les mots-clés de l'argumentation.} Les petits mots comme \all{einerseits / andererseits} (« d'une part / d'autre part »), \all{aber} (« mais »), \all{deshalb} (« c'est pourquoi »), \all{trotzdem} (« pourtant »), \all{auch} (« aussi ») sont les panneaux indicateurs du texte : ils montrent la logique du raisonnement.
\end{enumerate}
\end{methode}

\begin{attention}[Le piège de la traduction mot à mot]
Ne traduis jamais un texte allemand mot à mot : tu obtiendrais un français bizarre et tu perdrais le sens. Cherche d'abord le \textbf{sens global de la phrase} avant de consulter le glossaire. Exemple : \all{Das Handy lenkt vom Lernen ab.} Traduction mot à mot impossible (« le portable distrait de l'apprendre » ?) — sens global : « Le portable distrait des études. » Pense en \emph{idées}, pas en mots isolés.
\end{attention}

\courssub{Le texte support : Das Smartphone oder das Internetcafé?}

\textbf{Lecture silencieuse.} Lis maintenant le texte en silence, une première fois, sans t'arrêter sur les mots inconnus. Puis ton professeur le lira à voix haute : écoute la prononciation et l'intonation, en particulier le « ch » de \all{praktisch} (« prak-tiche »), le « V » de \all{Vorteil}, prononcé « f » (« for-taïl »), et le « z » de \all{Zeitung} (« ts »).

\begin{textebox}[Das Smartphone oder das Internetcafé? — Artikel aus einer Schülerzeitung]
Viele Jugendliche in Kamerun und in Deutschland benutzen jeden Tag ihr Smartphone. Einerseits ist das praktisch: Man kann chatten, Fotos machen und schnell Informationen finden. Andererseits ist die Verbindung oft teuer, \textbf{deshalb} lenkt das Handy vom Lernen ab. Das Internetcafé hat auch Vorteile: Dort ist der Computer groß, man kann Dokumente drucken, und eine Stunde kostet nur eine kleine Gebühr. \textbf{Trotzdem} muss man das Haus verlassen, und die Verbindung ist manchmal langsam. Ich denke, dass das Smartphone für die Kommunikation besser ist, aber für die Schularbeit gehe ich lieber ins Internetcafé.
\end{textebox}

\textbf{Glossaire (à garder sous les yeux pendant la lecture) :}

\begin{center}
\begin{tabularx}{\linewidth}{|X|X|X|}
\hline
\rowcolor{popblueL} \textbf{Allemand} & \textbf{Français} & \textbf{Remarque} \\
\hline
der Vorteil, -e & l'avantage & contraire : der Nachteil \\
\hline
der Nachteil, -e & l'inconvénient & \all{Vor- und Nachteile} = le pour et le contre \\
\hline
die Verbindung, -en & la connexion & famille : \all{verbinden} = connecter, relier \\
\hline
praktisch & pratique & adj. invariable \\
\hline
chatten & chatter, discuter en ligne & verbe régulier, empr. à l'anglais \\
\hline
teuer & cher, coûteux & adj. invariable \\
\hline
ablenken (trennbar) & distraire & \all{lenkt ... ab} = distrait (particule \all{ab} à la fin) \\
\hline
drucken & imprimer & le nom : der Drucker = l'imprimante \\
\hline
die Gebühr, -en & le tarif, les frais & \all{eine kleine Gebühr} = un petit prix \\
\hline
verlassen & quitter, laisser & verbe fort : \all{verlässt - verließ - hat verlassen} \\
\hline
langsam & lent, lente & adj. invariable \\
\hline
die Schularbeit, -en & le travail scolaire, le devoir & composé : \all{Schule} + \all{Arbeit} \\
\hline
lieber & préférablement, plus volontiers & comparatif de \all{gerne} \\
\hline
\end{tabularx}
\end{center}

\begin{savaistu}[Das Handy : un faux anglicisme]
En Allemagne, on dit familièrement \all{das Handy} pour « téléphone portable » — mais ce mot, les Anglais ne le comprennent pas ! C'est un mot inventé en Allemagne qui \emph{ressemble} à de l'anglais sans en être : on appelle cela un \emph{faux anglicisme} (\all{der Scheinanglizismus}). Le mot officiel est \all{das Mobiltelefon} ou, aujourd'hui, \all{das Smartphone}. Au Cameroun, tu entendras aussi \all{das Handy} dans les cours d'allemand : c'est le mot le plus courant dans la bouche des Allemands.
\end{savaistu}

\courssub{Compréhension globale : de quoi parle le texte ?}

Après la première lecture rapide, on répond aux questions fondamentales du lecteur. Ces questions sont les mêmes pour tout texte journalistique, en allemand comme en français.

\textbf{1. Le thème.} Le texte parle des \cle{médias} de la vie quotidienne des jeunes, et plus précisément de deux outils de communication et de travail : le \all{Smartphone} et l'\all{Internetcafé} (le cybercafé).

\textbf{2. Le type de texte.} C'est un \cle{article de journal scolaire} (\all{ein Artikel aus einer Schülerzeitung}). Indices : le sujet concerne les jeunes, le ton est personnel (\all{Ich denke...}), l'argumentation est simple et équilibrée (avantages et inconvénients des deux côtés), et l'auteur termine par son opinion personnelle.

\textbf{3. Le locuteur (qui parle ?).} Un élève ou un jeune journaliste scolaire, qui écrit à la première personne à la fin du texte : \all{Ich denke, dass...} (« Je pense que... »).

\textbf{4. Le destinataire (à qui ?).} Les lecteurs du journal de l'école, donc d'autres élèves — au Cameroun comme en Allemagne, puisque le texte mentionne les deux pays : \all{Viele Jugendliche in Kamerun und in Deutschland}.

\textbf{5. Les deux médias comparés.} Le titre pose la question avec \all{oder} (« ou ») : faut-il choisir le Smartphone \emph{ou} le cybercafé ? Tout le texte est construit sur cette comparaison.

\begin{exemplebox}[Question de compréhension globale et réponse modèle]
\textbf{Frage :} \all{Worum geht es in diesem Text?} (De quoi parle ce texte ?)\par
\textbf{Antwort :} \all{Der Text vergleicht das Smartphone und das Internetcafé. Er zeigt die Vorteile und die Nachteile der beiden Medien.}\par
« Le texte compare le smartphone et le cybercafé. Il montre les avantages et les inconvénients des deux médias. »
\end{exemplebox}

\courssub{Les mots-repères de l'argumentation}

Un texte argumentatif, même court, possède une ossature : de petits mots qui organisent les idées. Les repérer, c'est comprendre la logique de l'auteur sans même traduire toutes les phrases. Reprenons le texte et soulignons ces \cle{connecteurs}.

\begin{center}
\begin{tabularx}{\linewidth}{|X|X|X|}
\hline
\rowcolor{popblueL} \textbf{Deutsch} & \textbf{Français} & \textbf{Fonction dans le texte} \\
\hline
einerseits ... andererseits & d'une part ... d'autre part & présente le pour et le contre du Smartphone \\
\hline
deshalb & c'est pourquoi, donc & introduit une conséquence (lien cause-effet) \\
\hline
auch & aussi & ajoute un argument du même côté (Internetcafé) \\
\hline
trotzdem & pourtant, malgré cela & introduit une concession (inconvénient malgré avantages) \\
\hline
aber & mais & introduit une objection, un inconvénient \\
\hline
ich denke, dass ... & je pense que ... & exprime l'opinion personnelle (subordonnée) \\
\hline
\end{tabularx}
\end{center}

Dans notre texte, la structure argumentative est parfaitement visible :

\begin{itemize}
\item \all{Einerseits ist das praktisch} → avantage du Smartphone ;
\item \all{Andererseits ist die Verbindung oft teuer, deshalb lenkt das Handy vom Lernen ab} → inconvénients du Smartphone (cause + conséquence) ;
\item \all{Das Internetcafé hat auch Vorteile} → avantages du cybercafé (ajout) ;
\item \all{Trotzdem muss man das Haus verlassen, und die Verbindung ist manchmal langsam} → concession + inconvénients du cybercafé ;
\item \all{Ich denke, dass ...} → conclusion et opinion personnelle.
\end{itemize}

\begin{attention}[La place du verbe après ces connecteurs]
Attention, francophone ! Après \all{aber}, \all{deshalb}, \all{trotzdem} et \all{auch}, le verbe conjugué reste en \textbf{deuxième position}, comme dans une phrase principale normale : \all{Aber man \textbf{muss} das Haus verlassen.} / \all{Deshalb \textbf{lenkt} das Handy ab.} Ne dis jamais \all{aber man das Haus verlassen muss}. En revanche, après \all{dass} (subordonnée), le verbe va à la \textbf{fin} : \all{Ich denke, dass das Smartphone besser \textbf{ist}.} Nous étudierons cette règle en détail dans la section grammaire.
\end{attention}

\courssub{Carte mentale : organiser ce que le texte nous apprend}

Pour mémoriser le contenu du texte et préparer le résumé, rien de tel qu'une carte mentale. Voici celle de notre thème \all{Medien} (« les médias »), construite à partir du texte :

\begin{popfigure}
\begin{tikzpicture}[
  mindmap,
  grow cyclic,
  every node/.style={concept},
  root concept/.append style={concept color=popblue, font=\large\bfseries},
  level 1/.append style={level distance=3.2cm, sibling angle=180},
  level 2/.append style={level distance=2.6cm, sibling angle=70},
  text=white
]
\node[root concept]{Medien}
  child[concept color=popgreen] { node {Smartphone}
    child[concept color=popgreen!70] { node[align=center, font=\small] {Vorteile\\ praktisch, chatten,\\ Fotos, Infos} }
    child[concept color=poporange] { node[align=center, font=\small] {Nachteile\\ teuer, Ablenkung,\\ Datenvolumen} }
  }
  child[concept color=poppurple] { node {Internetcafé}
    child[concept color=popgreen!70] { node[align=center, font=\small] {Vorteile\\ großer PC, drucken,\\ kleine Gebühr} }
    child[concept color=poporange] { node[align=center, font=\small] {Nachteile\\ Haus verlassen,\\ langsam, öffentlich} }
  };
\end{tikzpicture}
\legende{Carte mentale du texte : les deux médias comparés, avec leurs avantages (\all{Vorteile}) et leurs inconvénients (\all{Nachteile}).}
\end{popfigure}

Cette carte mentale est aussi un \emph{plan de résumé} : en la parcourant de gauche à droite, tu peux raconter tout le texte sans l'avoir sous les yeux.

\courssub{Reformulation : dire le texte avec ses propres mots}

Dernière étape de la lecture globale : \cle{reformuler} le contenu en une phrase, avec ses propres mots. C'est la preuve que l'on a vraiment compris. La reformulation n'est pas une traduction : on garde l'idée, on change les mots.

\begin{exemplebox}[Trois reformulations possibles du texte]
\begin{itemize}
\item \all{Der Text zeigt das Für und Wider von Smartphone und Internetcafé.} « Le texte montre le pour et le contre du smartphone et du cybercafé. »
\item \all{Beide Medien haben Vorteile und Nachteile: Das Handy ist gut für die Kommunikation, das Internetcafé ist besser für die Schularbeit.} « Les deux médias ont des avantages et des inconvénients : le portable est bon pour la communication, le cybercafé est meilleur pour le travail scolaire. »
\item \all{Ein Jugendlicher vergleicht zwei Medien und sagt am Ende seine Meinung.} « Un jeune compare deux médias et donne son avis à la fin. »
\end{itemize}
\end{exemplebox}

\begin{exoresolu}[Questions de compréhension globale]
\textbf{Énoncé.} Réponds en allemand, par des phrases complètes, aux questions suivantes :
\begin{enumerate}
\item \all{Welche zwei Medien vergleicht der Text?}
\item \all{Nenne einen Vorteil des Smartphones!}
\item \all{Nenne einen Nachteil des Internetcafés!}
\item \all{Was ist die Meinung des Autors?}
\end{enumerate}
\tcblower
\textbf{Solution détaillée.}
\begin{enumerate}
\item \all{Der Text vergleicht das Smartphone und das Internetcafé.} « Le texte compare le smartphone et le cybercafé. » — On reprend les deux noms du titre, avec leurs articles : \all{das Smartphone}, \all{das Internetcafé}.
\item \all{Das Smartphone ist praktisch: Man kann chatten und schnell Informationen finden.} « Le smartphone est pratique : on peut chatter et trouver vite des informations. » — Un seul avantage suffisait ; on choisit le plus facile à formuler.
\item \all{Man muss das Haus verlassen, und die Verbindung ist manchmal langsam.} « Il faut quitter la maison, et la connexion est parfois lente. »
\item \all{Der Autor denkt, dass das Smartphone für die Kommunikation besser ist, aber er geht lieber ins Internetcafé für die Schularbeit.} « L'auteur pense que le smartphone est meilleur pour la communication, mais il préfère aller au cybercafé pour le travail scolaire. » — Attention à la subordonnée en \all{dass} : le verbe \all{ist} va à la fin.
\end{enumerate}
\end{exoresolu}

\begin{exercice}[À toi de jouer : lecture globale]
\begin{enumerate}
\item Lis le texte une dernière fois à voix haute, en respectant les pauses aux virgules et aux points.
\item Souligne dans le texte les six connecteurs : \all{einerseits}, \all{andererseits}, \all{deshalb}, \all{auch}, \all{trotzdem}, \all{aber}.
\item Recopie la carte mentale dans ton cahier et ajoute, en allemand, un avantage ou un inconvénient de ton invention à chaque branche (par exemple pour le Smartphone : \all{die Batterie ist schnell leer} — « la batterie se vide vite » ; pour l'Internetcafé : \all{man trifft Freunde} — « on retrouve des amis »).
\item Reformule le texte en une phrase allemande, sans regarder les modèles du cours.
\end{enumerate}
\end{exercice}

\begin{corrige}[Exercice « À toi de jouer : lecture globale »]
\textbf{1.} Lecture orale : veille à prononcer \all{Vorteil} « for-taïl », \all{Jugendliche} « you-guend-li-che », \all{Gebühr} « gueu-burr », \all{Schularbeit} « chou-lar-ba-ït ».
\textbf{2.} Les connecteurs se trouvent : phrase 2 (\all{Einerseits}), phrase 3 (\all{Andererseits}, \all{deshalb}), phrase 4 (\all{auch}), phrase 5 (\all{Trotzdem}, \all{aber}).
\textbf{3.} Exemples d'ajouts possibles — Smartphone, \all{Nachteile} : \all{die Batterie ist schnell leer}, \all{die Daten sind teuer}. Internetcafé, \all{Vorteile} : \all{man trifft Freunde}, \all{es gibt schnelles WLAN}.
\textbf{4.} Reformulation attendue, par exemple : \all{Der Text zeigt die Vorteile und Nachteile von Handy und Internetcafé und der Autor sagt seine Meinung.} Toute phrase correcte qui reprend l'idée générale (comparaison + opinion) est acceptée.
\end{corrige}

\begin{aretenir}[L'essentiel de cette section]
\begin{itemize}
\item Lire un texte allemand se fait en \textbf{quatre étapes} : titre et anticipation, lecture rapide globale, lecture détaillée avec le glossaire, repérage des connecteurs.
\item Notre texte est un \textbf{article de journal scolaire} qui compare le \all{Smartphone} et l'\all{Internetcafé} : chacun a ses \all{Vorteile} et ses \all{Nachteile}.
\item Les connecteurs \all{einerseits / andererseits}, \all{deshalb}, \all{trotzdem}, \all{aber}, \all{auch} sont les panneaux indicateurs de l'argumentation.
\item Après \all{aber, deshalb, trotzdem, auch} : verbe en \textbf{deuxième position} (phrase principale). Après \all{dass} : verbe à la \textbf{fin} (subordonnée).
\item Savoir \textbf{reformuler} en une phrase, c'est prouver qu'on a compris.
\end{itemize}
\end{aretenir}

\coursec{Compréhension détaillée : questions et exploitation}

Dans la section précédente, nous avons lu le texte support une première fois pour en dégager le sens global : un lycéen de Yaoundé compare le \all{Smartphone} et l'\all{Internetcafé} pour ses travaux scolaires. Nous passons maintenant à la \cle{lecture détaillée} : il s'agit de répondre à des questions précises, de justifier ses réponses par le texte et d'apprendre à rédiger des réponses complètes en allemand. C'est une compétence essentielle pour l'évaluation de la compréhension de l'écrit.

\courssub{Rappel du texte support}

Avant de répondre aux questions, relisons le texte. Le voici à nouveau, avec un glossaire des mots difficiles.

\begin{textebox}[Mein Smartphone und das Internetcafé — Kwame, 17 ans, Yaoundé]
Ich heiße Kwame und besuche die Première am Gymnasium in Yaoundé. Fast alle meine Kameraden haben ein Smartphone. Ich auch. Mit dem Handy kann ich chatten, Fotos machen, Musik hören und im Internet suchen. Das ist praktisch, aber das Handy lenkt mich oft vom Lernen ab: Wenn ich Hausaufgaben mache, schaue ich jede Minute auf den Bildschirm.

Für die Schularbeit gehe ich lieber ins Internetcafé. Eine Stunde kostet dort 300 Francs. Die Verbindung ist schnell, und ich kann dort Dokumente drucken. Das ist ein großer Vorteil, denn zu Hause haben wir keinen Drucker. Aber das Internetcafé hat auch Nachteile: Es ist oft laut, manchmal ist die Verbindung langsam, und abends sind alle Computer besetzt.

Mein Smartphone ist teuer, und das Internetcafé kostet auch Geld. Trotzdem brauche ich beide Medien: das Handy für die Kommunikation mit meinen Freunden und das Internetcafé für die Schule. Ich glaube, dass das Internetcafé für die Schularbeit besser ist, weil ich mich dort besser konzentrieren kann.
\end{textebox}

\textbf{Glossaire (Wortschatz) :}
\begin{itemize}
\item \all{der Kamerad, -en} (Pl. \all{die Kameraden}) : le camarade
\item \all{praktisch} : pratique
\item \all{ablenken (lenkt ab, hat abgelenkt)} : distraire, détourner l'attention (verbe à particule séparable)
\item \all{der Bildschirm, -e} (Pl. \all{die Bildschirme}) : l'écran
\item \all{drucken} : imprimer
\item \all{der Drucker, -} (Pl. \all{die Drucker}) : l'imprimante
\item \all{besetzt} : occupé
\item \all{trotzdem} : quand même, malgré cela
\item \all{sich konzentrieren (auf + Akk.)} : se concentrer (sur)
\end{itemize}

\courssub{La méthode : comment répondre à une question de compréhension}

Beaucoup d'élèves perdent des points non pas parce qu'ils n'ont pas compris le texte, mais parce qu'ils répondent mal : réponse d'un seul mot, phrase recopiée entièrement, réponse hors sujet. Voici la démarche rigoureuse à suivre.

\begin{methode}[Répondre à une question de compréhension en quatre étapes]
\begin{enumerate}
\item \textbf{Souligner le mot interrogatif} et les mots-clés de la question : \all{wer} (qui), \all{was} (quoi, que), \all{warum} (pourquoi), \all{wo} (où), \all{wann} (quand), \all{wie} (comment), \all{wie viel / wie lange} (combien / combien de temps).
\item \textbf{Localiser le passage du texte} qui contient l'information : repérer les mots identiques ou les synonymes.
\item \textbf{Reformuler} l'information avec ses propres mots, sans recopier toute la phrase du texte.
\item \textbf{Rédiger une phrase complète} en reprenant les termes de la question : sujet + verbe conjugué en deuxième position. Jamais de « ja » ou « nein » seul !
\end{enumerate}
\end{methode}

\begin{popfigure}
\begin{tikzpicture}[
node distance=0.55cm,
etape/.style={rectangle, rounded corners, draw=popblue, fill=popblueL, align=center, text width=5.2cm, inner sep=5pt},
fleche/.style={-{Stealth[length=3mm]}, thick, popdark}]
\node[etape, align=center] (q){\textbf{1. Question}\\ \all{Warum geht Kwame ins Internetcafé?}};
\node[etape, below=of q, align=center] (r){\textbf{2. Repérage dans le texte}\\ \all{...ich kann dort Dokumente drucken.}};
\node[etape, below=of r, align=center] (f){\textbf{3. Reformulation}\\ drucken = imprimer ses documents};
\node[etape, below=of f, fill=popgreenL, draw=popgreen, align=center] (p){\textbf{4. Phrase complète de réponse}\\ \all{Er geht ins Internetcafé, weil er dort Dokumente drucken kann.}};
\draw[fleche] (q) -- (r);
\draw[fleche] (r) -- (f);
\draw[fleche] (f) -- (p);
\end{tikzpicture}
\legende{De la question à la réponse : la chaîne en quatre étapes}
\end{popfigure}

\begin{exemplebox}[Question et réponse modèles]
Question : \all{Warum geht der Autor lieber ins Internetcafé?}\\
Réponse : \all{Er geht lieber ins Internetcafé, weil er dort Dokumente drucken kann.}\\
« Pourquoi l'auteur préfère-t-il aller au cybercafé ? — Parce qu'il peut y imprimer des documents. »
\end{exemplebox}

\begin{propriete}[La réponse avec « weil » : le verbe à la fin]
Quand on justifie une réponse par \all{weil} (parce que), on construit une \cle{subordonnée} : le verbe conjugué se place \textbf{à la fin} de la proposition.
\begin{itemize}
\item \all{..., weil die Verbindung teuer \textbf{ist}.} (parce que la connexion est chère)
\item \all{..., weil er dort Dokumente drucken \textbf{kann}.} (parce qu'il peut y imprimer des documents ; l'infinitif précède le verbe conjugué)
\end{itemize}
\textbf{Attention :} en français, l'ordre ne change jamais ; en allemand, \all{weil} envoie le verbe en dernière position. Ne dites jamais \all{weil die Verbindung ist teuer} !
\par\medskip
\textbf{Repère :} Le texte utilise aussi \all{denn} (car) : \all{„...denn zu Hause haben wir keinen Drucker.“} Avec \all{denn}, l'ordre reste celui de la phrase principale (verbe en 2\up{ème} position) : \all{denn wir haben...} Ne confondez pas \all{weil} (subordonnée) et \all{denn} (coordination).
\end{propriete}

\courssub{Questions de compréhension factuelle}

Ces questions portent sur des faits explicites du texte : l'information est écrite noir sur blanc. Il suffit de la repérer et de la reformuler.

\begin{center}
\begin{tabularx}{\linewidth}{|X|X|}
\hline
\rowcolor{popblueL} \textbf{Question} & \textbf{Réponse attendue (phrase complète)} \\
\hline
\all{Wer benutzt das Smartphone?} & \all{Fast alle Kameraden und auch Kwame benutzen ein Smartphone.} \\
\hline
\all{Was kann man mit dem Handy machen?} & \all{Mit dem Handy kann man chatten, Fotos machen, Musik hören und im Internet suchen.} \\
\hline
\all{Was kostet eine Stunde im Internetcafé?} & \all{Eine Stunde im Internetcafé kostet 300 Francs.} \\
\hline
\all{Wo geht Kwame für die Schularbeit hin?} & \all{Für die Schularbeit geht er ins Internetcafé.} \\
\hline
\end{tabularx}
\end{center}

Remarquez la technique : la réponse \textbf{reprend les termes de la question}. À \all{Was kostet eine Stunde?} on répond \all{Eine Stunde kostet ...} (sujet en tête, verbe en 2\up{ème} position). Si un complément de phrase ouvre la phrase (ex. \all{Für die Schularbeit}), le verbe reste en 2\up{ème} position et le sujet le suit immédiatement : \all{geht er ...}.

\courssub{Questions de compréhension fine}

Ces questions demandent de combiner plusieurs informations ou de comprendre une explication donnée dans le texte.

\begin{exemplebox}[Compréhension fine : exemples résolus]
\textbf{Question 1 :} \all{Warum lenkt das Handy vom Lernen ab?}\\
Réponse : \all{Das Handy lenkt vom Lernen ab, weil Kwame jede Minute auf den Bildschirm schaut.}\\
« Pourquoi le portable distrait-il de l'étude ? Parce que Kwame regarde l'écran chaque minute. »

\textbf{Question 2 :} \all{Welche Nachteile hat das Internetcafé?}\\
Réponse : \all{Das Internetcafé ist oft laut, die Verbindung ist manchmal langsam, und abends sind alle Computer besetzt.}\\
« Quels inconvénients le cybercafé présente-t-il ? Il est souvent bruyant, la connexion est parfois lente et le soir tous les ordinateurs sont occupés. »
\end{exemplebox}

\courssub{Questions d'interprétation et justification par le texte}

Les questions d'interprétation exigent de dégager l'opinion de l'auteur. Il faut alors \cle{justifier} sa réponse par une courte citation du texte, introduite par exemple par \all{Zeile ...} (« ligne ... ») ou \all{Der Autor sagt: „...“}.

\begin{exemplebox}[Interprétation : exemple résolu]
\textbf{Question :} \all{Welches Medium bevorzugt der Autor für die Schularbeit? Warum?}\\
Réponse : \all{Für die Schularbeit bevorzugt der Autor das Internetcafé, weil er sich dort besser konzentrieren kann. Er sagt: „Ich glaube, dass das Internetcafé für die Schularbeit besser ist.“}\\
« Quel média l'auteur préfère-t-il pour le travail scolaire ? Pourquoi ? — Pour le travail scolaire, il préfère le cybercafé, parce qu'il peut s'y mieux concentrer. Il dit : „Je crois que le cybercafé est meilleur pour le travail scolaire.“ »
\end{exemplebox}

\begin{attention}[« gern », « lieber », « bevorzugen » : ne pas confondre !]
\begin{itemize}
\item \all{gern + verbe} : faire quelque chose \textbf{avec plaisir}. \all{Ich chatte gern.} = J'aime discuter en ligne.
\item \all{lieber + verbe} : faire quelque chose \textbf{de préférence} (comparaison entre deux choix). \all{Ich gehe lieber ins Internetcafé.} = Je préfère aller au cybercafé (plutôt qu'utiliser le téléphone).
\item \all{bevorzugen} : verbe soutenu signifiant « préférer » (transitif, accusatif). \all{Er bevorzugt das Internetcafé.} = Il préfère le cybercafé.
\end{itemize}
Piège de traduction : « Il aime mieux le cybercafé » se traduit par \all{Er geht lieber ins Internetcafé} ou \all{Er bevorzugt das Internetcafé}, jamais par \all{besser lieben} !
\end{attention}

\courssub{Vrai ou faux : Richtig oder Falsch ?}

L'exercice « vrai/faux » vérifie la précision de votre lecture. Règle d'or : \textbf{toute réponse doit être justifiée} par une phrase ou une citation du texte (\all{Begründung}). Une réponse « richtig » sans justification ne vaut que la moitié des points.

\begin{exoresolu}[Richtig oder Falsch ? Begründen Sie !]
Énoncé : lisez les affirmations suivantes, dites si elles sont vraies (\all{richtig}) ou fausses (\all{falsch}) et justifiez par le texte.
\begin{enumerate}
\item \all{Kwame hat kein Smartphone.}
\item \all{Eine Stunde im Internetcafé kostet 300 Francs.}
\item \all{Kwame hat zu Hause einen Drucker.}
\item \all{Abends sind im Internetcafé alle Computer frei.}
\end{enumerate}
\tcblower
Solution détaillée :
\begin{enumerate}
\item \all{Falsch.} Le texte dit : \all{„Fast alle meine Kameraden haben ein Smartphone. Ich auch.“} Kwame possède donc un smartphone.
\item \all{Richtig.} Citation : \all{„Eine Stunde kostet dort 300 Francs.“}
\item \all{Falsch.} Le texte dit : \all{„...denn zu Hause haben wir keinen Drucker.“} Il n'y a pas d'imprimante à la maison ; c'est pour cela qu'il imprime au cybercafé.
\item \all{Falsch.} Citation : \all{„...abends sind alle Computer besetzt.“} \all{besetzt} signifie « occupé », c'est le contraire de \all{frei} (« libre »). Attention au faux sens !
\end{enumerate}
\end{exoresolu}

\courssub{Le résumé guidé avec des connecteurs}

Dernière étape de l'exploitation : résumer le texte en 3 ou 4 phrases. Pour enchaîner les idées, on utilise des \cle{connecteurs} : \all{zuerst} (d'abord), \all{dann} (ensuite), \all{außerdem} (de plus), \all{schließlich} (finalement). Attention à la place du verbe : après ces connecteurs placés en tête de phrase, le verbe conjugué vient en deuxième position, donc \textbf{avant le sujet}.

\begin{center}
\begin{tabularx}{\linewidth}{|X|X|X|}
\hline
\rowcolor{popblueL} \textbf{Connecteur} & \textbf{Français} & \textbf{Exemple} \\
\hline
\all{zuerst} & d'abord & \all{Zuerst beschreibt der Autor sein Smartphone.} \\
\hline
\all{dann} & ensuite & \all{Dann erklärt er die Vorteile des Internetcafés.} \\
\hline
\all{außerdem} & de plus & \all{Außerdem nennt er die Nachteile.} \\
\hline
\all{schließlich} & finalement & \all{Schließlich sagt er seine Meinung.} \\
\hline
\end{tabularx}
\end{center}

\begin{exemplebox}[Résumé modèle du texte]
\all{Zuerst sagt Kwame, dass fast alle Kameraden ein Smartphone haben. Dann erklärt er, dass das Handy vom Lernen ablenkt. Außerdem beschreibt er die Vorteile und die Nachteile des Internetcafés. Schließlich meint er, dass das Internetcafé für die Schularbeit besser ist.}\\
« D'abord Kwame dit que presque tous ses camarades ont un smartphone. Ensuite il explique que le portable distrait de l'étude. De plus, il décrit les avantages et les inconvénients du cybercafé. Finalement, il estime que le cybercafé est meilleur pour le travail scolaire. »
\end{exemplebox}

\begin{exercice}[À vous de jouer !]
\begin{enumerate}
\item Répondez par des phrases complètes : \all{Wann schaut Kwame auf den Bildschirm? Warum druckt er im Internetcafé?}
\item Vrai ou faux ? Justifiez : \all{Die Verbindung im Internetcafé ist immer schnell.}
\item Rédigez votre propre résumé du texte en 3 phrases avec \all{zuerst}, \all{dann} et \all{schließlich}.
\end{enumerate}
\end{exercice}

\begin{corrige}[Exercice « À vous de jouer ! »]
\begin{enumerate}
\item \all{Wenn er Hausaufgaben macht, schaut er jede Minute auf den Bildschirm. Er druckt im Internetcafé, weil er zu Hause keinen Drucker hat.}
\item \all{Falsch.} Le texte dit : \all{„...manchmal ist die Verbindung langsam.“} La connexion n'est donc pas toujours rapide.
\item Exemple : \all{Zuerst beschreibt Kwame sein Smartphone. Dann erklärt er die Vorteile und Nachteile des Internetcafés. Schließlich sagt er, dass er das Internetcafé für die Schularbeit bevorzugt.}
\end{enumerate}
\end{corrige}

\begin{aretenir}[Les cinq règles d'or de la compréhension détaillée]
\begin{itemize}
\item Souligner le mot interrogatif (\all{wer, was, warum, wo, wann, wie}) avant de chercher la réponse.
\item Localiser le passage du texte, puis reformuler sans tout recopier.
\item Toujours répondre par une phrase complète, jamais par « ja » ou « nein » seul.
\item Justifier par une citation courte : \all{Der Autor sagt: „...“}.
\item Dans une justification avec \all{weil}, le verbe conjugué va à la fin : \all{..., weil die Verbindung teuer ist.}
\end{itemize}
\end{aretenir}

\coursec{Lexique thématique : les médias et la communication}

Après avoir répondu aux questions de compréhension détaillée sur le texte support, nous allons maintenant organiser et approfondir le vocabulaire qui permet de parler des smartphones, des cybercafés et de donner son avis sur les médias. Un bon lexique, c'est la moitié de la réussite en allemand : chaque mot appris avec son article et son pluriel devient un outil prêt à l'emploi pour la compréhension, la traduction et l'expression.

\courssub{Observer d'abord : un petit texte de réinvestissement}

Lisons d'abord ce texte original, qui réutilise naturellement tout le vocabulaire du chapitre. Soulignez mentalement les mots que vous avez déjà rencontrés dans le texte support.

\begin{textebox}[Smartphone oder Internetcafé ? — Ein Text zum Wortschatz]
In Yaoundé hat fast jeder Jugendliche ein Smartphone. Das Handy ist praktisch: Man kann telefonieren, chatten, E-Mails schreiben und im Internet surfen. Der Bildschirm ist klein, aber die Tastatur ist einfach. Viele Schüler schicken jeden Tag Nachrichten an ihre Freunde. „Ich chatte jeden Abend mit meinen Freunden“, sagt Marie, eine Schülerin der Première. „Das ist günstig und schnell.“

Aber das Smartphone hat auch Nachteile. Es ist oft teuer, und es kann vom Lernen ablenken. Manche Schüler spielen oder chatten in der Nacht, und am Morgen sind sie müde. Der Vater von Marie meint: „Das Handy ist nützlich, aber manchmal auch gefährlich.“

Nicht alle Familien haben ein Smartphone oder einen Computer zu Hause. Diese Schüler gehen ins Internetcafé. Dort gibt es Computer, einen Drucker und eine gute Verbindung. Eine Stunde Internet kostet nicht viel Geld. Im Internetcafé kann man Websites besuchen, Informationen suchen, Hausaufgaben machen und Dokumente drucken. Paul, ein Freund von Marie, findet das Internetcafé praktisch: „Die Verbindung ist hier schnell, und ich kann meine Arbeiten drucken.“

Aber manchmal ist die Verbindung langsam, besonders am Abend, wenn viele Leute online sind. Dann denkt Paul: „Das Internet ist heute wirklich langsam!“

Was ist besser — das Smartphone oder das Internetcafé? Marie glaubt: „Beide haben Vorteile und Nachteile. Das Smartphone ist mobil, das Internetcafé ist gut für die Schule.“ Der Grund ist einfach: Man braucht beide Medien. Wichtig ist nur: Die Medien dürfen nicht vom Lernen ablenken!
\end{textebox}

\textbf{Glossaire (mots difficiles) :}
\begin{itemize}
\item \all{das Handy, -s} : le portable ; \all{der Bildschirm, -e} : l'écran ; \all{die Tastatur, -en} : le clavier
\item \all{die Nachricht, -en} : le message, la nouvelle ; \all{die E-Mail, -s} : le courriel
\item \all{das Internetcafé, -s} : le cybercafé ; \all{der Drucker, die Drucker} : l'imprimante
\item \all{die Verbindung, -en} : la connexion ; \all{die Website, -s} : le site web
\item \all{ablenken (von + Dativ)} : distraire (de) ; \all{günstig} : avantageux, bon marché
\item \all{der Vorteil, -e} : l'avantage ; \all{der Nachteil, -e} : l'inconvénient
\item \all{müde} : fatigué ; \all{besuchen} : visiter ; \all{mobil} : mobile
\end{itemize}

\textbf{Questions de vérification du lexique :}
\begin{enumerate}
\item \all{Was kann man mit dem Smartphone machen?} (Que peut-on faire avec le smartphone ?)
\item \all{Warum gehen manche Schüler ins Internetcafé?} (Pourquoi certains élèves vont-ils au cybercafé ?)
\item \all{Welche Nachteile hat das Smartphone?} (Quels inconvénients le smartphone a-t-il ?)
\end{enumerate}

\courssub{Le champ lexical du smartphone}

\begin{definition}[Das Smartphone und seine Welt]
\all{das Smartphone, -s} : le smartphone ; mot anglais emprunté, toujours neutre. Synonyme très courant dans la langue parlée : \all{das Handy, -s} (le portable). Attention : en anglais, \emph{handy} signifie « pratique » — c'est un faux ami international !
\end{definition}

\begin{center}
\begin{tabularx}{\linewidth}{|X|X|X|}\hline
\rowcolor{popblueL} Deutsch & Français & Exemple \\ \hline
das Smartphone, -s & le smartphone & \all{Das Smartphone ist teuer.} \\ \hline
das Handy, -s & le portable & \all{Mein Handy ist kaputt.} \\ \hline
der Bildschirm, -e & l'écran & \all{Der Bildschirm ist klein.} \\ \hline
die Tastatur, -en & le clavier & \all{Die Tastatur ist einfach.} \\ \hline
die Nachricht, -en & le message & \all{Ich schicke eine Nachricht.} \\ \hline
die App, -s & l'application & \all{Diese App ist nützlich.} \\ \hline
chatten & chatter & \all{Wir chatten am Abend.} \\ \hline
telefonieren & téléphoner & \all{Sie telefoniert mit ihrer Mutter.} \\ \hline
\end{tabularx}
\end{center}

\begin{exemplebox}[Exemples traduits]
\all{Ich chatte jeden Abend mit meinen Freunden.} « Je chatte chaque soir avec mes amis. »

\all{Das Smartphone hat viele Vorteile, aber auch Nachteile.} « Le smartphone a beaucoup d'avantages, mais aussi des inconvénients. »

\all{Marie schickt ihrem Bruder eine Nachricht.} « Marie envoie un message à son frère. »
\end{exemplebox}

\begin{attention}[La conjugaison de « chatten » et « telefonieren »]
\all{chatten} se conjugue comme un verbe faible régulier : \all{ich chatte, du chattest, er/sie/es chattet, wir chatten, ihr chattet, sie chatten}. Au parfait : \all{ich habe gechattet}. Ne pas écrire « schatten » : le « ch » initial se prononce « tch », comme en anglais. De même \all{telefonieren} est faible : \all{ich habe telefoniert} (les verbes en \all{-ieren} ne prennent jamais \all{ge-} au participe passé).
\end{attention}

\courssub{Le champ lexical de l'internetcafé et d'internet}

\begin{definition}[Die Verbindung et verbinden]
\all{die Verbindung, -en} : la connexion, la liaison. Elle vient du verbe \all{verbinden} (relier, connecter ; verbe fort : \all{verbindet — verband — hat verbunden}). Exemple : \all{Die Verbindung ist heute sehr langsam.} « La connexion est très lente aujourd'hui. »
\end{definition}

\begin{center}
\begin{tabularx}{\linewidth}{|X|X|X|}\hline
\rowcolor{popblueL} Deutsch & Français & Exemple \\ \hline
das Internetcafé, -s & le cybercafé & \all{Paul geht ins Internetcafé.} \\ \hline
der Computer, die Computer & l'ordinateur & \all{Der Computer ist neu.} \\ \hline
das Internet (meist ohne Plural) & l'internet & \all{Das Internet ist schnell.} \\ \hline
die Verbindung, -en & la connexion & \all{Die Verbindung ist langsam.} \\ \hline
die Website, -s & le site web & \all{Diese Website ist interessant.} \\ \hline
die E-Mail, -s & le courriel & \all{Ich schreibe eine E-Mail.} \\ \hline
surfen & surfer & \all{Wir surfen im Internet.} \\ \hline
drucken & imprimer & \all{Er druckt seine Hausaufgaben.} \\ \hline
der Drucker, die Drucker & l'imprimante & \all{Der Drucker druckt schnell.} \\ \hline
\end{tabularx}
\end{center}

\begin{attention}[Faux amis et pièges de genre]
\begin{itemize}
\item \all{der Computer} est masculin, comme « l'ordinateur » : facile ! Mais \all{das Internet} est neutre, alors que « l'internet » est masculin en français.
\item \all{die E-Mail} est féminin en allemand (comme \all{die Post}), alors que « le mail » est masculin en français.
\item On dit \all{im Internet surfen} (surfer \emph{sur} internet) : la préposition \all{in} + datif, jamais \all{auf dem Internet}.
\end{itemize}
\end{attention}

\courssub{Le champ lexical de l'argumentation : avantages, inconvénients, opinions}

Pour débattre des médias, il faut un lexique d'évaluation : nommer les avantages et les inconvénients, donner des raisons et qualifier.

\begin{center}
\begin{tabularx}{\linewidth}{|X|X|X|}\hline
\rowcolor{popblueL} Deutsch & Français & Remarque \\ \hline
der Vorteil, -e & l'avantage & \all{Das Handy hat viele Vorteile.} \\ \hline
der Nachteil, -e & l'inconvénient & contraire de \all{der Vorteil} \\ \hline
der Grund, -¨e & la raison, le motif & pluriel : \all{die Gründe} \\ \hline
die Meinung, -en & l'avis, l'opinion & \all{meiner Meinung nach} : à mon avis \\ \hline
nützlich & utile & adjectif en \all{-lich} \\ \hline
praktisch & pratique & attention : pas « pratique » au sens d'exercice \\ \hline
teuer & cher & \all{Das Smartphone ist zu teuer.} \\ \hline
günstig & avantageux, bon marché & contraire de \all{teuer} \\ \hline
gefährlich & dangereux & de \all{die Gefahr} (le danger) \\ \hline
langsam / schnell & lent / rapide & \all{Die Verbindung ist schnell.} \\ \hline
\end{tabularx}
\end{center}

\begin{propriete}[Les verbes d'opinion : denken, meinen, glauben, finden]
Pour exprimer son avis, l'allemand utilise quatre verbes principaux :
\begin{itemize}
\item \all{denken} (penser, verbe fort : \all{denkt — dachte — hat gedacht}) : \all{Ich denke, das Handy ist nützlich.}
\item \all{meinen} (penser, estimer, verbe faible) : \all{Er meint, das Internetcafé ist praktisch.}
\item \all{glauben} (croire, verbe faible) : \all{Ich glaube, die Verbindung ist heute schnell.}
\item \all{finden} + adjectif (trouver, verbe fort : \all{findet — fand — hat gefunden}) : \all{Ich finde das Smartphone praktisch.} « Je trouve le smartphone pratique. »
\end{itemize}
Remarque : après \all{finden}, l'adjectif reste invariable et se place en fin de phrase : \all{Paul findet das Internetcafé günstig.}
\end{propriete}

\begin{exemplebox}[Donner son avis : modèles de phrases]
\all{Meiner Meinung nach ist das Handy sehr nützlich.} « À mon avis, le portable est très utile. »

\all{Ich finde das Internetcafé praktisch, aber manchmal langsam.} « Je trouve le cybercafé pratique, mais parfois lent. »

\all{Sie glaubt, dass das Smartphone vom Lernen ablenkt.} « Elle croit que le smartphone distrait des études. »
\end{exemplebox}

\courssub{Familles de mots : construire son vocabulaire en réseau}

L'allemand forme facilement des mots nouveaux à partir d'une racine. Apprendre les familles de mots, c'est multiplier son vocabulaire sans effort.

\begin{center}
\begin{tabularx}{\linewidth}{|X|X|X|}\hline
\rowcolor{popblueL} Verbe / racine & Mot dérivé & Français \\ \hline
verbinden (relier) & die Verbindung, -en & la connexion \\ \hline
nützen (servir) & nützlich / der Nutzen & utile / l'utilité \\ \hline
drucken (imprimer) & der Drucker, die Drucker & l'imprimante \\ \hline
informieren (informer) & die Information, -en & l'information \\ \hline
telefonieren (téléphoner) & das Telefon, -e & le téléphone \\ \hline
schicken (envoyer) & die Nachricht, -en & le message \\ \hline
\end{tabularx}
\end{center}

\begin{propriete}[Les suffixes utiles]
\begin{itemize}
\item \all{-ung} forme des noms féminins à partir de verbes : \all{verbinden → die Verbindung}, \all{meinen → die Meinung}. Tous les noms en \all{-ung} sont féminins, sans exception !
\item \all{-ion} (souvent via le latin) forme aussi des noms féminins : \all{informieren → die Information}.
\item \all{-lich} forme des adjectifs : \all{der Nutzen → nützlich}, \all{die Gefahr → gefährlich}.
\item \all{-er} désigne souvent l'appareil ou la personne qui fait l'action : \all{drucken → der Drucker}, \all{der Computer}.
\end{itemize}
\end{propriete}

\courssub{Expressions utiles à mémoriser}

\begin{aretenir}[Les expressions clés du chapitre]
\begin{itemize}
\item \all{eine Nachricht schicken} : envoyer un message — \all{Ich schicke meiner Freundin eine Nachricht.}
\item \all{im Internet surfen} : surfer sur internet — \all{Wir surfen im Internet.}
\item \all{eine Stunde kosten} : coûter une heure — \all{Eine Stunde Internet kostet 200 Franken.}
\item \all{vom Lernen ablenken} : distraire des études — \all{Das Handy lenkt viele Schüler vom Lernen ab.} (verbe à particule séparable : \all{lenkt ... ab} !)
\item \all{meiner Meinung nach} : à mon avis — \all{Meiner Meinung nach sind Smartphones nützlich.}
\end{itemize}
\end{aretenir}

\begin{attention}[Verbe à particule : ablenken]
\all{ablenken} est un verbe à particule séparable : à l'indicatif présent, la particule \all{ab-} va en fin de phrase : \all{Das Smartphone lenkt die Schüler vom Lernen ab.} Au parfait : \all{hat abgelenkt}. Ne dites jamais \all{Das Smartphone ablenkt die Schüler}.
\end{attention}

\courssub{Carte mentale du vocabulaire}

\begin{popfigure}
\begin{center}
\begin{tikzpicture}[mindmap, grow cyclic,
  every node/.style={concept, align=center, font=\small},
  root concept/.append style={concept color=popblue, text=white, font=\bfseries},
  level 1/.append style={level distance=3.6cm, sibling angle=120},
  level 2/.append style={level distance=2.4cm, sibling angle=45}]
\node[root concept]{Medien und Kommunikation}
  child[concept color=poporange]{ node {Smartphone}
    child { node[concept color=poporangeL, text=popdark] {das Handy} }
    child { node[concept color=poporangeL, text=popdark] {der Bildschirm} }
    child { node[concept color=poporangeL, text=popdark] {chatten} }
    child { node[concept color=poporangeL, text=popdark] {die Nachricht} }
  }
  child[concept color=popgreen]{ node {Internetcafé}
    child { node[concept color=popgreenL, text=popdark] {der Computer} }
    child { node[concept color=popgreenL, text=popdark] {die Verbindung} }
    child { node[concept color=popgreenL, text=popdark] {surfen, drucken} }
    child { node[concept color=popgreenL, text=popdark] {die Website} }
  }
  child[concept color=poppurple]{ node {Meinung ausdrücken}
    child { node[concept color=poppurpleL, text=popdark] {der Vorteil / der Nachteil} }
    child { node[concept color=poppurpleL, text=popdark] {nützlich, teuer, günstig} }
    child { node[concept color=poppurpleL, text=popdark] {denken, meinen, finden} }
  };
\end{tikzpicture}
\end{center}
\legende{Carte mentale du lexique : trois familles de mots pour parler des médias.}
\end{popfigure}

\courssub{Exercice résolu et entraînement}

\begin{exoresolu}[Complétez avec le mot qui convient]
Complétez les phrases avec : \all{die Verbindung, der Nachteil, chatten, drucken, nützlich, die Nachricht}.
\begin{enumerate}
\item \all{Im Internetcafé kann ich meine Hausaufgaben \dots}
\item \all{\dots\ ist heute sehr langsam, ich kann nicht surfen.}
\item \all{Ich schicke meinem Bruder eine \dots}
\item \all{Am Abend \dots\ viele Jugendliche mit ihren Freunden.}
\item \all{Das Handy ist \dots, aber es lenkt vom Lernen ab.}
\item \all{Ein \dots\ des Smartphones: Es ist sehr teuer.}
\end{enumerate}
\tcblower
\textbf{Solution détaillée :}
\begin{enumerate}
\item \all{drucken} — après \all{können}, l'infinitif se place en fin de phrase : « imprimer mes devoirs ».
\item \all{Die Verbindung} — sujet féminin avec majuscule : « la connexion est lente ».
\item \all{Nachricht} — \all{eine Nachricht schicken} : « envoyer un message » (accusatif).
\item \all{chatten} — sujet pluriel (\all{viele Jugendliche}), verbe faible : \all{sie chatten}.
\item \all{nützlich} — adjectif invariable après \all{sein} : « utile ».
\item \all{Nachteil} — \all{ein Nachteil} : masculin, « un inconvénient » (contraire : \all{ein Vorteil}).
\end{enumerate}
\end{exoresolu}

\begin{exercice}[À vous de jouer !]
\begin{enumerate}
\item Donnez l'article et le pluriel : Smartphone, Tastatur, Internetcafé, Vorteil, Meinung, Website.
\item Traduisez : a) « Je trouve le cybercafé pratique. » b) « Le portable distrait les élèves des études. » c) « À mon avis, la connexion est trop lente. »
\item Formez le nom à partir du verbe : \all{verbinden, informieren, meinen}. Quel est le genre de ces noms ?
\end{enumerate}
\end{exercice}

\begin{corrige}[Exercice « À vous de jouer ! »]
\begin{enumerate}
\item \all{das Smartphone, -s} ; \all{die Tastatur, -en} ; \all{das Internetcafé, -s} ; \all{der Vorteil, -e} ; \all{die Meinung, -en} ; \all{die Website, -s}.
\item a) \all{Ich finde das Internetcafé praktisch.} b) \all{Das Handy lenkt die Schüler vom Lernen ab.} c) \all{Meiner Meinung nach ist die Verbindung zu langsam.}
\item \all{die Verbindung} (suffixe \all{-ung}), \all{die Information} (suffixe \all{-ion}), \all{die Meinung} (suffixe \all{-ung}) : tous les trois sont féminins.
\end{enumerate}
\end{corrige}

\begin{methode}[Comment mémoriser durablement le vocabulaire]
\begin{enumerate}
\item Apprenez chaque nom \textbf{avec son article et son pluriel} : \all{das Smartphone, -s} ; \all{die Verbindung, -en}. L'article décide de toute la grammaire de la phrase !
\item Classez les mots en familles (carte mentale ci-dessus) : le cerveau retient mieux les réseaux que les listes.
\item Intégrez chaque mot dans \textbf{une phrase personnelle} sur votre vie : \all{Ich chatte jeden Abend mit meinen Freunden.}
\item Réutilisez le vocabulaire à l'oral : décrivez votre téléphone ou le cybercafé de votre quartier en trois phrases.
\end{enumerate}
\end{methode}

\begin{savaistu}[Le portable s'appelle « Natel » en Suisse]
En Suisse alémanique, on dit \all{das Natel} pour le portable (de \emph{Nationales Auto-Telefon}), et en Autriche on entend \all{die Handytasche} pour la housse du téléphone. Le vocabulaire des médias varie donc selon les pays germanophones — mais \all{das Handy} est compris partout ! Au Cameroun, beaucoup d'élèves découvrent internet d'abord au cybercafé : un point commun avec les années 1990 en Allemagne, où les \all{Internetcafés} étaient aussi très nombreux.
\end{savaistu}

\begin{aretenir}[L'essentiel du lexique]
Trois familles de mots : le smartphone (\all{das Handy, der Bildschirm, chatten, die Nachricht}), l'internetcafé (\all{der Computer, die Verbindung, surfen, drucken}) et l'argumentation (\all{der Vorteil, der Nachteil, nützlich, teuer, günstig} + \all{denken, meinen, glauben, finden}). Chaque nom s'apprend avec article et pluriel ; les noms en \all{-ung} sont toujours féminins ; \all{finden} + adjectif invariable permet de donner son avis : \all{Ich finde das Handy praktisch.}
\end{aretenir}

\coursec{Grammaire 1 : la comparaison (Komparativ und Superlativ)}

Dans la section précédente, nous avons enrichi notre vocabulaire sur le thème des médias et de la communication : \all{das Smartphone}, \all{das Internetcafé}, \all{chatten}, \all{die Verbindung}... Mais un vocabulaire ne sert pleinement que si l'on peut \textbf{comparer} les réalités qu'il désigne. Or, dans notre texte support, le narrateur compare constamment le smartphone et le cybercafé : lequel est plus pratique ? plus rapide ? le meilleur ? C'est exactement ce que permet la \cle{comparaison} (\all{der Vergleich}), objet de cette section de grammaire.

\courssub{Observation : que remarque-t-on dans le texte ?}

\begin{textebox}[Extrait du texte support — à observer]
\all{Das Smartphone ist praktischer als das Internetcafé. Man kann es überall benutzen: zu Hause, in der Schule, im Bus. Das Internetcafé ist günstiger, aber man muss dorthin gehen. Am Abend ist das Internet am langsamsten, weil viele Leute online sind. Für viele Schüler in Yaoundé ist das Smartphone das beste Medium für die Kommunikation. Chatten ist so interessant wie Telefonieren, aber oft billiger.}
\par\medskip
\emph{Glossaire :} \all{praktisch} (pratique) ; \all{günstig} (avantageux, bon marché) ; \all{benutzen} (utiliser) ; \all{langsam} (lent) ; \all{die Kommunikation} (la communication) ; \all{billig} (bon marché).
\end{textebox}

\textbf{Questions d'observation :}
\begin{enumerate}
\item Souligne dans l'extrait les adjectifs qui comparent deux réalités. Quelle terminaison portent-ils ?
\item Quel mot introduit le second terme de la comparaison après \all{praktischer} ?
\item Dans \all{am langsamsten} et \all{das beste Medium}, quelle forme particulière reconnais-tu ?
\item Dans \all{so interessant wie}, la comparaison exprime-t-elle une supériorité ou une égalité ?
\end{enumerate}

On observe trois constructions différentes :
\begin{itemize}
\item \all{praktischer als} : l'adjectif prend la terminaison \textbf{-er}, suivi de \all{als} ;
\item \all{am langsamsten}, \all{das beste} : l'adjectif prend la terminaison \textbf{-sten} après \all{am} ou après l'article ;
\item \all{so interessant wie} : l'adjectif reste à sa forme de base, encadré par \all{so ... wie}.
\end{itemize}

Ce sont les trois degrés de la comparaison : le \cle{positif} (\all{Positiv}), le \cle{comparatif} (\all{Komparativ}) et le \cle{superlatif} (\all{Superlativ}).

\begin{definition}[Les trois degrés de la comparaison]
\begin{itemize}
\item \textbf{Positiv} (degré positif) : forme de base de l'adjectif (\all{schnell}).
\item \textbf{Komparativ} (comparatif) : compare deux éléments (\all{schneller als}).
\item \textbf{Superlativ} (superlatif) : exprime le degré le plus élevé parmi un ensemble (\all{am schnellsten} / \all{das schnellste}).
\end{itemize}
\end{definition}

\courssub{Le comparatif : adjectif + -er ... als}

\begin{propriete}[Formation du comparatif]
Le comparatif se forme en ajoutant le suffixe \textbf{-er} à l'adjectif. Le second terme de la comparaison est introduit par \all{als} (= « que »).
\begin{center}
\all{Adjectif + -er + als}
\end{center}
\all{schnell} $\rightarrow$ \all{schneller als} (plus rapide que) ; \all{teuer} $\rightarrow$ \all{teurer als} (plus cher que).
\end{propriete}

\begin{exemplebox}[Exemples traduits]
\all{Das Smartphone ist praktischer als das Internetcafé.} « Le smartphone est plus pratique que le cybercafé. »
\par
\all{Das Internetcafé ist günstiger als das Smartphone.} « Le cybercafé est plus avantageux que le smartphone. »
\par
\all{Der Bus ist langsamer als das Motorrad.} « Le bus est plus lent que la moto. »
\par
\all{Chatten ist billiger als Telefonieren.} « Chatter est moins cher que téléphoner. »
\end{exemplebox}

\textbf{Remarques de forme :}
\begin{itemize}
\item Les adjectifs en \textbf{-er} comme \all{teuer} perdent le \emph{e} muet : \all{teuer} $\rightarrow$ \all{teurer} (et non \emph{teuerer}).
\item Les adjectifs en \textbf{-el} perdent aussi le \emph{e} muet : \all{dunkel} $\rightarrow$ \all{dunkler}.
\item L'adjectif comparatif \textbf{attributif} (devant le nom) prend en plus la désinence adjectivale selon le genre, le nombre et le cas : \all{ein schnelleres Handy} (un téléphone plus rapide), \all{das billigere Angebot} (l'offre la moins chère).
\end{itemize}

\courssub{Umlaut au comparatif : alt $\rightarrow$ älter}

\begin{propriete}[Umlaut des adjectifs monosyllabiques]
Les adjectifs \textbf{monosyllabiques} (une seule syllabe) dont la voyelle radicale est \textbf{a}, \textbf{o} ou \textbf{u} prennent généralement l'\textbf{Umlaut} (tréma) au comparatif et au superlatif :
\begin{center}
\all{alt} $\rightarrow$ \all{älter} $\rightarrow$ \all{am ältesten} \quad ; \quad \all{groß} $\rightarrow$ \all{größer} $\rightarrow$ \all{am größten}
\end{center}
Autres exemples courants : \all{jung} $\rightarrow$ \all{jünger} ; \all{lang} $\rightarrow$ \all{länger} ; \all{kurz} $\rightarrow$ \all{kürzer} ; \all{warm} $\rightarrow$ \all{wärmer} ; \all{stark} $\rightarrow$ \all{stärker} ; \all{kalt} $\rightarrow$ \all{kälter} ; \all{hart} $\rightarrow$ \all{härter}.
\end{propriete}

\begin{exemplebox}[Exemples traduits]
\all{Mein Bruder ist älter als ich.} « Mon frère est plus âgé que moi. »
\par
\all{Das Internetcafé in der Stadt ist größer als das Café im Viertel.} « Le cybercafé en ville est plus grand que le café du quartier. »
\par
\all{Im Dezember sind die Nächte länger als im Juni.} « En décembre, les nuits sont plus longues qu'en juin. »
\par
\all{Der Kaffee ist wärmer als der Tee.} « Le café est plus chaud que le thé. »
\end{exemplebox}

\begin{attention}[Exceptions notables — pas d'Umlaut]
Quelques adjectifs monosyllabiques en \emph{a}, \emph{o}, \emph{u} ne prennent \textbf{pas} l'Umlaut : \all{klar} $\rightarrow$ \all{klarer} ; \all{voll} $\rightarrow$ \all{voller} ; \all{stolz} $\rightarrow$ \all{stolzer} ; \all{flach} $\rightarrow$ \all{flacher}. 
\par
Les adjectifs de plus d'une syllabe (\all{praktisch}, \all{interessant}, \all{günstig}) ne prennent \textbf{jamais} l'Umlaut : \all{praktischer} (jamais \emph{präktischer}), \all{günstiger} (jamais \emph{günstiger} avec Umlaut sur le ü qui existe déjà).
\end{attention}

\courssub{Le superlatif : am ... -sten et der/die/das ... -ste}

\begin{propriete}[Formation du superlatif]
Deux constructions selon la fonction syntaxique de l'adjectif :
\begin{enumerate}
\item \textbf{Adjectif attribut (prédicatif), après le verbe \all{sein} / \all{werden} / \all{bleiben}} : \all{am + adjectif + -sten}.\\
\all{Das Internet ist am Abend am langsamsten.} « Internet est le plus lent le soir. »
\item \textbf{Adjectif épithète (attributif), devant un nom} : \all{der/die/das + adjectif + -ste} + désinence adjectivale + nom.\\
\all{Das ist das schnellste Internet der Stadt.} « C'est l'Internet le plus rapide de la ville. »
\end{enumerate}
\emph{Règle de prononciation :} après les adjectifs terminés par \textbf{-d, -t, -s, -ß, -z, -sch}, on ajoute \textbf{-esten} (au lieu de -sten) pour faciliter la prononciation : \all{alt} $\rightarrow$ \all{am ältesten} ; \all{heiß} $\rightarrow$ \all{am heißesten} ; \all{kurz} $\rightarrow$ \all{am kürzesten} ; \all{groß} $\rightarrow$ \all{am größten}.
\end{propriete}

\begin{exemplebox}[Exemples traduits]
\all{Der Bendskin ist am schnellsten im Verkehr von Douala.} « Le moto-taxi est le plus rapide dans la circulation de Douala. »
\par
\all{Die jüngste Schülerin der Klasse hat das beste Smartphone.} « L'élève la plus jeune de la classe a le meilleur smartphone. »
\par
\all{Samstag ist der Markt am vollsten.} « Le samedi, le marché est le plus plein. »
\par
\all{Das ist das teuerste Handy im Laden.} « C'est le portable le plus cher du magasin. »
\end{exemplebox}

\begin{popfigure}
\begin{center}
\begin{tikzpicture}
\node[draw=popblue, fill=popblueL, rounded corners, minimum width=3.4cm, minimum height=1.1cm, align=center, font=\bfseries] (pos) at (0,0) {Positiv\\ \all{schnell}};
\node[draw=poporange, fill=poporangeL, rounded corners, minimum width=3.4cm, minimum height=1.1cm, align=center, font=\bfseries] (komp) at (5.2,0) {Komparativ\\ \all{schneller als}};
\node[draw=popgreen, fill=popgreenL, rounded corners, minimum width=3.4cm, minimum height=1.1cm, align=center, font=\bfseries] (sup) at (10.4,0) {Superlativ\\ \all{am schnellsten}};
\draw[-{Stealth}, very thick, popdark] (pos) -- (komp) node[midway, above, font=\small] {+ \texttt{-er}};
\draw[-{Stealth}, very thick, popdark] (komp) -- (sup) node[midway, above, font=\small] {\all{am} + \texttt{-sten}};
\node[align=center, font=\small, text=popmuted] at (5.2,-1.3) {\all{teuer} $\rightarrow$ \all{teurer} $\rightarrow$ \all{am teuersten} \quad ; \quad \all{groß} $\rightarrow$ \all{größer} $\rightarrow$ \all{am größten}};
\end{tikzpicture}
\end{center}
\legende{Les trois degrés de la comparaison : Positiv, Komparativ, Superlativ.}
\end{popfigure}

\courssub{Comparaison d'égalité et d'inégalité}

\begin{propriete}[Égalité : so ... wie ; inégalité : -er ... als]
\begin{itemize}
\item \textbf{Égalité} (« aussi... que ») : \all{so + adjectif (positif) + wie}.\\
\all{Das Handy ist so teuer wie der Computer.} « Le portable est aussi cher que l'ordinateur. »
\item \textbf{Inégalité / Supériorité} (« plus... que ») : \all{comparatif + als}.\\
\all{Das Smartphone ist schneller als das Internetcafé.} « Le smartphone est plus rapide que le cybercafé. »
\item \textbf{Négation de l'égalité} (« pas aussi... que ») : \all{nicht so ... wie}.\\
\all{Das Internetcafé ist nicht so praktisch wie das Smartphone.} « Le cybercafé n'est pas aussi pratique que le smartphone. »
\end{itemize}
\end{propriete}

\begin{popfigure}
\begin{center}
\begin{tikzpicture}
\node[draw=popgreen, fill=popgreenL, rounded corners, align=center, minimum width=5.6cm, minimum height=1.6cm, font=\small] (egal) at (0,0) {\textbf{ÉGALITÉ}\\ \all{so + Adjektiv + wie}\\ \all{so interessant wie}};
\node[draw=poporange, fill=poporangeL, rounded corners, align=center, minimum width=5.6cm, minimum height=1.6cm, font=\small] (ineg) at (8,0) {\textbf{INÉGALITÉ}\\ \all{Adjektiv + -er + als}\\ \all{interessanter als}};
\node[draw=popblue, fill=popblueL, rounded corners, align=center, minimum width=4.2cm, minimum height=1cm, font=\small] (adj) at (4,2.4) {\textbf{ADJECTIF (Positiv)}\\ \all{interessant}};
\draw[-{Stealth}, thick, popdark] (adj.south west) -- (egal.north);
\draw[-{Stealth}, thick, popdark] (adj.south east) -- (ineg.north);
\end{tikzpicture}
\end{center}
\legende{Deux constructions à partir de l'adjectif : égalité (\all{so ... wie}) ou inégalité (\all{-er ... als}).}
\end{popfigure}

\courssub{Les adjectifs et adverbes irréguliers}

\begin{aretenir}[Tableau des formes irrégulières — à apprendre par cœur]
\begin{center}
\begin{tabular}{|l|l|l|l|}
\hline
\rowcolor{popblueL} \textbf{Positiv} & \textbf{Komparativ} & \textbf{Superlativ} & \textbf{Français} \\
\hline
\all{gut} & \all{besser} & \all{am besten} & bon / meilleur / le meilleur \\
\hline
\all{gern} & \all{lieber} & \all{am liebsten} & volontiers / plutôt / de préférence \\
\hline
\all{viel} & \all{mehr} & \all{am meisten} & beaucoup / plus / le plus \\
\hline
\all{hoch} & \all{höher} & \all{am höchsten} & haut / plus haut / le plus haut \\
\hline
\all{nah} & \all{näher} & \all{am nächsten} & proche / plus proche / le plus proche \\
\hline
\all{bald} & \all{eher} & \all{am ehesten} & bientôt / plus tôt / le plus tôt \\
\hline
\end{tabular}
\end{center}
\end{aretenir}

\begin{exemplebox}[Exemples traduits — formes irrégulières]
\all{Ich gehe lieber ins Internetcafé.} « Je préfère aller au cybercafé. »
\par
\all{Das Smartphone ist das beste Medium für die Kommunikation.} « Le smartphone est le meilleur média pour la communication. »
\par
\all{Im Internetcafé bezahlt man weniger, aber das Handy benutzt man mehr.} « Au cybercafé, on paie moins, mais on utilise davantage le portable. »
\par
\all{Das Internetcafé ist näher als die Bibliothek.} « Le cybercafé est plus proche que la bibliothèque. »
\par
\all{Er kommt eher als erwartet.} « Il arrive plus tôt que prévu. »
\end{exemplebox}

\courssub{Tableau récapitulatif et comparaison français/allemand}

\begin{center}
\begin{tabular}{|l|l|l|l|}
\hline
\rowcolor{popblueL} \textbf{Positiv} & \textbf{Komparativ} & \textbf{Superlativ (prédicatif)} & \textbf{Remarque} \\
\hline
\all{schnell} & \all{schneller} & \all{am schnellsten} & régulier \\
\hline
\all{teuer} & \all{teurer} & \all{am teuersten} & perte du \emph{e} muet \\
\hline
\all{groß} & \all{größer} & \all{am größten} & Umlaut \\
\hline
\all{alt} & \all{älter} & \all{am ältesten} & Umlaut + \texttt{-esten} \\
\hline
\all{gut} & \all{besser} & \all{am besten} & irrégulier \\
\hline
\all{gern} & \all{lieber} & \all{am liebsten} & irrégulier (adverbe) \\
\hline
\end{tabular}
\end{center}

\begin{center}
\begin{tabularx}{\linewidth}{|X|X|X|}
\hline
\rowcolor{popblueL} \textbf{Français} & \textbf{Deutsch} & \textbf{Remarque} \\
\hline
plus rapide que & \all{schneller als} & « que » = \all{als} après un comparatif \\
\hline
aussi cher que & \all{so teuer wie} & « aussi... que » = \all{so ... wie} \\
\hline
le plus rapide (prédicatif) & \all{am schnellsten} & \all{am} + \texttt{-sten} (après \all{sein}) \\
\hline
le meilleur média (attributif) & \all{das beste Medium} & article + \texttt{-ste} + désinence devant le nom \\
\hline
pas aussi pratique que & \all{nicht so praktisch wie} & négation de l'égalité \\
\hline
\end{tabularx}
\end{center}

\begin{attention}[Piège n° 1 : « als » et non « wie » après un comparatif]
Erreur fréquente des francophones : traduire « que » par \all{wie} après un comparatif. On dit \all{besser ALS}, jamais \emph{besser wie}. \all{wie} ne s'emploie qu'avec \all{so} (égalité) : \all{so gut wie}. Retiens : \textbf{comparatif $\rightarrow$ als ; égalité $\rightarrow$ so ... wie}.
\end{attention}

\begin{attention}[Piège n° 2 : jamais « mehr » + comparatif]
En français on dit « plus rapide », mais en allemand le « plus » est déjà contenu dans le suffixe \textbf{-er}. On ne dit donc jamais \emph{mehr schneller} : la bonne forme est \all{schneller}. De même : \all{besser} (et non \emph{mehr besser}), \all{größer} (et non \emph{mehr größer}). \all{mehr} est lui-même le comparatif de \all{viel} : \all{Ich habe mehr Zeit als du.} « J'ai plus de temps que toi. »
\end{attention}

\courssub{Exercice résolu et entraînement}

\begin{exoresolu}[Mets l'adjectif au degré demandé]
Énoncé : complète les phrases avec la forme correcte de l'adjectif entre parenthèses.
\begin{enumerate}
\item \all{Das Smartphone ist \dots\ als der Brief. (schnell)}
\item \all{Das Internetcafé ist \dots\ als das Smartphone. (günstig)}
\item \all{Mein Vater ist \dots\ als mein Bruder. (alt)}
\item \all{Für mich ist Chatten das \dots\ Medium. (gut)}
\item \all{Der Computer ist \dots\ teuer \dots\ das Handy. (so)}
\item \all{Am Abend ist die Verbindung am \dots\ . (langsam)}
\end{enumerate}
\tcblower
Solution détaillée :
\begin{enumerate}
\item \all{schneller} : comparatif régulier, adjectif + \texttt{-er}, suivi de \all{als}. « Le smartphone est plus rapide que la lettre. »
\item \all{günstiger} : comparatif régulier (perte du \emph{e} muet). « Le cybercafé est plus avantageux que le smartphone. »
\item \all{älter} : adjectif monosyllabique en \emph{a} $\rightarrow$ Umlaut. « Mon père est plus âgé que mon frère. »
\item \all{beste} : superlatif irrégulier de \all{gut}, attributif devant le nom neutre \all{das Medium} $\rightarrow$ \all{das beste} (désinence \texttt{-e} après article défini). « Pour moi, chatter est le meilleur média. »
\item \all{so ... wie} : comparaison d'égalité. \all{Der Computer ist so teuer wie das Handy.} « L'ordinateur est aussi cher que le portable. »
\item \all{langsamsten} : superlatif prédicatif après \all{am} $\rightarrow$ \all{am langsamsten}. « Le soir, la connexion est la plus lente. »
\end{enumerate}
\end{exoresolu}

\begin{exercice}[À toi de comparer !]
\begin{enumerate}
\item Traduis en allemand : 
      a) « Le smartphone est plus pratique que le cybercafé. » 
      b) « Chatter est aussi intéressant que téléphoner. » 
      c) « Internet est le plus lent le soir. »
\item Forme le comparatif et le superlatif (prédicatif \all{am ... -sten}) de : \all{billig}, \all{jung}, \all{hoch}, \all{viel}, \all{gut}.
\item Rédige trois phrases pour comparer le smartphone et le cybercafé de ton quartier, en utilisant au moins un comparatif et un superlatif.
\end{enumerate}
\end{exercice}

\begin{corrige}[Exercice « À toi de comparer ! »]
1. a) \all{Das Smartphone ist praktischer als das Internetcafé.} 
   b) \all{Chatten ist so interessant wie Telefonieren.} 
   c) \all{Das Internet ist am Abend am langsamsten.}
\par
2. \all{billig $\rightarrow$ billiger $\rightarrow$ am billigsten} ; 
   \all{jung $\rightarrow$ jünger $\rightarrow$ am jüngsten} ; 
   \all{hoch $\rightarrow$ höher $\rightarrow$ am höchsten} ; 
   \all{viel $\rightarrow$ mehr $\rightarrow$ am meisten} ; 
   \all{gut $\rightarrow$ besser $\rightarrow$ am besten}.
\par
3. Modèle : \all{Das Smartphone ist schneller als das Internetcafé, aber das Internetcafé ist billiger. Für mich ist das Smartphone das praktischste Medium.}
\end{corrige}

\begin{savaistu}[La comparaison au Bac]
Au baccalauréat, la comparaison est très souvent demandée dans les questions de grammaire ou de production : une consigne comme \all{„Vergleichen Sie Smartphone und Internetcafé“} exige au moins \textbf{deux comparatifs corrects} (par exemple \all{praktischer als} et \all{billiger als}) et, idéalement, un superlatif (\all{das beste Medium} ou \all{am praktischsten}). Maîtriser \all{-er ... als}, \all{so ... wie} et \all{am ... -sten}, c'est donc gagner des points sûrs le jour de l'examen !
\end{savaistu}

En résumé : le comparatif se forme avec \textbf{-er} suivi de \all{als}, le superlatif avec \all{am} + \textbf{-sten} (prédicatif) ou article + \textbf{-ste} + désinence (attributif), l'égalité avec \all{so ... wie}. Retiens les irréguliers \all{gut $\rightarrow$ besser $\rightarrow$ am besten} et \all{gern $\rightarrow$ lieber $\rightarrow$ am liebsten}, et méfie-toi des deux pièges : jamais \emph{mehr schneller}, jamais \emph{besser wie}. Dans la section suivante, nous verrons comment exprimer ses opinions avec les subordonnées en \all{dass} : \all{Ich denke, dass das Smartphone praktischer ist.}

\coursec{Grammaire 2 : les subordonnées en « dass » et le parfait des verbes forts}

Dans la section précédente, nous avons appris à comparer : \all{Das Smartphone ist praktischer als das Internetcafé.} (Le smartphone est plus pratique que le cybercafé.) Mais pour exprimer une opinion complète sur les médias, il ne suffit pas de comparer : il faut aussi savoir dire \emph{ce que l'on pense}, \emph{ce que l'on croit}, \emph{ce que l'on dit}. En allemand, cela passe par une construction très particulière : la subordonnée en \all{dass}, où le verbe conjugué « migre » à la fin de la phrase. Nous verrons ensuite comment raconter au passé une expérience vécue avec les médias grâce au parfait des verbes forts : \all{Ich habe eine E-Mail geschrieben.} (J'ai écrit un courriel.)

\courssub{Observation : où est passé le verbe ?}

Lisons attentivement ces phrases tirées de notre thème :

\begin{exemplebox}[Phrases d'observation]
\all{Ich denke, dass das Smartphone besser ist.}\\
« Je pense que le smartphone est meilleur. »

\all{Sie meint, dass die Verbindung zu langsam ist.}\\
« Elle estime que la connexion est trop lente. »

\all{Wir glauben, dass das Internetcafé viele Vorteile hat.}\\
« Nous croyons que le cybercafé présente beaucoup d'avantages. »

\all{Der Lehrer sagt, dass die Schüler zu viel chatten.}\\
« Le professeur dit que les élèves tchattent trop. »
\end{exemplebox}

Observez la place du verbe conjugué (\all{ist}, \all{hat}, \all{chatten}) : dans la première partie de la phrase (proposition principale), il occupe la deuxième position, comme toujours en allemand. Mais après \all{dass}, il se retrouve tout à la fin de la phrase. C'est la grande différence avec le français : après « que », le français garde l'ordre normal (sujet + verbe), tandis que l'allemand envoie le verbe conjugué en dernière position.

\begin{popfigure}
\begin{tikzpicture}[
  node distance=8mm and 4mm,
  block/.style={draw=popblue, fill=popblueL, rounded corners, align=center, font=\small, inner sep=4pt},
  subblock/.style={draw=poporange, fill=poporangeL, rounded corners, align=center, font=\small, inner sep=4pt},
  verb/.style={draw=poppink, fill=poppinkL, rounded corners, align=center, font=\small\bfseries, inner sep=4pt},
  lbl/.style={font=\scriptsize\itshape, text=popmuted, anchor=south}
]
\node[block] (main) {Ich denke,};
\node[lbl] at (main.north) {proposition principale};
\node[subblock, right=of main] (dass) {dass};
\node[subblock, right=of dass] (subj) {das Smartphone};
\node[lbl] at (subj.north) {sujet};
\node[subblock, right=of subj] (rest) {besser};
\node[verb, right=of rest] (verb) {ist.};
\node[lbl] at (verb.north) {verbe conjugué};
\draw[-{Stealth[length=3mm]}, very thick, poppink] (dass.south) .. controls +(0,-1.5) and +(0,-1.5) .. (verb.south);
\node[lbl, text=poppink] at ($(dass)!0.5!(verb) - (0,1.8)$) {le verbe migre en position finale};
\end{tikzpicture}
\legende{Structure de la subordonnée en « dass » : le verbe conjugué quitte la deuxième place et migre à la fin.}
\end{popfigure}

\courssub{La règle : la subordonnée en « dass »}

\begin{propriete}[Le verbe conjugué à la fin]
Dans une subordonnée introduite par \all{dass} (que), le verbe conjugué occupe la \cle{dernière place} de la phrase. La construction est :

proposition principale + \textbf{virgule} + \all{dass} + sujet + autres éléments + \textbf{verbe conjugué en fin}.

\all{Ich glaube, dass das Internetcafé viele Vorteile hat.}\\
« Je crois que le cybercafé a beaucoup d'avantages. »
\end{propriete}

Trois points sont à retenir absolument :

\begin{itemize}
\item La \cle{virgule} est obligatoire en allemand avant \all{dass}, même si le français n'en met pas devant « que ».
\item Le verbe conjugué se place \emph{après} tous les autres éléments de la subordonnée (compléments de temps, de lieu, d'objet).
\item Si la subordonnée contient un auxiliaire et un participe (perfect, futur, passif), l'auxiliaire conjugué vient en tout dernier : \all{Ich weiß, dass er eine E-Mail geschrieben hat.} (Je sais qu'il a écrit un courriel.)
\end{itemize}

\textbf{Emplois : après les verbes d'opinion et de déclaration.}\par
La subordonnée en \all{dass} s'emploie surtout après les verbes qui expriment une pensée, une opinion, une parole ou une connaissance :

\begin{center}
\begin{tabularx}{\linewidth}{|l|X|X|}\hline
\rowcolor{popblueL} \textbf{Verbe} & \textbf{Exemple allemand} & \textbf{Traduction} \\ \hline
\all{denken} (penser) & \all{Ich denke, dass das Smartphone praktischer ist.} & Je pense que le smartphone est plus pratique. \\ \hline
\all{meinen} (estimer) & \all{Sie meint, dass die Verbindung zu langsam ist.} & Elle estime que la connexion est trop lente. \\ \hline
\all{glauben} (croire) & \all{Wir glauben, dass das Internet nützlich ist.} & Nous croyons qu'internet est utile. \\ \hline
\all{sagen} (dire) & \all{Er sagt, dass er oft chattet.} & Il dit qu'il tchatte souvent. \\ \hline
\all{hoffen} (espérer) & \all{Ich hoffe, dass das Internetcafé bald öffnet.} & J'espère que le cybercafé ouvre bientôt. \\ \hline
\all{wissen} (savoir) & \all{Sie weiß, dass der Computer kaputt ist.} & Elle sait que l'ordinateur est en panne. \\ \hline
\all{finden} (trouver/avoir l'avis) & \all{Ich finde, dass Smartphones zu teuer sind.} & Je trouve que les smartphones sont trop chers. \\ \hline
\end{tabularx}
\end{center}

L'expression \all{ich finde, dass...} (je trouve que...) est très utile pour donner son avis dans un débat sur les avantages et les inconvénients des médias.

\begin{attention}[Le piège du calque français]
L'erreur la plus fréquente des francophones est de laisser le verbe en deuxième position après \all{dass}, par calque du français :

\emph{Faux} : \all{Ich denke, dass das Smartphone \textbf{ist} besser.}\\
\emph{Correct} : \all{Ich denke, dass das Smartphone besser \textbf{ist}.}

En français, on dit « je pense que le smartphone \emph{est} meilleur » : le verbe reste à sa place habituelle. En allemand, il doit migrer à la fin. Réflexe à acquérir : dès que vous écrivez \textbf{virgule + dass}, vérifiez que le verbe conjugué est bien en dernière position.
\end{attention}

\begin{attention}[Ne pas confondre « dass » et « das »]
\all{dass} (avec \textbf{ss}) est la conjonction « que » ; \all{das} (avec un seul \textbf{s}) est l'article neutre ou le pronom démonstratif. Test simple : si vous pouvez remplacer le mot par \all{dieses} (ce/cet), c'est \all{das} ; sinon, c'est \all{dass}.

\all{Das Smartphone ist teuer.} (Le smartphone est cher. → \all{dieses Smartphone} : possible, donc \all{das}.)\\
\all{Ich weiß, dass das Smartphone teuer ist.} (Je sais que le smartphone est cher. → le premier « que » ne peut pas être remplacé : c'est \all{dass}.)
\end{attention}

\begin{savaistu}[Pourquoi deux « s » ?]
Avant la réforme orthographique de 1996, on écrivait \all{daß} avec un \all{ß}. Depuis la réforme, la règle est simple : après une voyelle courte, on écrit \all{ss} (\all{dass}, \all{muss}), et après une voyelle longue ou une diphtongue, on garde \all{ß} (\all{Straße}, \all{heißen}). Vous rencontrerez encore l'ancienne graphie \all{daß} dans les livres anciens : ne vous étonnez pas, le sens est identique.
\end{savaistu}

\courssub{Le parfait des verbes forts}

Pour raconter ce que vous avez fait hier au cybercafé ou avec votre smartphone, on utilise le \cle{Perfekt} (parfait), le temps du passé le plus fréquent à l'oral. Vous connaissez déjà le parfait des verbes faibles : \all{chatten → hat gechattet}. Les verbes forts, eux, changent de voyelle (alternance apophonique) et prennent la terminaison \all{-en} au participe II.

\begin{propriete}[Formation du parfait des verbes forts]
Parfait = auxiliaire \all{haben} ou \all{sein} (conjugué en 2\textsuperscript{e} position) + participe II en \cle{ge-...-en} (à la fin), souvent avec changement de voyelle :

\all{gehen → ist gegangen} ; \all{schreiben → hat geschrieben} ; \all{lesen → hat gelesen} ; \all{finden → hat gefunden} ; \all{nehmen → hat genommen}.
\end{propriete}

\begin{exemplebox}[Le parfait dans le contexte des médias]
\all{Gestern habe ich eine Nachricht geschrieben.}\\
« Hier, j'ai écrit un message. »

\all{Wir sind ins Internetcafé gegangen.}\\
« Nous sommes allés au cybercafé. »

\all{Er hat die Information im Internet gefunden.}\\
« Il a trouvé l'information sur internet. »

\all{Sie hat den ganzen Artikel gelesen.}\\
« Elle a lu tout l'article. »

\all{Ich habe mein Handy genommen und ein Foto gemacht.}\\
« J'ai pris mon portable et j'ai fait une photo. »
\end{exemplebox}

\textbf{Le choix de l'auxiliaire : haben ou sein ?}\par

\begin{propriete}[Auxiliaire « sein » : mouvement et changement d'état]
On utilise \all{sein} au parfait avec les verbes de \cle{déplacement} (qui indiquent un mouvement d'un point A vers un point B) et les verbes de \cle{changement d'état} : \all{gehen} (aller), \all{kommen} (venir), \all{fahren} (aller en véhicule), \all{bleiben} (rester), \all{aufstehen} (se lever), \all{passieren} (se passer/arriver). Tous les autres verbes prennent \all{haben}.
\end{propriete}

\begin{attention}[Haben ou sein : le réflexe à acquérir]
\all{sein} pour les verbes de déplacement : \all{gehen, kommen, fahren} → \all{Wir \textbf{sind} ins Internetcafé gegangen.}\\
\all{haben} pour la grande majorité des autres verbes : \all{schreiben, lesen, chatten, finden} → \all{Ich \textbf{habe} eine Nachricht geschrieben.}

Erreur typique du francophone : dire \all{ich habe gegangen} sous l'influence du français « j'ai allé ». En allemand, tout déplacement exige \all{sein} : \all{ich \textbf{bin} gegangen}.
\end{attention}

\textbf{Les verbes forts sans préfixe ge-.}\par
Deux catégories de verbes forts forment leur participe II \emph{sans} le préfixe \all{ge-} :

\begin{itemize}
\item Les verbes en \all{-ieren} (souvent d'origine étrangère) : \all{studieren → hat studiert} (étudier), \all{reparieren → hat repariert} (réparer), \all{telefonieren → hat telefoniert} (téléphoner).
\item Les verbes à préfixe inséparable (\all{ver-, be-, er-, ent-, ge-, zer-, miss-}) : \all{verlieren → hat verloren} (perdre), \all{verstehen → hat verstanden} (comprendre), \all{bekommen → hat bekommen} (recevoir), \all{erzählen → hat erzählt} (raconter).
\end{itemize}

\all{Er hat sein Handy verloren.} (Il a perdu son portable.) — \all{Ich habe die Nachricht nicht verstanden.} (Je n'ai pas compris le message.)

\begin{popfigure}
\begin{center}
\begin{tabular}{|l|l|l|l|}\hline
\rowcolor{popblueL} \textbf{Infinitiv} & \textbf{Präteritum} & \textbf{Perfekt} & \textbf{Aux.} \\ \hline
\all{schreiben} & \all{schrieb} & \all{hat geschrieben} & \all{haben} \\ \hline
\all{lesen} & \all{las} & \all{hat gelesen} & \all{haben} \\ \hline
\all{finden} & \all{fand} & \all{hat gefunden} & \all{haben} \\ \hline
\all{gehen} & \all{ging} & \all{\textcolor{poppink}{ist gegangen}} & \all{sein} \\ \hline
\all{nehmen} & \all{nahm} & \all{hat genommen} & \all{haben} \\ \hline
\all{verlieren} & \all{verlor} & \all{hat verloren} & \all{haben} \\ \hline
\all{kommen} & \all{kam} & \all{\textcolor{poppink}{ist gekommen}} & \all{sein} \\ \hline
\all{verstehen} & \all{verstand} & \all{hat verstanden} & \all{haben} \\ \hline
\end{tabular}
\end{center}
\legende{Huit verbes forts du thème : infinitif, prétérit, parfait. En rose : auxiliaire « sein » (verbes de mouvement).}
\end{popfigure}

\courssub{Texte de réinvestissement}

Lisons maintenant un court texte qui réutilise les subordonnées en \all{dass} et le parfait des verbes forts dans notre thème.

\begin{textebox}[Ein Tag ohne Smartphone]
Gestern hat Marie ihr Handy verloren. Zuerst war sie traurig, aber dann ist sie ins Internetcafé gegangen. Dort hat sie eine E-Mail an ihre Freundin geschrieben und hat die Nachrichten im Internet gelesen. „Ich finde, dass das Internetcafé sehr praktisch ist“, sagt sie. „Viele Leute denken, dass man ohne Smartphone nicht leben kann, aber ich glaube, dass das nicht stimmt.“ Marie hat auch einen interessanten Artikel über Medien gefunden. Sie meint, dass das Internet viele Vorteile hat, aber sie weiß auch, dass zu viel Internet gefährlich sein kann. Am Abend ist sie nach Hause gegangen und hat ihr Handy unter dem Bett gefunden! „Ich hoffe, dass ich es nie wieder verliere“, sagt sie und lacht.
\end{textebox}

\textbf{Glossaire :} \all{verlieren, verlor, hat verloren} (perdre) ; \all{traurig} (triste) ; \all{zuerst} (d'abord) ; \all{die Nachricht, -en} (le message, la nouvelle) ; \all{stimmen} (être exact, être vrai) ; \all{der Artikel, -} (l'article) ; \all{gefährlich} (dangereux) ; \all{das Bett, -en} (le lit) ; \all{lachen} (rire) ; \all{nie wieder} (jamais plus).

\textbf{Questions de compréhension :}
\begin{enumerate}
\item \all{Was hat Marie verloren?} (Qu'est-ce que Marie a perdu ?)
\item \all{Wo ist sie gegangen?} (Où est-elle allée ?)
\item \all{Was denken viele Leute?} (Que pensent beaucoup de gens ?)
\item \all{Wo hat sie ihr Handy gefunden?} (Où a-t-elle trouvé son portable ?)
\end{enumerate}

\courssub{Exercices résolus et entraînement}

\begin{exoresolu}[Mettez les phrases dans le bon ordre]
Énoncé : remettez les mots dans l'ordre correct pour former des phrases correctes.\\
a) \all{denke / dass / ich / praktisch / das Smartphone / ist / sehr}\\
b) \all{meint / dass / sie / die Verbindung / langsam / zu / ist}\\
c) \all{glauben / dass / wir / viele Vorteile / das Internetcafé / hat}
\tcblower
Solution : après \all{dass}, le verbe conjugué va à la fin.\\
a) \all{Ich denke, dass das Smartphone sehr praktisch ist.} (Je pense que le smartphone est très pratique.)\\
b) \all{Sie meint, dass die Verbindung zu langsam ist.} (Elle estime que la connexion est trop lente.)\\
c) \all{Wir glauben, dass das Internetcafé viele Vorteile hat.} (Nous croyons que le cybercafé a beaucoup d'avantages.)
\end{exoresolu}

\begin{exoresolu}[Mettez au parfait]
Énoncé : conjuguez les verbes entre parenthèses au parfait. Attention au choix de l'auxiliaire et à la forme du participe.\\
a) \all{Ich (schreiben) eine E-Mail.}\\
b) \all{Wir (gehen) ins Internetcafé.}\\
c) \all{Er (finden) die Information im Internet.}\\
d) \all{Sie (verlieren) ihr Handy.}
\tcblower
Solution :\\
a) \all{Ich habe eine E-Mail geschrieben.} (verbe fort \all{schreiben} → \all{geschrieben}, auxiliaire \all{haben})\\
b) \all{Wir sind ins Internetcafé gegangen.} — mouvement → auxiliaire \all{sein}. (Nous sommes allés au cybercafé.)\\
c) \all{Er hat die Information im Internet gefunden.} (verbe fort \all{finden} → \all{gefunden}, auxiliaire \all{haben})\\
d) \all{Sie hat ihr Handy verloren.} — préfixe inséparable \all{ver-} : pas de \all{ge-}. (Elle a perdu son portable.)
\end{exoresolu}

\begin{exercice}[À vous de jouer]
\begin{enumerate}
\item Traduisez en allemand :\\
a) Je pense que le smartphone est trop cher.\\
b) Elle dit que le cybercafé est fermé aujourd'hui.\\
c) Nous espérons que la connexion est meilleure demain.
\item Mettez au parfait :\\
a) \all{Er (lesen) den Artikel.}\\
b) \all{Ich (nehmen) mein Handy.}\\
c) \all{Wir (kommen) ins Internetcafé.}\\
d) \all{Sie (verstehen) die Nachricht nicht.}
\item Choisissez \all{das} ou \all{dass} :\\
a) \all{Ich weiß, ... Smartphone teuer ist.}\\
b) \all{... Internetcafé ist heute geschlossen.}
\end{enumerate}
\end{exercice}

\begin{corrige}[Exercice « À vous de jouer »]
1. a) \all{Ich denke, dass das Smartphone zu teuer ist.} b) \all{Sie sagt, dass das Internetcafé heute geschlossen ist.} c) \all{Wir hoffen, dass die Verbindung morgen besser ist.}\\
2. a) \all{Er hat den Artikel gelesen.} b) \all{Ich habe mein Handy genommen.} c) \all{Wir sind ins Internetcafé gekommen.} (mouvement → \all{sein}) d) \all{Sie hat die Nachricht nicht verstanden.} (préfixe \all{ver-} → pas de \all{ge-})\\
3. a) \all{Ich weiß, dass das Smartphone teuer est.} (premier point : conjonction « que » → \all{dass} ; deuxième : article → \all{das}) b) \all{Das Internetcafé ist heute geschlossen.} (article → \all{das}.)
\end{corrige}

\courssub{Méthode : Exprimer son opinion sur les médias (écrit/oral)}

\begin{methode}[Rédiger un paragraphe d'opinion structuré]
Pour réussir une production écrite ou orale sur les avantages/inconvénients des médias (sujet type du Bac), suivez ces 4 étapes :
\begin{enumerate}
\item \textbf{Préparez le vocabulaire :} listez 3 avantages et 3 inconvénients en allemand (\all{der Vorteil}, \all{der Nachteil}, \all{praktisch}, \all{teuer}, \all{gefährlich}, \all{nützlich}, \all{abhängig machen}).
\item \textbf{Structurez votre paragraphe :} 
\begin{itemize}
\item Phrase d'introduction : \all{Meiner Meinung nach hat das Smartphone Vor- und Nachteile.} (À mon avis, le smartphone a des avantages et des inconvénients.)
\item Argument 1 (+ \all{dass}) : \all{Ein Vorteil ist, dass man immer erreichbar ist.}
\item Argument 2 (+ \all{dass}) : \all{Ein Nachteil ist, dass man zu viel Zeit verliert.}
\item Conclusion : \all{Ich finde, dass man das Smartphone bewusst nutzen sollte.}
\end{itemize}
\item \textbf{Vérifiez la grammaire :} 
\begin{itemize}
\item Virgule avant \all{dass} ? Verbe à la fin de la subordonnée ? (\all{...ist}, \all{...verliert}, \all{...sollte}).
\item Comparatifs corrects ? (\all{besser}, \all{schlechter}, \all{teurer}).
\item Parfait correct si vous racontez une expérience ? (\all{Ich habe ... geschrieben / bin ... gegangen}).
\end{itemize}
\item \textbf{Relisez à voix haute} pour vérifier la fluidité et la prononciation (Umlaute, \all{ß}, accent tonique).
\end{enumerate}
\end{methode}

\begin{experience}[Débat guidé : « Smartphone ou Cybercafé ? »]
En binôme. L'élève A défend le smartphone, l'élève B le cybercafé. Utilisez impérativement :
\begin{itemize}
\item 2 subordonnées en \all{dass} (\all{Ich denke, dass... / Ich finde, dass...}).
\item 1 phrase au parfait (\all{Gestern habe ich... / Ich bin... gegangen}).
\item 1 comparatif (\all{besser als}, \all{schneller als}, \all{günstiger als}).
\end{itemize}
Exemple de début : \all{A: Ich finde, dass das Smartphone besser ist, weil ich überall chatten kann. B: Das stimmt, aber ich denke, dass das Internetcafé günstiger ist. Gestern bin ich ins Internetcafé gegangen und habe nur 500 FCFA bezahlt.}
\end{experience}

\begin{aretenir}[L'essentiel de la section]
\begin{itemize}
\item Après \all{dass} (que) : virgule obligatoire + sujet + ... + \textbf{verbe conjugué à la fin} : \all{Ich denke, dass das Smartphone praktisch ist.}
\item \all{dass} = conjonction « que » ; \all{das} = article ou pronom (test : remplaçable par \all{dieses}).
\item Parfait des verbes forts : \all{haben}/\all{sein} + participe en \all{ge-...-en} : \all{geschrieben, gelesen, gefunden, genommen, gegangen, gekommen, verstanden, verloren}.
\item Auxiliaire \all{sein} pour les déplacements et changements d'état : \all{gehen, kommen, fahren, bleiben, aufstehen, passieren}.
\item Pas de \all{ge-} pour les verbes en \all{-ieren} et à préfixe inséparable (\all{ver-, be-, er-, ent-...}) : \all{studiert, verloren, verstanden, erzählt}.
\end{itemize}
\end{aretenir}

\coursec{Méthode du commentaire, commentaire modèle et production guidée}

Dans la section précédente, nous avons appris à construire des subordonnées en \all{dass} et à conjuguer les verbes forts au parfait : \all{Ich habe eine E-Mail geschrieben.} Ces deux outils grammaticaux ne sont pas des exercices isolés : ce sont précisément les phrases que l'on utilise pour \emph{commenter} un texte. Dire ce que l'auteur pense (\all{Der Autor meint, dass...}), raconter ce que l'on a fait (\all{Ich habe gestern gechattet...}) : c'est cela, le commentaire de texte. Dans cette dernière section, nous allons donc assembler toutes les pièces de la leçon — le vocabulaire des médias, la comparaison, les subordonnées en \all{dass}, le parfait des verbes forts — pour produire deux choses concrètes : un \cle{commentaire de texte} et une \cle{production personnelle guidée}.

\courssub{La méthode du commentaire en quatre étapes}

Un commentaire de texte (\all{die Textbesprechung}) n'est pas un résumé mot à mot. C'est une petite dissertation orale ou écrite qui suit toujours le même plan, en allemand comme en français. Retenez ces quatre étapes.

\begin{methode}[Rédiger un commentaire de texte en 4 étapes]
\begin{enumerate}
\item \textbf{Présenter le texte} : de quoi parle-t-on ? Commencez toujours par la formule \all{Der Text handelt von...} (Le texte parle de...). Ajoutez si possible le type de texte (\all{ein Artikel}, \all{ein Interview}, \all{ein Dialog}) et le thème.
\item \textbf{Dégager l'idée principale} : que veut dire l'auteur ? Utilisez \all{Der Autor meint, dass...} ou \all{Der Autor zeigt, dass...} — attention, le verbe conjugué va à la fin de la subordonnée.
\item \textbf{Analyser les arguments} : citez brièvement le texte et opposez les points de vue avec \all{Einerseits..., andererseits...} (d'une part..., d'autre part...). C'est ici que vous placez vos comparatifs : \all{Das Smartphone ist praktischer als das Internetcafé.}
\item \textbf{Donner son avis personnel argumenté} : concluez avec \all{Meiner Meinung nach...} (à mon avis) suivi d'une justification avec \all{weil} ou \all{denn}.
\end{enumerate}
\end{methode}

\begin{popfigure}
\begin{center}
\begin{tikzpicture}[
box/.style={draw=popblue, fill=popblueL, rounded corners=3pt, align=center, text width=4.2cm, minimum height=1.1cm, font=\small},
fleche/.style={-{Stealth[length=3mm]}, thick, poporange}]
\node[box, align=center] (a) at (0,0){\textbf{1. Présentation}\\ \all{Der Text handelt von...}};
\node[box, fill=popgreenL, draw=popgreen, align=center] (b) at (0,-1.9){\textbf{2. Idée principale}\\ \all{Der Autor meint, dass...}};
\node[box, fill=poppurpleL, draw=poppurple, align=center] (c) at (0,-3.8){\textbf{3. Analyse}\\ \all{Einerseits... andererseits...}};
\node[box, fill=popgoldL, draw=popgold, align=center] (d) at (0,-5.7){\textbf{4. Avis personnel}\\ \all{Meiner Meinung nach...}};
\draw[fleche] (a) -- (b);
\draw[fleche] (b) -- (c);
\draw[fleche] (c) -- (d);
\end{tikzpicture}
\end{center}
\legende{Organigramme du commentaire de texte : quatre étapes dans l'ordre, quatre phrases-types.}
\end{popfigure}

\courssub{Les phrases-types du commentaire}

Voici la boîte à outils du commentateur. Apprenez ces formules par cœur : elles fonctionnent avec n'importe quel texte, au Bac comme en classe.

\begin{center}
\begin{tabularx}{\linewidth}{|X|X|X|}
\hline
\rowcolor{popblueL} Deutsch & Français & Fonction \\
\hline
\all{Der Text handelt von Smartphones und Internetcafés.} & Le texte parle des smartphones et des cybercafés. & présenter \\
\hline
\all{Der Autor meint, dass beide Medien wichtig sind.} & L'auteur pense que les deux médias sont importants. & idée principale \\
\hline
\all{Einerseits ist das Smartphone praktischer.} & D'une part, le smartphone est plus pratique. & argument + \\
\hline
\all{Andererseits ist das Internetcafé günstiger.} & D'autre part, le cybercafé est plus économique. & argument -- \\
\hline
\all{Ich habe gestern eine E-Mail geschrieben.} & J'ai écrit un e-mail hier. & expérience personnelle \\
\hline
\all{Meiner Meinung nach ist das Smartphone nützlicher.} & À mon avis, le smartphone est plus utile. & avis personnel \\
\hline
\end{tabularx}
\end{center}

\begin{exemplebox}[Phrases du commentaire modèle, traduites]
\all{Der Autor zeigt, dass beide Medien Vorteile haben.} « L'auteur montre que les deux médias ont des avantages. » — Observez : le verbe \all{haben} est rejeté à la fin de la subordonnée en \all{dass}.\par
\all{Meiner Meinung nach ist das Smartphone nützlicher, weil man es überall benutzen kann.} « À mon avis, le smartphone est plus utile parce qu'on peut l'utiliser partout. » — Observez : le comparatif \all{nützlicher} et le verbe \all{kann} à la fin après \all{weil}.
\end{exemplebox}

\courssub{Commentaire modèle}

Lisons maintenant un commentaire complet, rédigé selon la méthode en quatre étapes, sur le texte support de la leçon.

\begin{textebox}[Commentaire modèle — lu et analysé en classe]
Der Text handelt von Smartphones und Internetcafés in Kamerun. Der Autor meint, dass beide Medien für die Jugendlichen wichtig sind. Einerseits ist das Smartphone praktischer, weil man überall chatten und telefonieren kann. Andererseits ist das Internetcafé günstiger für die Schularbeit, denn man kann dort drucken und lange im Internet surfen. Ich habe gestern eine E-Mail im Internetcafé geschrieben, und das war schnell und billig. Meiner Meinung nach ist das Smartphone das beste Medium für die Kommunikation, aber für die Schule gehe ich lieber ins Internetcafé.
\end{textebox}

\textbf{Traduction.} « Le texte parle des smartphones et des cybercafés au Cameroun. L'auteur pense que les deux médias sont importants pour les jeunes. D'une part, le smartphone est plus pratique, parce qu'on peut chatter et téléphoner partout. D'autre part, le cybercafé est plus économique pour le travail scolaire, car on peut y imprimer et surfer longtemps sur Internet. Hier, j'ai écrit un e-mail au cybercafé, et c'était rapide et bon marché. À mon avis, le smartphone est le meilleur média pour la communication, mais pour l'école, je préfère aller au cybercafé. »

\textbf{Repérage des outils de la leçon dans le modèle.} Cherchons ensemble les trois outils grammaticaux étudiés dans cette leçon :

\begin{itemize}
\item \textbf{Un comparatif} : \all{praktischer} (plus pratique), \all{günstiger} (plus économique) — adjectif + \all{-er}.
\item \textbf{Une subordonnée en \all{dass}} : \all{Der Autor meint, \textbf{dass} beide Medien für die Jugendlichen wichtig \textbf{sind}.} — le verbe conjugué \all{sind} est bien à la fin.
\item \textbf{Un verbe fort au parfait} : \all{Ich habe gestern eine E-Mail \textbf{geschrieben}.} — \all{schreiben -- schrieb -- hat geschrieben}, participe irrégulier avec \all{ge-} et changement de voyelle.
\end{itemize}

\begin{savaistu}[« Meiner Meinung nach » : la formule qui plaît au correcteur]
\all{Meiner Meinung nach} (à mon avis) est une tournure très appréciée au Bac pour exprimer le point de vue personnel : elle montre que l'élève sait conclure un commentaire. On peut aussi dire \all{Ich bin der Meinung, dass...} (Je suis d'avis que...) — avec, là encore, le verbe à la fin : \all{Ich bin der Meinung, dass das Smartphone zu teuer ist.} Attention à la place de \all{nach} : on dit \all{meiner Meinung \textbf{nach}}, jamais \all{nach meiner Meinung} dans l'usage scolaire courant.
\end{savaistu}

\courssub{Exercice résolu : compléter un commentaire}

\begin{exoresolu}[Complétez le commentaire avec les mots de la boîte]
Mots à placer : \all{handelt -- dass -- einerseits -- andererseits -- geschrieben -- Meiner Meinung nach}\par
Der Text \rule{2.5cm}{0.4pt} von den Medien in Kamerun. Der Autor zeigt, \rule{2.5cm}{0.4pt} viele Jugendliche ein Smartphone haben. \rule{2.5cm}{0.4pt} ist das Smartphone schneller, \rule{2.5cm}{0.4pt} ist das Internetcafé billiger. Ich habe gestern eine E-Mail \rule{2.5cm}{0.4pt}. \rule{2.5cm}{0.4pt} ist das Internetcafé besser für die Schule.
\tcblower
\textbf{Solution détaillée.}\par
1. Der Text \textbf{handelt} von den Medien in Kamerun. — Formule de présentation : \all{handeln von} (parler de), ici conjugué à la 3\textsuperscript{e} personne.\par
2. Der Autor zeigt, \textbf{dass} viele Jugendliche ein Smartphone haben. — Subordonnée : conjonction \all{dass}, verbe \all{haben} à la fin.\par
3. \textbf{Einerseits} ist das Smartphone schneller, \textbf{andererseits} ist das Internetcafé billiger. — Le couple de connecteurs d'opposition ; notez les deux comparatifs \all{schneller} et \all{billiger}.\par
4. Ich habe gestern eine E-Mail \textbf{geschrieben}. — Parfait du verbe fort \all{schreiben} : auxiliaire \all{habe} en position 2, participe \all{geschrieben} à la fin.\par
5. \textbf{Meiner Meinung nach} ist das Internetcafé besser für die Schule. — Formule d'avis personnel, suivie de l'inversion (verbe avant sujet) car elle occupe la première position.
\end{exoresolu}

\courssub{Production guidée : « Das Smartphone : Vorteile und Nachteile für mich »}

À votre tour ! Rédigez un court paragraphe de 6 à 8 lignes en suivant le \cle{plan imposé} ci-dessous. Ce plan reprend exactement la méthode du commentaire, appliquée à votre expérience personnelle.

\begin{methode}[Plan imposé de la production]
\begin{enumerate}
\item \textbf{Introduction} (1 phrase) : \all{Das Smartphone ist für mich sehr wichtig.} ou \all{Ich habe ein Smartphone.}
\item \textbf{Avantage 1} (1-2 phrases) : la communication — \all{Ich kann mit meinen Freunden chatten.}
\item \textbf{Avantage 2} (1-2 phrases) : l'information ou l'école — \all{Ich finde schnell Informationen für die Hausaufgaben.}
\item \textbf{Inconvénient} (1-2 phrases) : \all{Aber das Smartphone ist teuer, und die Verbindung ist manchmal schlecht.}
\item \textbf{Avis personnel} (1 phrase) : \all{Meiner Meinung nach ist das Smartphone nützlicher als das Internetcafé, weil...}
\end{enumerate}
\end{methode}

\begin{attention}[Une idée par phrase !]
Dans votre production, évitez les phrases trop longues : c'est l'erreur numéro un des francophones, qui traduisent mentalement de longues phrases françaises. En allemand, préférez \textbf{une idée par phrase}, reliée par un connecteur simple : \all{und} (et), \all{aber} (mais), \all{denn} (car), \all{weil} (parce que, verbe à la fin), \all{deshalb} (c'est pourquoi, inversion). Mauvais : une phrase de quatre lignes avec trois \all{und}. Bon : \all{Das Smartphone ist teuer. Deshalb benutze ich manchmal das Internetcafé.}
\end{attention}

\textbf{Grille d'auto-évaluation.} Avant de rendre votre paragraphe, vérifiez chaque point et cochez :

\begin{center}
\begin{tabularx}{\linewidth}{|X|l|}
\hline
\rowcolor{popgreenL} Critère & Oui / Non \\
\hline
J'ai utilisé au moins \textbf{3 mots du vocabulaire du thème} (\all{das Smartphone, das Internetcafé, der Vorteil, der Nachteil, chatten, die Verbindung...}). & \\
\hline
J'ai écrit \textbf{un comparatif} correct (adjectif + \all{-er} : \all{nützlicher, billiger, schneller}). & \\
\hline
J'ai construit \textbf{une phrase en \all{dass}} avec le verbe à la fin. & \\
\hline
J'ai employé \textbf{un verbe fort au parfait} (\all{habe geschrieben, habe gelesen, bin gegangen...}). & \\
\hline
J'ai respecté le plan imposé (introduction, 2 avantages, 1 inconvénient, avis). & \\
\hline
Mes noms allemands ont tous une \textbf{majuscule}. & \\
\hline
\end{tabularx}
\end{center}

\courssub{Expression orale guidée : présenter son paragraphe}

\begin{experience}[Présentation orale en une minute]
Travail par binômes puis devant la classe.\par
\textbf{Étape 1.} Relisez votre paragraphe à voix haute à votre voisin, lentement, en regardant votre feuille le moins possible. Votre voisin vérifie avec la grille d'auto-évaluation.\par
\textbf{Étape 2.} Présentez votre texte à la classe en une minute. Commencez par \all{Ich spreche über das Smartphone.} (Je parle du smartphone.)\par
\textbf{Étape 3.} Les camarades posent deux questions simples : \all{Was machst du mit dem Smartphone?} (Que fais-tu avec le smartphone ?), \all{Wie oft gehst du ins Internetcafé?} (Combien de fois vas-tu au cybercafé ?), \all{Hast du ein Smartphone?} Répondez par une phrase complète.\par
\textbf{Conseil de prononciation} : \all{Smartphone} se dit « smarte-fone », \all{chatten} se dit « tchatt-n », \all{Verbindung} se dit « fer-bin-doung ».
\end{experience}

\begin{exercice}[À vous de jouer]
Rédigez votre paragraphe « \all{Das Smartphone: Vorteile und Nachteile für mich} » (6 à 8 lignes) en suivant le plan imposé, puis auto-évaluez-le avec la grille. Préparez ensuite votre présentation orale d'une minute.
\end{exercice}

\begin{corrige}[Exercice — modèle de production attendue]
\begin{textebox}[Modèle de production]
Ich habe ein Smartphone, und es ist für mich sehr wichtig. Einerseits kann ich mit meinen Freunden und meiner Familie chatten. Andererseits finde ich schnell Informationen für die Schule. Das ist ein großer Vorteil. Aber das Smartphone hat auch einen Nachteil: es ist teuer, und die Verbindung ist in meinem Quartier manchmal schlecht. Ich habe gestern eine E-Mail an meinen Lehrer geschrieben, und das war praktisch. Meiner Meinung nach ist das Smartphone nützlicher als das Internetcafé, weil man es überall benutzen kann.
\end{textebox}

\textbf{Traduction.} « J'ai un smartphone, et il est très important pour moi. D'une part, je peux chatter avec mes amis et ma famille. D'autre part, je trouve vite des informations pour l'école. C'est un grand avantage. Mais le smartphone a aussi un inconvénient : il est cher, et la connexion est parfois mauvaise dans mon quartier. Hier, j'ai écrit un e-mail à mon professeur, et c'était pratique. À mon avis, le smartphone est plus utile que le cybercafé, parce qu'on peut l'utiliser partout. »

\textbf{Vérification avec la grille.}
\begin{itemize}
\item \textbf{Vocabulaire (5 mots)} : \all{Smartphone, chatten, Vorteil, Nachteil, Verbindung, Internetcafé, E-Mail, Lehrer, Schule}.
\item \textbf{Comparatif} : \all{nützlicher} (comparatif de supériorité).
\item \textbf{Subordonnée en \all{dass}} (alternative au \all{weil} utilisé) : \all{Ich finde, dass das Smartphone nützlich ist.}
\item \textbf{Verbe fort au parfait} : \all{habe geschrieben} (\all{schreiben -- schrieb -- geschrieben}).
\item \textbf{Plan respecté} : Introduction, 2 avantages (\all{Einerseits... Andererseits...}), 1 inconvénient (\all{Aber...}), Avis (\all{Meiner Meinung nach...}).
\item \textbf{Majuscules} : Tous les noms (\all{Smartphone, Vorteil, Nachteil, Verbindung, Quartier, E-Mail, Lehrer, Internetcafé, Meinung}) sont capitalisés.
\end{itemize}
\end{corrige}

\begin{aretenir}[L'essentiel de la section]
Un commentaire de texte suit quatre étapes : présenter (\all{Der Text handelt von...}), dégager l'idée principale (\all{Der Autor meint, dass...}), analyser (\all{Einerseits... andererseits...}), conclure (\all{Meiner Meinung nach...}). Chaque étape mobilise un outil de la leçon : le vocabulaire des médias, la subordonnée en \all{dass}, le comparatif, le parfait des verbes forts. Pour la production personnelle : plan imposé, une idée par phrase, un connecteur par phrase, et auto-évaluation avec la grille avant de rendre.
\end{aretenir}

\newpage

\coursec{Activité d'intégration}

\courssub{Objectifs de l'activité}

À la fin de cette activité d'intégration, l'élève sera capable de :

\begin{itemize}
\item mobiliser le \cle{vocabulaire} des médias et de la communication (das Smartphone, das Internetcafé, der Vorteil, der Nachteil, chatten, die Verbindung) dans un échange oral et un texte écrit ;
\item employer correctement le \cle{comparatif} et le \cle{superlatif} pour comparer deux moyens de communication ;
\item construire des \cle{subordonnées en dass} avec le verbe conjugué à la fin ;
\item utiliser le \cle{parfait des verbes forts} pour raconter des expériences personnelles ;
\item défendre un point de vue à l'oral dans un débat réglé, puis rédiger un court \cle{article} structuré (introduction, arguments, avis personnel) pour le journal du lycée.
\end{itemize}

\courssub{Situation}

Le journal de ton lycée à Yaoundé, \all{„Unsere Schülerzeitung“}, prépare un numéro spécial sur les médias. La rédaction lance un débat dans les classes de Première : vaut-il mieux, pour un élève camerounais, posséder un smartphone ou fréquenter le cybercafé du quartier ? Le meilleur article sera publié. Voici l'appel à contributions affiché dans la salle des fêtes :

\begin{textebox}[Appel à contributions — Aufruf der Schülerzeitung]
Liebe Schülerinnen und Schüler!

Fast alle Jugendlichen in Yaoundé sind heute online. Viele haben ein eigenes Smartphone, andere gehen ins Internetcafé am Markt. Aber was ist besser für uns Schüler?

Das Smartphone ist klein und praktisch: Man kann überall chatten, Fotos machen und Hausaufgaben recherchieren. Aber es ist teuer, und die Verbindung ist manchmal schlecht. Das Internetcafé ist billiger und der Computer hat einen großen Bildschirm. Aber man muss laufen, warten und bezahlen.

Schreibt einen kurzen Artikel (circa 10 Zeilen) für unsere Zeitung! Vergleicht das Smartphone und das Internetcafé: Nennt zwei Vorteile und einen Nachteil und sagt eure Meinung. Der beste Artikel kommt in die nächste Ausgabe!

Eure Redaktion
\end{textebox}

\begin{definition}[Glossaire du texte]
\begin{itemize}
\item \all{praktisch} : pratique ; \all{recherchieren} : faire des recherches
\item \all{die Verbindung, -en} : la connexion ; \all{billig} : bon marché
\item \all{der Bildschirm, -e} : l'écran ; \all{die Ausgabe, -n} : le numéro (d'un journal)
\item \all{die Meinung, -en} : l'avis, l'opinion ; \all{vergleichen (verglich, hat verglichen)} : comparer
\end{itemize}
\end{definition}

\courssub{Consignes}

\begin{enumerate}
\item \textbf{Lecture et repérage.} Lis l'appel à contributions à voix haute avec ton voisin. Souligne en rouge deux \cle{comparatifs} et en bleu deux mots du champ lexical des médias. Recopie-les sur ton cahier avec leur traduction.

\item \textbf{Préparation en équipe.} La classe est divisée en deux équipes : l'équipe A défend le smartphone, l'équipe B défend le cybercafé. Par groupes de quatre, rédigez \textbf{quatre arguments écrits}. Chaque argument doit respecter au moins une de ces contraintes : contenir un comparatif (ex. \all{Das Smartphone ist schneller als das Internetcafé.}), contenir une subordonnée en \all{dass} (ex. \all{Ich denke, dass das Internetcafé billiger ist.}), contenir un verbe fort au parfait (ex. \all{Ich bin gestern ins Internetcafé gegangen.}).

\item \textbf{Le débat.} Chaque équipe présente ses arguments à l'oral, à tour de rôle. L'équipe adverse doit répondre à chaque argument par une phrase commençant par \all{Aber ich glaube, dass ...} ou \all{Ich finde ... besser, weil ...}. Le professeur arbitre et note la correction de la langue (place du verbe, comparatif, parfait).

\item \textbf{Production écrite individuelle.} Rédige maintenant seul(e) ton article pour \all{„Unsere Schülerzeitung“} : environ 10 lignes, intitulé \all{„Smartphone oder Internetcafé?“}. Ton article doit contenir : une introduction (1--2 phrases), deux avantages de ta solution préférée, un inconvénient honnêtement reconnu, et un avis personnel argumenté. Utilise au minimum : deux comparatifs, deux subordonnées en \all{dass} et deux verbes forts au parfait.

\item \textbf{Relecture.} Échange ton article avec ton voisin. Vérifie : la majuscule à tous les noms, le verbe à la fin dans les phrases en \all{dass}, la forme du participe passé des verbes forts (\all{ge-...-en}), la place du verbe en position 2 dans les phrases principales.
\end{enumerate}

\courssub{Ressources à mobiliser}

\begin{propriete}[Rappel : le comparatif et le superlatif]
Comparatif : adjectif + \textbf{-er} + \all{als} : \all{Das Smartphone ist \textbf{kleiner als} der Computer.} « Le smartphone est plus petit que l'ordinateur. »

Superlatif : \all{am} + adjectif + \textbf{-sten} : \all{Das Internetcafé ist \textbf{am billigsten}.} « Le cybercafé est le moins cher. »

Irréguliers à connaître : \all{gut -- besser -- am besten} ; \all{gern -- lieber -- am liebsten} ; \all{viel -- mehr -- am meisten}. Umlaut sur les adjectifs courts : \all{alt -- älter}, \all{jung -- jünger}, \all{groß -- größer}.
\end{propriete}

\begin{propriete}[Rappel : la subordonnée en « dass » et le parfait des verbes forts]
Dans la subordonnée introduite par \all{dass}, le verbe conjugué va \textbf{à la fin} : \all{Ich glaube, dass das Smartphone \textbf{teurer ist}.} « Je crois que le smartphone est plus cher. »

Parfait des verbes forts : auxiliaire (\all{haben} ou \all{sein}) + participe en \textbf{ge-...-en} : \all{gehen -- ist gegangen}, \all{schreiben -- hat geschrieben}, \all{finden -- hat gefunden}, \all{nehmen -- hat genommen}. Auxiliaire \all{sein} pour les verbes de déplacement : \all{Ich \textbf{bin} ins Internetcafé \textbf{gegangen}.}
\end{propriete}

\begin{center}
\begin{tabularx}{\linewidth}{|X|X|X|}
\hline
\rowcolor{popblueL} Français & Deutsch & Remarque \\
\hline
le smartphone & das Smartphone, -s & mot anglais intégré, neutre \\
\hline
le cybercafé & das Internetcafé, -s & aussi : das Cybercafé \\
\hline
l'avantage & der Vorteil, -e & contraire : der Nachteil, -e \\
\hline
la connexion & die Verbindung, -en & verbe : verbinden \\
\hline
chatter & chatten & verbe faible : ich chatte \\
\hline
comparer & vergleichen & fort : verglich, hat verglichen \\
\hline
à mon avis & meiner Meinung nach & suivi du verbe en position 2 \\
\hline
\end{tabularx}
\end{center}

\begin{attention}[Pièges à éviter dans ton article]
\begin{itemize}
\item Ne dis pas \all{Ich denke, dass das Smartphone ist besser.} : le verbe doit aller à la fin : \all{..., dass das Smartphone besser \textbf{ist}.}
\item Le comparatif ne se construit pas avec \all{mehr} devant l'adjectif : pas \all{mehr schnell}, mais \all{\textbf{schneller}}.
\item Participe des verbes forts : pas \all{hat gegangen}, mais \all{\textbf{ist} gegangen} (mouvement $\rightarrow$ \all{sein}).
\item N'oublie pas la majuscule : \all{das Smartphone, der Vorteil, die Verbindung}.
\end{itemize}
\end{attention}

\courssub{Proposition de production et corrigé commenté}

\begin{exoresolu}[Article modèle pour „Unsere Schülerzeitung“]
Rédige un article d'environ 10 lignes intitulé \all{„Smartphone oder Internetcafé?“} avec introduction, deux avantages, un inconvénient et un avis personnel argumenté.
\tcblower
\textbf{Production modèle :}

\all{\textbf{Smartphone oder Internetcafé?}}

\all{Heute haben fast alle Jugendlichen in Yaoundé ein Handy, aber viele gehen auch noch ins Internetcafé. Was ist besser für uns Schüler?}

\all{Ich finde das Smartphone besser als das Internetcafé. Erstens ist es praktischer: Man kann überall chatten und schneller Informationen finden. Zweitens glaube ich, dass man mit dem Smartphone mehr lernt, weil man auch zu Hause recherchieren kann. Letzte Woche habe ich zum Beispiel einen Text für den Deutschkurs im Internet gefunden und sofort gelesen.}

\all{Aber das Smartphone hat auch einen Nachteil: Es ist teurer, und ich denke, dass die Verbindung manchmal schlecht ist. Früher bin ich oft ins Internetcafé gegangen, weil es billiger war.}

\all{Meiner Meinung nach ist das Smartphone die beste Lösung für Schüler, denn es ist kleiner, schneller und immer dabei.}

\bigskip
\textbf{Commentaire des choix de langue :}
\begin{itemize}
\item \textbf{Comparatifs} : \all{besser als} (irrégulier : \all{gut -- besser}), \all{praktischer}, \all{schneller}, \all{teurer}, \all{billiger}, \all{kleiner} ; superlatif : \all{die beste Lösung}.
\item \textbf{Subordonnées en dass} : \all{..., dass man mit dem Smartphone mehr lernt} et \all{..., dass die Verbindung manchmal schlecht ist} : le verbe conjugué est bien à la fin.
\item \textbf{Parfait des verbes forts} : \all{habe ich ... gefunden} (\all{finden -- hat gefunden}), \all{bin ich ... gegangen} (\all{gehen -- ist gegangen}, auxiliaire \all{sein} car mouvement).
\item \textbf{Structure} : introduction avec question, deux avantages introduits par \all{Erstens / Zweitens}, un inconvénient introduit par \all{Aber}, avis personnel avec \all{Meiner Meinung nach} (verbe en position 2 juste après).
\item \textbf{Vocabulaire du thème} : \all{das Handy, das Internetcafé, chatten, die Verbindung, recherchieren}.
\end{itemize}
\end{exoresolu}

\courssub{Bilan de l'activité}

Cette activité t'a permis de réinvestir l'ensemble des acquis de la leçon dans une tâche de communication authentique : comprendre un appel à contributions, débattre à l'oral en équipe et rédiger un article d'opinion structuré. Vérifie que tu sais désormais :

\begin{itemize}
\item citer au moins six mots du vocabulaire des médias avec leur article et leur pluriel ;
\item former et employer le comparatif (\all{besser als}, \all{billiger als}) et le superlatif (\all{am besten}) ;
\item construire une subordonnée en \all{dass} avec le verbe à la fin ;
\item conjuguer au parfait les verbes forts fréquents (\all{gehen, finden, schreiben, nehmen, lesen}) en choisissant le bon auxiliaire ;
\item organiser un court texte d'opinion : introduction, arguments, nuance, conclusion personnelle.
\end{itemize}

Ces compétences seront réutilisées dans les prochaines leçons du chapitre « Média et communication », notamment pour commenter d'autres textes de presse et pour l'épreuve de production écrite.

\newpage

\coursec{Méthodes et exercices résolus}

\coursec{ALA15 : Smartphones et cybercafés — Avantages et inconvénients des médias}

\courssub{Texte support et lecture globale}

\begin{textebox}[Article de blog : La vie numérique des jeunes à Yaoundé]
\all{Heute besitzen fast alle Jugendlichen in Yaoundé ein Smartphone. Sie benutzen es jeden Tag, um mit Freunden zu chatten, Nachrichten zu lesen oder Videos anzusehen. Die Geräte sind sehr praktisch, weil man damit überall online sein kann. Aber es gibt auch Probleme: Nicht jede Familie hat einen schnellen Internetanschluss zu Hause. Deshalb gehen viele Schüler nach dem Unterricht in ein Internetcafé. Dort ist die Verbindung oft besser und die Computer sind schneller. Ein Nachteil ist jedoch, dass man pro Stunde bezahlen muss. Außerdem sind die Daten in öffentlichen Netzen manchmal nicht sicher. Meiner Meinung nach ist das Smartphone ein tolles Gerät, aber man muss vorsichtig sein und nicht zu viel Zeit damit verbringen.}
\end{textebox}

\begin{exemplebox}[Lecture globale — Questions de repérage]
\begin{enumerate}
\item Quel est le thème du texte ? \all{Smartphones und Internetcafés in Yaoundé.}
\item Quel type de texte est-ce ? \all{Ein Blogartikel / Ein Meinungsartikel.} (Article de blog / Article d'opinion)
\item Quelle est l'idée principale en une phrase ? \all{Der Text beschreibt die Vorteile und Nachteile von Smartphones und Internetcafés für Jugendliche in Yaoundé.}
\end{enumerate}
\end{exemplebox}

\courssub{Compréhension détaillée : questions et exploitation}

\begin{methode}[Lire et comprendre un texte sur les médias (global puis détaillé)]
\begin{enumerate}
\item \textbf{Première lecture globale} : lis le texte une fois sans dictionnaire. Repère le thème (\all{Smartphones und Internetcafés}), le type de texte (article de blog) et l'idée générale : \all{Der Text handelt von den Vor- und Nachteilen digitaler Medien für Jugendliche in Kamerun.}
\item \textbf{Repérage des mots-clés} : souligne le vocabulaire du thème : \all{das Smartphone}, \all{das Internetcafé}, \all{der Vorteil}, \all{der Nachteil}, \all{chatten}, \all{die Verbindung}, \all{der Internetanschluss}, \all{sicher}. Devine le sens par le contexte et les familles de mots (\all{verbinden} $\rightarrow$ \all{die Verbindung}, \all{sicher} $\rightarrow$ \all{die Sicherheit}).
\item \textbf{Lecture détaillée} : relis phrase par phrase. Pour chaque question, localise la ligne exacte de la réponse.
\item \textbf{Réponse rédigée} : réponds par une phrase complète en allemand, en reprenant les mots de la question. Respecte la place du verbe (2\up{e} position en principale, fin en subordonnée).
\item \textbf{Vérification} : majuscules aux noms (\all{Smartphone}, \all{Vorteil}), accords sujet-verbe, orthographe \all{dass} (double \all{s}).
\end{enumerate}
\end{methode}

\begin{exoresolu}[Questions de compréhension détaillée]
\textbf{Texte} : (voir texte support ci-dessus)

\textbf{Énoncé} : Beantworten Sie die Fragen auf Deutsch mit ganzen Sätzen.
\begin{enumerate}
\item Was besitzen fast alle Jugendlichen in Yaoundé heute?
\item Wofür benutzen sie das Smartphone? (Nennen Sie drei Aktivitäten.)
\item Warum gehen viele Schüler in ein Internetcafé?
\item Nennen Sie zwei Nachteile des Internetcafés, die im Text stehen.
\item Was rät der Autor am Ende des Textes?
\end{enumerate}
\tcblower
\textbf{Raisonnement} :
\begin{enumerate}
\item Réponse ligne 1 : \all{fast alle Jugendlichen... besitzen ein Smartphone}.
\item Réponse ligne 2 : \all{chatten, Nachrichten lesen, Videos ansehen}.
\item Réponse ligne 4-5 : \all{weil... kein schneller Internetanschluss zu Hause} (subordonnée \all{weil} $\rightarrow$ verbe à la fin).
\item Réponse ligne 5-6 : \all{bezahlen pro Stunde}, \all{Daten nicht sicher}. Structure : \all{Ein Nachteil ist, dass...} (subordonnée \all{dass} $\rightarrow$ verbe à la fin).
\item Réponse dernière phrase : \all{man muss vorsichtig sein und nicht zu viel Zeit damit verbringen}.
\end{enumerate}
\textbf{Réponses rédigées} :
\begin{enumerate}
\item \all{Fast alle Jugendlichen in Yaoundé besitzen heute ein Smartphone.}
\item \all{Sie benutzen es, um mit Freunden zu chatten, Nachrichten zu lesen und Videos anzusehen.}
\item \all{Viele Schüler gehen in ein Internetcafé, weil sie zu Hause keinen schnellen Internetanschluss haben.}
\item \all{Ein Nachteil ist, dass man pro Stunde bezahlen muss. Ein weiterer Nachteil ist, dass die Daten in öffentlichen Netzen manchmal nicht sicher sind.}
\item \all{Der Autor rät, vorsichtig zu sein und nicht zu viel Zeit mit dem Smartphone zu verbringen.}
\end{enumerate}
\end{exoresolu}

\courssub{Lexique thématique : les médias et la communication}

\begin{propriete}[Vocabulaire de base — Noms avec article et pluriel]
\begin{center}
\begin{tabularx}{\linewidth}{|X|X|X|}
\hline
\rowcolor{popblueL} \textbf{Allemand (Singulier)} & \textbf{Pluriel} & \textbf{Français} \\
\hline
\all{das Smartphone} & \all{die Smartphones} & le smartphone \\
\hline
\all{das Internetcafé} & \all{die Internetcafés} & le cybercafé \\
\hline
\all{der Vorteil} & \all{die Vorteile} & l'avantage \\
\hline
\all{der Nachteil} & \all{die Nachteile} & l'inconvénient \\
\hline
\all{die Verbindung} & \all{die Verbindungen} & la connexion, le lien \\
\hline
\all{der Internetanschluss} & \all{die Internetanschlüsse} & la connexion Internet (abonn.) \\
\hline
\all{die Nachricht} & \all{die Nachrichten} & le message, la nouvelle \\
\hline
\all{das Gerät} & \all{die Geräte} & l'appareil \\
\hline
\all{das Netz} & \all{die Netze} & le réseau (Internet) \\
\hline
\all{die Zeit} & \all{die Zeiten} & le temps \\
\hline
\end{tabularx}
\end{center}
\end{propriete}

\begin{propriete}[Verbes utiles et leurs constructions]
\begin{center}
\begin{tabularx}{\linewidth}{|X|X|X|}
\hline
\rowcolor{popblueL} \textbf{Verbe} & \textbf{Construction / Cas} & \textbf{Exemple} \\
\hline
\all{chatten} & \all{mit + Dativ} & \all{Ich chatte mit meinem Freund.} \\
\hline
\all{telefonieren} & \all{mit + Dativ} & \all{Sie telefoniert mit ihrer Mutter.} \\
\hline
\all{simsen / eine SMS schreiben} & \all{+ Akkusativ} & \all{Er schreibt eine SMS.} \\
\hline
\all{surfen} & \all{im Internet (in + Dem)} & \all{Wir surfen im Internet.} \\
\hline
\all{herunterladen / downloaden} & \all{+ Akkusativ} & \all{Ich lade ein Video herunter.} \\
\hline
\all{verbringen (Zeit)} & \all{+ Akkusativ (Zeit)} & \all{Er verbringt drei Stunden am Handy.} \\
\hline
\all{besitzen / haben} & \all{+ Akkusativ} & \all{Sie besitzt ein neues Handy.} \\
\hline
\end{tabularx}
\end{center}
\end{propriete}

\begin{aretenir}[Familles de mots et antonymes]
\begin{itemize}
\item \all{verbinden} (v.) $\rightarrow$ \all{die Verbindung} (n.) — \all{verbunden} (adj.) « connecté ».
\item \all{sicher} (adj.) « sûr » $\leftrightarrow$ \all{unsicher} « insécurisé » ; \all{die Sicherheit} (n.f.) « la sécurité ».
\item \all{der Vorteil} $\leftrightarrow$ \all{der Nachteil} (antonymes).
\item \all{nützlich} « utile » $\leftrightarrow$ \all{nutzlos} « inutile ».
\item \all{teuer} « cher » $\leftrightarrow$ \all{billig / günstig} « bon marché ».
\item \all{praktisch} « pratique » ; \all{gefährlich} « dangereux ».
\end{itemize}
\end{aretenir}

\begin{exoresolu}[Lexique : texte à trous et traduction]
\textbf{Énoncé} : Ergänzen Sie die Lücken mit les mots du tableau (forme correcte). Traduisez ensuite le texte.
\all{Aminas Bruder hat ein neues (1) \dots\ Er (2) \dots\ jeden Abend mit seinen Freunden. Zu Hause ist die (3) \dots\ sehr langsam. Deshalb geht er oft ins (4) \dots\ in der Stadt. Ein (5) \dots\ ist, dass es dort teuer ist. Die (6) \dots\ sind dort manchmal auch nicht (7) \dots.}
\tcblower
\textbf{Solution} :
\begin{enumerate}
\item \all{Smartphone} (neutre, indéfini \all{ein} $\rightarrow$ \all{ein neues Smartphone})
\item \all{chattet} (3\up{e} pers. sg. présent : \all{er chattet})
\item \all{Verbindung} (fém., déf. \all{die} $\rightarrow$ \all{die Verbindung})
\item \all{Internetcafé} (mouvement vers lieu \all{in das} = \all{ins} + accusatif neutre \all{das Internetcafé})
\item \all{Nachteil} (masc., indéfini \all{ein} $\rightarrow$ \all{ein Nachteil})
\item \all{Daten} (pluriel de \all{das Datum} / \all{die Daten} — ici « les données »)
\item \all{sicher} (adjectif attribut après \all{sind} $\rightarrow$ pas de déclinaison)
\end{enumerate}
\textbf{Texte complet} : \all{Aminas Bruder hat ein neues Smartphone. Er chattet jeden Abend mit seinen Freunden. Zu Hause ist die Verbindung sehr langsam. Deshalb geht er oft ins Internetcafé in der Stadt. Ein Nachteil ist, dass es dort teuer ist. Die Daten sind dort manchmal auch nicht sicher.}
\textbf{Traduction} : « Le frère d'Amina a un nouveau smartphone. Il chatte chaque soir avec ses amis. À la maison, la connexion est très lente. C'est pourquoi il va souvent au cybercafé en ville. Un inconvénient est qu'il faut payer à l'heure. Les données y sont parfois aussi pas sûres. »
\end{exoresolu}

\courssub{Grammaire 1 : La comparaison (Komparativ und Superlativ)}

\begin{propriete}[Formation du comparatif et du superlatif]
\begin{enumerate}
\item \textbf{Positif} : \all{schnell, teuer, praktisch, groß, gut}.
\item \textbf{Comparatif} : adjectif + \textbf{\all{-er}}.
      \begin{itemize}
      \item Règle générale : \all{schnell $\rightarrow$ schneller}, \all{teuer $\rightarrow$ teurer}, \all{praktisch $\rightarrow$ praktischer}.
      \item \textbf{Umlaut} pour la plupart des monosyllabes en \all{a, o, u} : \all{alt $\rightarrow$ älter}, \all{groß $\rightarrow$ größer} (\all{ß} $\rightarrow$ \all{ss}), \all{kurz $\rightarrow$ kürzer}, \all{lang $\rightarrow$ länger}, \all{stark $\rightarrow$ stärker}, \all{warm $\rightarrow$ wärmer}, \all{klug $\rightarrow$ klüger}, \all{jung $\rightarrow$ jünger}.
      \item \textbf{Irréguliers importants} : \all{gut $\rightarrow$ besser}, \all{viel $\rightarrow$ mehr}, \all{gerne $\rightarrow$ lieber}, \all{hoch $\rightarrow$ höher}, \all{nah $\rightarrow$ näher}.
      \end{itemize}
\item \textbf{Superlatif} : \textbf{\all{am}} + adjectif + \textbf{\all{-sten}} (prédicatif) / article défini + adjectif + \textbf{\all{-ste}} + terminaison (attribut).
      \begin{itemize}
      \item \all{schnell $\rightarrow$ am schnellsten} / \all{der schnellste}.
      \item \all{groß $\rightarrow$ am größten} / \all{der größte}.
      \item \all{gut $\rightarrow$ am besten} / \all{der beste}.
      \end{itemize}
\end{enumerate}
\end{propriete}

\begin{exemplebox}[Comparatif d'inégalité : \all{als} (« que »)]
\all{Das Smartphone ist \textbf{schneller als} der alte Computer.} \\
\all{Ein Internetcafé ist oft \textbf{billiger als} ein eigenes Gerät.} \\
\all{Die Verbindung zu Hause ist \textbf{langsamer als} im Café.}
\emph{Attention : jamais \all{wie} après un comparatif en \all{-er} !}
\end{exemplebox}

\begin{exemplebox}[Comparatif d'égalité : \all{so ... wie} (« aussi ... que »)]
\all{Das neue Handy ist \textbf{so teuer wie} das alte.} \\
\all{Chatten ist \textbf{so beliebt wie} telefonieren.} \\
\emph{Négation (infériorité) :} \all{nicht so ... wie} — \all{Das Internetcafé ist \textbf{nicht so schnell wie} mein WLAN zu Hause.}
\end{exemplebox}

\begin{exemplebox}[Superlatif : attribut vs prédicatif]
\begin{itemize}
\item \textbf{Prédicatif (après verbe d'état)} : \all{am + adj. + -sten} — \all{WhatsApp ist \textbf{am beliebtesten}.}
\item \textbf{Attribut (devant nom)} : article + adj. + \all{-ste} + terminaison — \all{Das ist \textbf{die beliebteste App}.} / \all{Das \textbf{beste Smartphone} ist teuer.}
\end{itemize}
\end{exemplebox}

\begin{exoresolu}[Comparaison : exercices variés]
\textbf{Énoncé} : a) Formez le comparatif et le superlatif (prédicatif) : \all{schnell, billig, groß, gut, praktisch, beliebt}. \\
b) Complétez : \all{Das Smartphone ist \dots\ (praktisch) als der Brief.} — \all{Das Internetcafé ist \dots\ (billig) als das Smartphone.} — \all{WhatsApp ist \dots\ (beliebt) App in meiner Klasse.} \\
c) Traduisez : « Le smartphone est plus cher que le téléphone portable classique, mais il est plus utile. » — « C'est l'application la plus populaire de ma classe. »
\tcblower
\textbf{Solution} :
\begin{enumerate}
\item[a)] \all{schnell – schneller – am schnellsten} ; \all{billig – billiger – am billigsten} ; \all{groß – größer – am größten} ; \all{gut – besser – am besten} ; \all{praktisch – praktischer – am praktischsten} ; \all{beliebt – beliebter – am beliebtesten}.
\item[b)] \all{praktischer als} ; \all{billiger als} ; \all{die beliebteste} (attribut fém. sg. : \all{die} + \all{beliebt} + \all{ste} + \all{e}).
\item[c)] \all{Das Smartphone ist teurer als das klassische Handy, aber es ist nützlicher.} (\all{teuer $\rightarrow$ teurer}, \all{nützlich $\rightarrow$ nützlicher}, \all{als} pour « que » comparatif). — \all{Das ist die beliebteste App in meiner Klasse.}
\end{enumerate}
\end{exoresolu}

\courssub{Grammaire 2 : Les subordonnées en « dass »}

\begin{propriete}[La subordonnée complétive en \all{dass}]
\begin{enumerate}
\item \textbf{Fonction} : Complément d'objet après verbes de parole, pensée, sentiment : \all{sagen, denken, glauben, meinen, hoffen, wissen, verstehen, finden}. Aussi après noms : \all{Ein Vorteil/Nachteil ist, dass...}
\item \textbf{Structure} : \all{Hauptsatz, dass + Subjekt + ... + VerbeConjugué.}
\item \textbf{Règle d'or} : Le verbe conjugué est \textbf{rejeté à la fin} de la subordonnée.
\item \textbf{Virgule} : Obligatoire avant \all{dass}.
\item \textbf{Orthographe} : \all{dass} = conjonction (double \all{s}) ; \all{das} = article (neutre) ou pronom démonstratif/relatif.
\end{enumerate}
\end{propriete}

\begin{exemplebox}[Exemples avec verbes modaux et verbes à particule]
\begin{itemize}
\item Verbe modal : \all{Sie sagt, dass sie jeden Tag chatten \textbf{will}.} (Modal \all{will} à la fin, infinitif \all{chatten} avant).
\item Verbe à particule séparable : \all{Er denkt, dass er das Video \textbf{ansieht}.} (Particule \all{an} collée au verbe \all{sieht} à la fin).
\item Perfekt dans \all{dass} : \all{Wir wissen, dass er gestern ins Internetcafé \textbf{gegangen ist}.} (Auxiliaire \all{ist} à la fin, participe \all{gegangen} avant).
\end{itemize}
\end{exemplebox}

\begin{exoresolu}[Subordonnées en \all{dass} : transformation et thème]
\textbf{Énoncé} : a) Transformez le discours direct en subordonnée \all{dass} : \all{Der Lehrer sagt: „Die Schüler chatten zu viel.“} — \all{Meine Mutter fragt: „Hast du ein Smartphone?“} (indirecte : \all{ob}). \\
b) Traduisez en allemand : « Je crois que les cybercafés sont importants pour les élèves. » — « Elle dit qu'elle n'a pas de smartphone. » — « Nous espérons que la connexion sera meilleure demain. »
\tcblower
\textbf{Solution} :
\begin{enumerate}
\item[a)] \all{Der Lehrer sagt, dass die Schüler zu viel chatten.} (Virgule, \all{dass}, verbe \all{chatten} à la fin). — \all{Meine Mutter fragt, ob ich ein Smartphone habe.} (\all{ob} pour question fermée, verbe à la fin).
\item[b)] \all{Ich glaube, dass die Internetcafés für die Schüler wichtig sind.} (\all{glauben + dass}, sujet pluriel \all{Internetcafés} $\rightarrow$ \all{sind} à la fin ; \all{wichtig für + Akk}). — \all{Sie sagt, dass sie kein Smartphone hat.} (\all{kein} pour nier nom indéfini \all{ein Smartphone}, verbe \all{hat} à la fin). — \all{Wir hoffen, dass die Verbindung morgen besser sein wird.} (\all{werden} + infinitif \all{sein} à la fin = Futur I).
\end{enumerate}
\end{exoresolu}

\courssub{Conjugaison : Le Parfait des verbes forts (Perfekt)}

\begin{propriete}[Structure du Perfekt]
\all{Subjekt + Auxiliaire (haben/sein) conjugué (Pos. 2) + ... + Participe Passé (à la fin).}
\end{propriete}

\begin{propriete}[Choix de l'auxiliaire]
\begin{itemize}
\item \textbf{\all{haben}} : Verbes transitifs (avec objet accusatif) + la majorité des verbes intransitifs. \all{Ich habe ein Smartphone gekauft.}
\item \textbf{\all{sein}} : Verbes de \textbf{mouvement} vers un lieu (\all{gehen, fahren, fliegen, laufen, rennen, schwimmen, kommen, ankommen, einreisen}) + verbes de \textbf{changement d'état} (\all{sterben, einschlafen, aufwachen, werden, bleiben, sein}). \all{Ich bin ins Internetcafé gegangen.}
\end{itemize}
\end{propriete}

\begin{propriete}[Formation du Participe Passé des verbes FORTS]
\textbf{\all{ge-} + Radical modifié + \all{-en}} (pas de \all{-t} final !)
\begin{center}
\begin{tabular}{|l|l|l|l|}
\hline
\rowcolor{popblueL} \textbf{Infinitif} & \textbf{Präteritum (3ème ps sg)} & \textbf{Participe Passé} & \textbf{Auxiliaire} \\
\hline
\all{gehen} & \all{ging} & \all{gegangen} & \all{sein} \\
\hline
\all{kommen} & \all{kam} & \all{gekommen} & \all{sein} \\
\hline
\all{finden} & \all{fand} & \all{gefunden} & \all{haben} \\
\hline
\all{geben} & \all{gab} & \all{gegeben} & \all{haben} \\
\hline
\all{nehmen} & \all{nahm} & \all{genommen} & \all{haben} \\
\hline
\all{sprechen} & \all{sprach} & \all{gesprochen} & \all{haben} \\
\hline
\all{lesen} & \all{las} & \all{gelesen} & \all{haben} \\
\hline
\all{sehen} & \all{sah} & \all{gesehen} & \all{haben} \\
\hline
\all{schreiben} & \all{schrieb} & \all{geschrieben} & \all{haben} \\
\hline
\all{treffen} & \all{traf} & \all{getroffen} & \all{haben} \\
\hline
\all{verstehen} & \all{verstand} & \all{verstanden} & \all{haben} \\
\hline
\all{empfehlen} & \all{empfahl} & \all{empfohlen} & \all{haben} \\
\hline
\end{tabular}
\end{center}
\emph{Note : Verbes à préfixe inséparable (\all{ver-, be-, er-, ent-, ge-, zer-, miss-}) : pas de \all{ge-} $\rightarrow$ \all{verstehen $\rightarrow$ verstanden}. Verbes en \all{-ieren} : pas de \all{ge-}, finissent en \all{-t} $\rightarrow$ \all{telefonieren $\rightarrow$ telefoniert} (verbe faible).}
\end{propriete}

\begin{exoresolu}[Perfekt : mise au passé et thème]
\textbf{Énoncé} : Mettez au Perfekt : 1. \all{Ich gehe ins Internetcafé.} 2. \all{Sie schreibt eine E-Mail.} 3. \all{Wir lesen die Nachrichten.} 4. \all{Er findet ein günstiges Handy.} \\
Traduisez : « Hier, mon frère a trouvé un cybercafé près de l'école. » — « Nous avons passé trois heures en ligne. »
\tcblower
\textbf{Solution} :
\begin{enumerate}
\item \all{Ich \textbf{bin} ins Internetcafé \textbf{gegangen}.} (\all{gehen} mouvement $\rightarrow$ \all{sein} + \all{gegangen}).
\item \all{Sie \textbf{hat} eine E-Mail \textbf{geschrieben}.} (\all{schreiben} transitif $\rightarrow$ \all{haben} + \all{geschrieben}).
\item \all{Wir \textbf{haben} die Nachrichten \textbf{gelesen}.} (\all{lesen} transitif $\rightarrow$ \all{haben} + \all{gelesen}).
\item \all{Er \textbf{hat} ein günstiges Handy \textbf{gefunden}.} (\all{finden} transitif $\rightarrow$ \all{haben} + \all{gefunden}).
\end{enumerate}
\textbf{Traduction} :
\begin{itemize}
\item \all{Gestern \textbf{hat} mein Bruder ein Internetcafé in der Nähe der Schule \textbf{gefunden}.} (\all{gestern} = « hier » temporel ; \all{hier} = « ici » spatial — \textbf{faux ami} ! ; \all{in der Nähe + Genitiv}).
\item \all{Wir \textbf{haben} drei Stunden online \textbf{verbracht}.} (\all{verbringen} faible $\rightarrow$ \all{verbracht}, aux. \all{haben}).
\end{itemize}
\end{exoresolu}

\courssub{Méthode du commentaire et production guidée}

\begin{methode}[Rédiger un paragraphe d'opinion structuré (Vor- und Nachteile)]
\begin{enumerate}
\item \textbf{Analyse du sujet} : Identifie le thème (Smartphone / Internetcafé) et la tâche (avantages/inconvénients + avis).
\item \textbf{Plan en 4 étapes} :
      \begin{enumerate}
      \item \textbf{Introduction} (1 phrase) : \all{Heute ist das Smartphone für Jugendliche unverzichtbar.}
      \item \textbf{Avantages} (2-3 phrases) : \all{Ein Vorteil ist, dass... Außerdem...} + \textbf{un comparatif}.
      \item \textbf{Inconvénients} (2-3 phrases) : \all{Ein Nachteil ist, dass... Auch...} + \textbf{une phrase au Perfekt} (exemple concret).
      \item \textbf{Conclusion personnelle} (1 phrase) : \all{Meiner Meinung nach... / Ich finde..., aber...}
      \end{enumerate}
\item \textbf{Connecteurs logiques} : \all{zuerst, außerdem, außerdem, aber, deshalb, weil, dass, meiner Meinung nach, zum Schluss}.
\item \textbf{Contraintes grammaticales obligatoires (Bac)} : Au moins 1 subordonnée \all{dass/weil} (verbe fin), 1 comparatif \all{...er als}, 1 superlatif \all{am ...sten}, 1 Perfekt (fort), majuscules aux noms, \all{dass} (ss).
\item \textbf{Relecture} : Vérifie l'ordre des mots (V2 / verbe fin), les cas (accusatif/datif après prépositions), l'orthographe \all{ß/ss}.
\end{enumerate}
\end{methode}

\begin{exoresolu}[Rédaction guidée : « Das Smartphone : Vorteile und Nachteile »]
\textbf{Consigne} : Écrivez 7-9 phrases. Utilisez : \all{ein Vorteil ist, dass...}, \all{außerdem}, \all{ein Nachteil ist, dass...}, 1 comparatif, 1 Perfekt (verbe fort), 1 superlatif, \all{meiner Meinung nach}.
\tcblower
\textbf{Texte modèle} :
\all{Heute ist das Smartphone für fast alle Jugendlichen in Yaoundé das \textbf{wichtigste} Gerät. \textbf{Ein Vorteil ist,} dass man damit \textbf{schneller als} frueher Informationen \textbf{finden kann}. \textbf{Außerdem} kann man \textbf{überall} mit Freunden chatten und Fotos teilen. \textbf{Ein Nachteil ist,} dass viele Smartphones \textbf{sehr teuer} sind. \textbf{Ein weiterer Nachteil ist,} dass manche Jugendliche \textbf{zu viel Zeit} damit \textbf{verbringen}. \textbf{Gestern hat} mein Freund \textbf{vier Stunden} nur \textbf{gespielt} und \textbf{gechattet}! \textbf{Meiner Meinung nach} ist das Smartphone \textbf{sehr nützlich}, \textbf{aber} man muss es \textbf{mit Maß} benutzen.}

\textbf{Analyse des contraintes respectées} :
\begin{itemize}
\item Superlatif attribut : \all{das wichtigste Gerät}.
\item Subordonnée \all{dass} (verbe fin \all{kann}) : \all{dass man damit schneller als früher Informationen finden kann}.
\item Comparatif : \all{schneller als}.
\item Connecteur : \all{Außerdem}.
\item Subordonnée \all{dass} (verbe fin \all{sind}) : \all{dass viele Smartphones sehr teuer sind}.
\item Perfekt verbe fort (\all{verbringen} $\rightarrow$ \all{verbracht}) + exemple concret : \all{Gestern hat mein Freund vier Stunden nur gespielt und gechattet!} (\all{spielen} faible, \all{chatten} faible — \textit{note: pour verbe fort, ex: \all{gefunden} ou \all{geschrieben}}). \textbf{Correction pour contrainte verbe fort} : \all{Gestern hat mein Freund ein Video \textbf{heruntergeladen} und \textbf{geschrieben}.} (ou \all{gefunden}).
\item Opinion + concession : \all{Meiner Meinung nach... aber...}.
\end{itemize}
\end{exoresolu}

\courssub{Exercices d'entraînement (Devoirs / Classe)}

\begin{exercice}[Grammaire : Comparaison et Subordonnées]
\begin{enumerate}
\item Mettez l'adjectif au comparatif ou superlatif (prédicatif) : \all{Das Internet ist (schnell) \dots\ als das Buch.} — \all{Das Smartphone ist (teuer) \dots\ Gerät.} — \all{Die Verbindung im Café ist (gut) \dots\ als zu Hause.}
\item Reliez par \all{dass} (attention à la place du verbe) : \all{Ich denke. Das Handy ist praktisch.} — \all{Er sagt. Er hat kein Geld.} — \all{Wir hoffen. Die Verbindung wird besser.}
\item Traduisez : « Je pense que le cybercafé est moins cher qu'un smartphone. » — « Elle dit qu'elle est allée au cybercafé hier. » (Perfekt \all{gehen}).
\end{enumerate}
\end{exercice}

\begin{exercice}[Conjugaison : Perfekt des verbes forts]
Conjuguez au Perfekt (1\up{e} personne du singulier) : \all{finden, geben, nehmen, sprechen, treffen, lesen, sehen, schreiben, verstehen, empfehlen}. Indiquez l'auxiliaire.
\end{exercice}

\begin{exercice}[Expression écrite]
Rédigez un court paragraphe (6-8 phrases) sur : \all{„Das Internetcafé – Vor- und Nachteile“}. Utilisez le vocabulaire de la leçon, au moins une comparaison, une subordonnée \all{dass}, une phrase au Perfekt (verbe fort) et donnez votre avis.
\end{exercice}

\begin{corrige}[Exercice Grammaire]
\begin{enumerate}
\item \all{schneller als} — \all{das teuerste} (superlatif attribut neutre : \all{das teuerste Gerät}) ou \all{am teuersten} (prédicatif) — \all{besser als} (irrégulier \all{gut}).
\item \all{Ich denke, dass das Handy praktisch ist.} — \all{Er sagt, dass er kein Geld hat.} — \all{Wir hoffen, dass die Verbindung besser wird.}
\item \all{Ich denke, dass das Internetcafé billiger ist als ein Smartphone.} (ou \all{nicht so teuer wie}). — \all{Sie sagt, dass sie gestern ins Internetcafé gegangen ist.}
\end{enumerate}
\end{corrige}

\begin{corrige}[Exercice Conjugaison]
\all{ich habe gefunden} (haben) — \all{ich habe gegeben} (haben) — \all{ich habe genommen} (haben) — \all{ich habe gesprochen} (haben) — \all{ich habe getroffen} (haben) — \all{ich habe gelesen} (haben) — \all{ich habe gesehen} (haben) — \all{ich habe geschrieben} (haben) — \all{ich habe verstanden} (haben, pas de ge-) — \all{ich habe empfohlen} (haben, pas de ge-).
\end{corrige}

\begin{corrige}[Exercice Expression écrite — Proposition de correction]
\all{Ein Internetcafé hat Vor- und Nachteile. Ein Vorteil ist, dass die Computer dort oft schneller sind als alte Handys. Außerdem hat man eine gute Verbindung, wenn man zu Hause kein WLAN hat. Ein Nachteil ist, dass man pro Stunde bezahlen muss. Das ist teurer als ein eigener Vertrag. Ein weiterer Nachteil ist, dass die Daten nicht so sicher sind. Gestern habe ich im Internetcafé eine wichtige E-Mail geschrieben. Meiner Meinung nach ist das Internetcafé nützlich für Schüler ohne Internet zu Hause, aber man muss auf seine Daten aufpassen.}
\end{corrige}

\courssub{Attention : Les 6 pièges classiques du Bac (niveau Première)}

\begin{attention}[Erreurs éliminatoires à éviter absolument]
\begin{enumerate}
\item \textbf{\all{das} vs \all{dass}} : \all{Ich denke, \textbf{dass}...} (conjonction, double \all{s}) $\neq$ \all{\textbf{Das} ist mein Handy.} (article/pronom). \textbf{Astuce} : on peut remplacer \all{das} (article) par \all{ein} ou \all{dieses} ; \all{dass} ne se remplace pas.
\item \textbf{Verbe à la fin en subordonnée} : Après \all{dass, weil, ob, wenn, als} (temporel passé), le verbe conjugué va \textbf{à la toute fin}. \all{...weil die Verbindung \textbf{langsam ist}.} (Pas \all{...weil die Verbindung \textbf{ist} langsam.})
\item \textbf{\all{als} vs \all{wie} en comparaison} : Comparatif (\all{-er}) + \textbf{\all{als}} (« que »). Égalité (\all{so ...}) + \textbf{\all{wie}} (« que »). \textbf{Jamais} \all{schneller wie} ! \all{Als} aussi pour le passé ponctuel (\all{Als ich ankam...}).
\item \textbf{Majuscules aux noms} : \textbf{Tous} les noms prennent une majuscule : \all{das Smartphone, die Verbindung, der Vorteil, das Internetcafé, die Zeit}. Oublier la majuscule = faute d'orthographe grave.
\item \textbf{Auxiliaire \all{sein} au Perfekt} : Verbes de mouvement (\all{gehen, fahren, kommen, laufen, fliegen, schwimmen, reisen}) + changement d'état (\all{sterben, werden, bleiben, sein, aufwachen, einschlafen}) $\rightarrow$ \all{sein}. \all{Ich \textbf{bin} ins Café \textbf{gegangen}} (pas \all{habe gegangen}).
\item \textbf{Faux amis temporels} : \all{hier} = « \textbf{ici} » (lieu). « \textbf{Hier} » (temps) = \all{gestern}. \all{aktuell} = « actuel / en ce moment » (pas « actuel » au sens de « réel » $\rightarrow$ \all{tatsächlich}).
\end{enumerate}
\end{attention}

\begin{savaistu}[Contexte culturel : L'Allemagne et le numérique]
\begin{itemize}
\item \textbf{Allemagne} : Très attachée à la protection des données (\all{Datenschutz}, RGPD/\all{DSGVO}). Le paiement sans contact (\all{kontaktloses Bezahlen}) progresse mais l'argent liquide (\all{Bargeld}) reste roi. Le terme \all{Handy} (téléphone portable) est un pseudo-anglicisme très utilisé.
\item \textbf{Autriche / Suisse} : Vocabulaire proche, mais \all{das Natel} (Suisse) pour portable. Infrastructure Internet très dense.
\item \textbf{Cameroun} : Explosion du mobile money (\all{Mobile Money} / \all{Orange Money}, \all{MTN MoMo}) via smartphone. Les cybercafés (\all{Internetcafés} ou \all{Callboxes}) restent essentiels pour l'impression, le scan et la connexion haut débit pour les étudiants. Le bendskin (moto-taxi) utilise WhatsApp pour la commande.
\end{itemize}
\end{savaistu}

\begin{experience}[Activité orale : Débat « Smartphone à l'école ? »]
\textbf{Consigne} : En binôme. L'un défend le \all{Verbot} (interdiction), l'autre l' \all{Erlaubnis} (autorisation). Utilisez : \all{Ich bin dafür/dagegen, weil...}, \all{Ein Argument ist, dass...}, \all{Aber man muss bedenken, dass...}. Durée : 3 min. Le prof note : vocabulaire thème, \all{dass/weil} (verbe fin), comparatif, fluidité.
\end{experience}

\newpage

\coursec{Exercices}

\courssub{Je vérifie mes connaissances}

\begin{exercice}[QCM : vocabulaire des médias]
Entourez la bonne réponse.
\begin{enumerate}
\item \all{Das Internetcafé} signifie : a) le téléphone portable \quad b) le cybercafé \quad c) l'ordinateur portable
\item \all{Der Vorteil} signifie : a) l'inconvénient \quad b) le prix \quad c) l'avantage
\item \all{chatten} signifie : a) téléphoner \quad b) discuter en ligne \quad c) imprimer
\item \all{Die Verbindung ist schlecht.} signifie : a) La connexion est mauvaise. \quad b) La connexion est rapide. \quad c) La connexion est gratuite.
\item Quel est le pluriel de \all{das Smartphone} ? a) die Smartphones \quad b) die Smartphone \quad c) die Smartphonen
\item \all{Der Nachteil} a pour contraire : a) \all{der Preis} \quad b) \all{der Vorteil} \quad c) \all{das Internet}
\end{enumerate}
\end{exercice}

\begin{exercice}[Vrai ou faux ? Justifiez]
Lisez les phrases suivantes et dites si elles sont vraies (richtig) ou fausses (falsch) d'après le texte étudié. Justifiez chaque réponse par une courte phrase en allemand.
\begin{enumerate}
\item \all{Man kann mit dem Smartphone nur telefonieren.}
\item \all{Im Internetcafé kann man Dokumente drucken.}
\item \all{Das Smartphone ist immer billiger als das Internetcafé.}
\item \all{Viele Jugendliche in Kamerun benutzen das Internet zum Lernen.}
\end{enumerate}
\end{exercice}

\begin{exercice}[Articles et pluriels]
Complétez le tableau avec l'article et le pluriel de chaque nom.
\begin{center}
\begin{tabular}{|l|l|l|}
\hline
\rowcolor{popblueL} Nom & Article & Pluriel \\
\hline
Smartphone & \ldots & \ldots \\
\hline
Internetcafé & \ldots & \ldots \\
\hline
Vorteil & \ldots & \ldots \\
\hline
Nachteil & \ldots & \ldots \\
\hline
Verbindung & \ldots & \ldots \\
\hline
Computer & \ldots & \ldots \\
\hline
\end{tabular}
\end{center}
\end{exercice}

\begin{exercice}[Conjugaison : le parfait des verbes forts]
Conjuguez les verbes entre parenthèses au \all{Perfekt} (auxiliaire + participe passé).
\begin{enumerate}
\item Gestern \ldots\ ich eine Nachricht \ldots\ (schreiben).
\item Wir \ldots\ einen Text im Internet \ldots\ (lesen).
\item Paul \ldots\ gestern ins Internetcafé \ldots\ (gehen).
\item Marie \ldots\ viele Informationen \ldots\ (finden).
\item Ich \ldots\ mein Handy leider \ldots\ (verlieren).
\item Die Schüler \ldots\ den Bus in die Stadt \ldots\ (nehmen).
\end{enumerate}
\end{exercice}

\courssub{Je m'entraîne}

\begin{exercice}[Comparatif et superlatif]
Mettez les adjectifs suivants au comparatif et au superlatif, puis employez chaque forme dans une phrase complète sur le thème du smartphone ou du cybercafé.
\begin{center}
\begin{tabular}{|l|l|l|}
\hline
\rowcolor{popblueL} Positif & Komparativ & Superlativ \\
\hline
schnell & \ldots & \ldots \\
\hline
teuer & \ldots & \ldots \\
\hline
gut & \ldots & \ldots \\
\hline
groß & \ldots & \ldots \\
\hline
gern & \ldots & \ldots \\
\hline
praktisch & \ldots & \ldots \\
\hline
\end{tabular}
\end{center}
Exemple : \all{schnell} : \all{Das Smartphone ist schneller als der Brief. Das Internet ist am schnellsten.}
\end{exercice}

\begin{exercice}[Les subordonnées en dass]
Transformez les phrases suivantes : commencez par \all{Ich denke, dass \ldots} ou \all{Ich glaube, dass \ldots} Attention à la place du verbe !
\begin{enumerate}
\item Das Smartphone ist teuer.
\item Das Internetcafé ist billiger.
\item Viele Jugendliche chatten zu viel.
\item Die Verbindung ist heute schlecht.
\item Das Internet hilft den Schülern beim Lernen.
\item Mein Bruder hat sein Handy verloren.
\end{enumerate}
Exemple : \all{Das Wetter ist schön.} $\rightarrow$ \all{Ich denke, dass das Wetter schön ist.}
\end{exercice}

\begin{exercice}[Texte à trous]
Complétez le texte avec les mots suivants : \all{Verbindung, Vorteile, Nachteile, chattet, drucken, Internetcafé, Smartphone, teuer}.

\begin{textebox}[Texte à trous]
Mein Bruder hat ein neues \ldots\ . Er \ldots\ jeden Tag mit seinen Freunden. Aber das Handy ist sehr \ldots\ . Ich habe kein Handy. Ich gehe oft ins \ldots\ in der Stadt. Dort kann ich im Internet surfen und meine Hausaufgaben \ldots\ . Die \ldots\ ist dort schnell. Das Internet hat viele \ldots\ : Man lernt viel und findet Informationen. Aber es hat auch \ldots\ : Manche Jugendliche spielen zu viel und lernen nicht mehr.
\end{textebox}
\end{exercice}

\begin{exercice}[Appariement et traduction]
\textbf{A.} Reliez chaque mot allemand à sa traduction française.
\begin{center}
\begin{tabular}{|l|l|}
\hline
\rowcolor{popblueL} Deutsch & Français \\
\hline
1. die Verbindung & a) imprimer \\
\hline
2. drucken & b) la connexion \\
\hline
3. der Nachteil & c) bon marché \\
\hline
4. billig & d) l'inconvénient \\
\hline
5. die Nachricht & e) le message \\
\hline
\end{tabular}
\end{center}
\textbf{B.} Traduisez en allemand.
\begin{enumerate}
\item Je pense que le cybercafé est moins cher que le smartphone.
\item Le smartphone est plus pratique que l'ordinateur.
\item Hier, nous sommes allés au cybercafé et nous avons imprimé nos devoirs.
\item Ma sœur a écrit un message et a envoyé des photos.
\item Je crois qu'Internet a beaucoup d'avantages pour les élèves.
\end{enumerate}
\end{exercice}

\courssub{Je me prépare au Bac}

\begin{exercice}[Texte et questions de compréhension (type Bac)]
Lisez le texte suivant, puis répondez aux questions \textbf{en allemand, par des phrases complètes}, en justifiant vos réponses par le texte.

\begin{textebox}[Smartphones und Internetcafés in Kamerun]
In Kamerun benutzen heute viele Jugendliche ein Smartphone. Mit dem Handy können sie telefonieren, Nachrichten schreiben, Fotos machen und im Internet surfen. Das Smartphone ist praktisch, aber es ist auch teuer, und nicht alle Familien können es kaufen. Deshalb gehen viele Schüler ins Internetcafé. Dort ist der Preis niedrig: Für wenig Geld kann man eine Stunde im Internet surfen, E-Mails lesen und Dokumente drucken. Vor den Prüfungen sind die Internetcafés in Yaoundé und Douala oft voll. Die Schüler suchen Informationen für ihre Hausaufgaben und lernen mit ihren Freunden zusammen. Aber das Internet hat auch Nachteile. Manche Jugendliche chatten stundenlang und spielen nur. Sie lernen nicht mehr und haben schlechte Noten. Auch die Verbindung ist manchmal langsam, besonders am Abend. Die Lehrer sagen: „Das Internet ist ein gutes Werkzeug, aber man muss es richtig benutzen.“ Die Jugendlichen müssen also lernen, das Handy und das Internet klug zu nutzen: zum Lernen, zur Information und zur Kommunikation — nicht nur zum Spielen.
\end{textebox}

\textbf{Fragen zum Text :}
\begin{enumerate}
\item Was können die Jugendlichen mit dem Smartphone machen ? (Nennen Sie drei Dinge.)
\item Warum gehen viele Schüler ins Internetcafé ?
\item Was machen die Schüler im Internetcafé vor den Prüfungen ?
\item Welche Nachteile des Internets nennt der Text ? (Nennen Sie zwei.)
\item Was sagen die Lehrer über das Internet ?
\end{enumerate}

\textbf{Fragen zur Sprache :}
\begin{enumerate}
\item Trouvez dans le texte deux adjectifs et mettez-les au comparatif et au superlatif.
\item Relevez dans le texte un verbe fort au présent et conjuguez-le au \all{Perfekt} à toutes les personnes du singulier.
\item Transformez : \all{Das Internet ist ein gutes Werkzeug.} $\rightarrow$ \all{Die Lehrer sagen, dass \ldots}
\end{enumerate}
\end{exercice}

\begin{exercice}[Thème et version (type Bac)]
\textbf{A. Version.} Traduisez en français le paragraphe suivant.

\begin{textebox}[Die Gefahren der sozialen Netzwerke]
Viele Jugendliche verbringen jeden Tag viele Stunden in den sozialen Netzwerken. Das ist gefährlich: Sie schlafen wenig, lernen nicht genug und sprechen nicht mehr mit ihrer Familie. Manche zeigen auch zu viele persönliche Fotos und Informationen. Experten sagen, dass die Jugendlichen vorsichtiger sein müssen. Das Handy ist nützlich, aber es darf das Leben nicht kontrollieren.
\end{textebox}

\textbf{B. Thème.} Traduisez en allemand.
\begin{enumerate}
\item Je pense que le smartphone est plus cher que le cybercafé.
\item Hier, mon frère est allé au cybercafé et il a imprimé ses devoirs.
\item Les jeunes Camerounais aiment beaucoup discuter en ligne.
\item Ma sœur a trouvé beaucoup d'informations sur Internet.
\item Je crois qu'Internet est utile, mais il a aussi des inconvénients.
\end{enumerate}
\end{exercice}

\begin{exercice}[Rédaction guidée (type Bac)]
\textbf{Sujet :} \all{Die Vorteile und Nachteile des Smartphones für Jugendliche in Kamerun.}

Rédigez un court texte en allemand (environ 10 lignes, 80 à 100 mots) en suivant le plan ci-dessous. Employez au moins deux comparatifs, une subordonnée en \all{dass} et deux verbes forts au \all{Perfekt}.
\begin{enumerate}
\item \textbf{Einleitung :} Viele Jugendliche in Kamerun haben heute ein Smartphone. (Présentez le sujet en une ou deux phrases.)
\item \textbf{Vorteile :} Was kann man mit dem Smartphone machen ? (telefonieren, lernen, Fotos machen, im Internet surfen \ldots)
\item \textbf{Nachteile :} Was ist negativ ? (teuer, zu viel chatten, weniger lernen \ldots)
\item \textbf{Schluss :} Ihre Meinung : \all{Ich denke, dass \ldots}
\end{enumerate}
\end{exercice}

\newpage

\coursec{Corrigés des exercices}

\begin{corrige}[Exercice 1 — QCM : vocabulaire des médias]
\begin{enumerate}
\item \textbf{b) le cybercafé.} \all{Das Internetcafé} (pluriel : \all{die Internetcafés}) désigne le lieu où l'on paie pour utiliser Internet.
\item \textbf{c) l'avantage.} \all{Der Vorteil} (pluriel : \all{die Vorteile}). Son contraire est \all{der Nachteil} (pluriel : \all{die Nachteile}).
\item \textbf{b) discuter en ligne.} \all{chatten} est un verbe faible (régulier) emprunté à l'anglais : \all{ich chatte, er chattet, er hat gechattet}.
\item \textbf{a) La connexion est mauvaise.} \all{schlecht} = mauvais ; \all{schnell} = rapide ; \all{gratis} = gratuit.
\item \textbf{a) die Smartphones.} Pluriel en \all{-s}, typique des mots étrangers récents : \all{das Smartphone, die Smartphones} ; de même \all{das Handy, die Handys}.
\item \textbf{b) der Vorteil.} \all{Der Nachteil} (l'inconvénient) a pour contraire \all{der Vorteil} (l'avantage).
\end{enumerate}
\end{corrige}

\begin{corrige}[Exercice 2 — Vrai ou faux ? Justifiez]
\begin{enumerate}
\item \textbf{Falsch.} \all{Man kann mit dem Smartphone nicht nur telefonieren, sondern auch Nachrichten schreiben, Fotos machen und im Internet surfen.} (On peut aussi écrire des messages, prendre des photos et surfer sur Internet.) Noter la tournure de correction : \all{nicht nur \ldots\ sondern auch} (non seulement \ldots\ mais aussi).
\item \textbf{Richtig.} \all{Im Internetcafé kann man Dokumente drucken.} (Au cybercafé, on peut imprimer des documents.)
\item \textbf{Falsch.} \all{Das Smartphone ist teuer, das Internetcafé ist billiger.} (Le smartphone est cher, le cybercafé est moins cher.) Attention : le mot \all{immer} (toujours) rend la phrase fausse, car le texte ne dit pas que le cybercafé est \emph{toujours} moins cher.
\item \textbf{Richtig.} \all{Viele Jugendliche in Kamerun benutzen das Internet zum Lernen und für die Hausaufgaben.} (Beaucoup de jeunes au Cameroun utilisent Internet pour apprendre et pour les devoirs.)
\end{enumerate}
\end{corrige}

\begin{corrige}[Exercice 3 — Articles et pluriels]
\begin{center}
\begin{tabular}{|l|l|l|}
\hline
\rowcolor{popblueL} Nom & Article & Pluriel \\
\hline
Smartphone & das & die Smartphones \\
\hline
Internetcafé & das & die Internetcafés \\
\hline
Vorteil & der & die Vorteile \\
\hline
Nachteil & der & die Nachteile \\
\hline
Verbindung & die & die Verbindungen \\
\hline
Computer & der & die Computer \\
\hline
\end{tabular}
\end{center}
Points de vigilance :
\begin{itemize}
\item Les noms en \all{-ung} sont toujours féminins (\all{die Verbindung}) et prennent \all{-en} au pluriel : \all{die Verbindung, die Verbindungen}.
\item Les mots anglais récents prennent souvent \all{-s} au pluriel : \all{das Smartphone, die Smartphones} ; \all{das Internetcafé, die Internetcafés}.
\item \all{Der Computer} est invariable au pluriel : \all{der Computer, die Computer} (comme \all{der Lehrer, die Lehrer}).
\item Au pluriel, l'article est toujours \all{die}, quel que soit le genre du singulier.
\end{itemize}
\end{corrige}

\begin{corrige}[Exercice 4 — Conjugaison : le parfait des verbes forts]
\begin{enumerate}
\item Gestern \textbf{habe} ich eine Nachricht \textbf{geschrieben}. (Hier, j'ai écrit un message.)
\item Wir \textbf{haben} einen Text im Internet \textbf{gelesen}. (Nous avons lu un texte sur Internet.)
\item Paul \textbf{ist} gestern ins Internetcafé \textbf{gegangen}. (Paul est allé hier au cybercafé.)
\item Marie \textbf{hat} viele Informationen \textbf{gefunden}. (Marie a trouvé beaucoup d'informations.)
\item Ich \textbf{habe} mein Handy leider \textbf{verloren}. (J'ai malheureusement perdu mon portable.)
\item Die Schüler \textbf{haben} den Bus in die Stadt \textbf{genommen}. (Les élèves ont pris le bus pour la ville.)
\end{enumerate}
Justifications :
\begin{itemize}
\item \all{schreiben -- schrieb -- hat geschrieben} ; \all{lesen -- las -- hat gelesen} ; \all{finden -- fand -- hat gefunden} ; \all{verlieren -- verlor -- hat verloren} ; \all{nehmen -- nahm -- hat genommen} : verbes forts conjugués avec l'auxiliaire \all{haben}.
\item \all{gehen -- ging -- ist gegangen} : verbe de déplacement, il se conjugue avec l'auxiliaire \all{sein}.
\end{itemize}
Attention à la place du participe passé : il se place \textbf{à la fin de la phrase}, tandis que l'auxiliaire conjugué occupe la deuxième position. Ne dites jamais \all{ich habe geschrieben eine Nachricht} : le participe doit rester en fin de phrase.
\end{corrige}

\begin{corrige}[Exercice 5 — Comparatif et superlatif]
\begin{center}
\begin{tabular}{|l|l|l|}
\hline
\rowcolor{popblueL} Positif & Komparativ & Superlativ \\
\hline
schnell & schneller & am schnellsten \\
\hline
teuer & teurer & am teuersten \\
\hline
gut & besser & am besten \\
\hline
groß & größer & am größten \\
\hline
gern & lieber & am liebsten \\
\hline
praktisch & praktischer & am praktischsten \\
\hline
\end{tabular}
\end{center}
Phrases possibles (modèles) :
\begin{itemize}
\item \all{teuer} : \all{Das Smartphone ist teurer als das Internetcafé. Das neue Handy ist am teuersten.} (Le smartphone est plus cher que le cybercafé. Le nouveau portable est le plus cher.)
\item \all{gut} : \all{Die Verbindung im Internetcafé ist besser als zu Hause. Am besten ist die Verbindung am Morgen.} (La connexion au cybercafé est meilleure qu'à la maison. La meilleure connexion est le matin.)
\item \all{groß} : \all{Das Internetcafé in Douala ist größer als das Internetcafé in meinem Viertel. Es ist am größten in der Stadt.} (Le cybercafé de Douala est plus grand que celui de mon quartier. Il est le plus grand de la ville.)
\item \all{gern} : \all{Ich chatte lieber mit meinen Freunden. Am liebsten lerne ich im Internet.} (Je préfère discuter en ligne avec mes amis. Ce que je préfère, c'est apprendre sur Internet.)
\item \all{praktisch} : \all{Das Smartphone ist praktischer als der Computer. Es ist am praktischsten für unterwegs.} (Le smartphone est plus pratique que l'ordinateur. Il est le plus pratique en déplacement.)
\end{itemize}
Points de vigilance :
\begin{itemize}
\item \all{gut -- besser -- am besten} et \all{gern -- lieber -- am liebsten} sont irréguliers : ils s'apprennent par cœur.
\item Les adjectifs en \all{-er} perdent le \all{e} au comparatif : \all{teuer -- teurer} (et non \all{teuerer}).
\item Les adjectifs monosyllabiques à voyelle \all{a, o, u} prennent souvent l'Umlaut : \all{groß -- größer -- am größten}.
\item Le superlatif se construit avec \all{am \ldots\ -sten} : \all{am schnellsten}.
\item Le second terme de la comparaison est introduit par \all{als} (et non par \all{wie}) : \all{teurer als das Internetcafé}.
\end{itemize}
\end{corrige}

\begin{corrige}[Exercice 6 — Les subordonnées en dass]
\begin{enumerate}
\item \all{Ich denke, dass das Smartphone teuer ist.} (Je pense que le smartphone est cher.)
\item \all{Ich glaube, dass das Internetcafé billiger ist.} (Je crois que le cybercafé est moins cher.)
\item \all{Ich denke, dass viele Jugendliche zu viel chatten.} (Je pense que beaucoup de jeunes discutent trop en ligne.)
\item \all{Ich glaube, dass die Verbindung heute schlecht ist.} (Je crois que la connexion est mauvaise aujourd'hui.)
\item \all{Ich denke, dass das Internet den Schülern beim Lernen hilft.} (Je pense qu'Internet aide les élèves à apprendre.) Noter : \all{helfen} se construit avec le datif : \all{den Schülern}.
\item \all{Ich glaube, dass mein Bruder sein Handy verloren hat.} (Je crois que mon frère a perdu son portable.)
\end{enumerate}
Règle appliquée : dans la subordonnée introduite par \all{dass}, le verbe conjugué se place \textbf{à la fin}. Au \all{Perfekt} (phrase 6), l'ordre est : participe passé puis auxiliaire (\all{verloren hat}), et non \all{hat verloren}.
\end{corrige}

\begin{corrige}[Exercice 7 — Texte à trous]
Mein Bruder hat ein neues \textbf{Smartphone}. Er \textbf{chattet} jeden Tag mit seinen Freunden. Aber das Handy ist sehr \textbf{teuer}. Ich habe kein Handy. Ich gehe oft ins \textbf{Internetcafé} in der Stadt. Dort kann ich im Internet surfen und meine Hausaufgaben \textbf{drucken}. Die \textbf{Verbindung} ist dort schnell. Das Internet hat viele \textbf{Vorteile} : Man lernt viel und findet Informationen. Aber es hat auch \textbf{Nachteile} : Manche Jugendliche spielen zu viel und lernen nicht mehr.

Traduction : Mon frère a un nouveau smartphone. Il discute en ligne tous les jours avec ses amis. Mais le portable est très cher. Je n'ai pas de portable. Je vais souvent au cybercafé en ville. Là-bas, je peux surfer sur Internet et imprimer mes devoirs. La connexion y est rapide. Internet a beaucoup d'avantages : on apprend beaucoup et on trouve des informations. Mais il a aussi des inconvénients : certains jeunes jouent trop et n'apprennent plus.

Points de vigilance :
\begin{itemize}
\item \all{chattet} : 3\up{e} personne du singulier du verbe faible \all{chatten}.
\item \all{Vorteile} et \all{Nachteile} au pluriel après \all{viele}.
\item \all{ins Internetcafé} = \all{in das Internetcafé} : accusatif après \all{in} quand il y a idée de mouvement (on va \emph{vers} le cybercafé) ; on aurait le datif pour le lieu où l'on se trouve : \all{im Internetcafé} (= \all{in dem Internetcafé}).
\end{itemize}
\end{corrige}

\begin{corrige}[Exercice 8 — Appariement et traduction]
\textbf{A. Appariement :} 1 -- b ; 2 -- a ; 3 -- d ; 4 -- c ; 5 -- e.

\textbf{B. Traductions :}
\begin{enumerate}
\item \all{Ich denke, dass das Internetcafé billiger ist als das Smartphone.} (Verbe conjugué \all{ist} en fin de subordonnée ; \all{billiger als} = moins cher que.)
\item \all{Das Smartphone ist praktischer als der Computer.} (Comparatif + \all{als}.)
\item \all{Gestern sind wir ins Internetcafé gegangen und haben unsere Hausaufgaben gedruckt.} (\all{gehen} au \all{Perfekt} avec \all{sein} : \all{sind \ldots\ gegangen} ; \all{drucken} est un verbe faible : \all{hat gedruckt}.)
\item \all{Meine Schwester hat eine Nachricht geschrieben und Fotos geschickt.} (\all{schreiben -- hat geschrieben} ; \all{schicken -- hat geschickt}.)
\item \all{Ich glaube, dass das Internet viele Vorteile für die Schüler hat.} (\all{das Internet} est sujet : le verbe \all{hat} est au singulier et rejeté à la fin de la subordonnée.)
\end{enumerate}
\end{corrige}

\begin{corrige}[Exercice 9 — Texte et questions de compréhension (type Bac)]
\textbf{Fragen zum Text — réponses modèles :}
\begin{enumerate}
\item \all{Mit dem Smartphone können die Jugendlichen telefonieren, Nachrichten schreiben und Fotos machen.} (Ils peuvent aussi surfer sur Internet.)
\item \all{Viele Schüler gehen ins Internetcafé, weil der Preis dort niedrig ist und weil nicht alle Familien ein Smartphone kaufen können.} (Parce que le prix y est bas et que toutes les familles ne peuvent pas acheter un smartphone.) Noter : après \all{weil}, le verbe conjugué va à la fin (\all{ist}, \all{können}).
\item \all{Vor den Prüfungen suchen die Schüler Informationen für ihre Hausaufgaben und lernen mit ihren Freunden zusammen.} (Avant les examens, ils cherchent des informations pour leurs devoirs et apprennent ensemble avec leurs amis.)
\item \all{Der Text nennt zwei Nachteile: Manche Jugendliche chatten stundenlang, spielen nur und lernen nicht mehr. Außerdem ist die Verbindung manchmal langsam.} (Certains jeunes bavardent des heures en ligne, ne font que jouer et n'apprennent plus ; de plus, la connexion est parfois lente.)
\item \all{Die Lehrer sagen, dass das Internet ein gutes Werkzeug ist, aber dass man es richtig benutzen muss.} (Les professeurs disent qu'Internet est un bon outil, mais qu'il faut l'utiliser correctement.)
\end{enumerate}
\textbf{Fragen zur Sprache :}
\begin{enumerate}
\item Exemples : \all{praktisch -- praktischer -- am praktischsten} ; \all{teuer -- teurer -- am teuersten} ; \all{niedrig -- niedriger -- am niedrigsten} ; \all{langsam -- langsamer -- am langsamsten} ; \all{gut -- besser -- am besten}.
\item Exemple avec \all{gehen} (verbe fort du texte : \all{viele Schüler gehen}) : \all{ich bin gegangen, du bist gegangen, er/sie/es ist gegangen}. Autre possibilité : \all{schreiben} : \all{ich habe geschrieben, du hast geschrieben, er hat geschrieben}. Attention : \all{suchen} est un verbe faible (\all{ich habe gesucht}) ; il ne convient pas si l'on demande un verbe fort.
\item \all{Die Lehrer sagen, dass das Internet ein gutes Werkzeug ist.} (Le verbe conjugué \all{ist} est rejeté à la fin de la subordonnée.)
\end{enumerate}
Barème indicatif (sur 20) : questions de compréhension 10 points (2 par question : 1 pour le sens, 1 pour la langue) ; questions de langue 6 points (2 par item) ; présentation et orthographe 4 points.

Points de vigilance : répondre par des \textbf{phrases complètes} ; ne pas recopier de longs passages sans les adapter ; vérifier la place du verbe dans les subordonnées ; majuscule à tous les noms.
\end{corrige}

\begin{corrige}[Exercice 10 — Thème et version (type Bac)]
\textbf{A. Version — traduction modèle :}

Beaucoup de jeunes passent chaque jour de nombreuses heures sur les réseaux sociaux. C'est dangereux : ils dorment peu, n'apprennent pas assez et ne parlent plus avec leur famille. Certains montrent aussi trop de photos et d'informations personnelles. Les experts disent que les jeunes doivent être plus prudents. Le portable est utile, mais il ne doit pas contrôler la vie.

Glossaire : \all{verbringen} = passer (du temps) ; \all{gefährlich} = dangereux ; \all{vorsichtig} = prudent ; \all{nützlich} = utile ; \all{kontrollieren} = contrôler.

\textbf{B. Thème — traductions modèles :}
\begin{enumerate}
\item \all{Ich denke, dass das Smartphone teurer ist als das Internetcafé.}
\item \all{Gestern ist mein Bruder ins Internetcafé gegangen und hat seine Hausaufgaben gedruckt.}
\item \all{Die jungen Kameruner chatten sehr gern.} (ou : \all{Die Jugendlichen in Kamerun chatten sehr gern.})
\item \all{Meine Schwester hat viele Informationen im Internet gefunden.}
\item \all{Ich glaube, dass das Internet nützlich ist, aber dass es auch Nachteile hat.}
\end{enumerate}
Barème indicatif (sur 20) : version 8 points ; thème 10 points (2 par phrase : vocabulaire 1, grammaire 1) ; présentation 2 points.

Points de vigilance : auxiliaire \all{sein} avec \all{gehen} (phrase 2 : \all{ist \ldots\ gegangen}) ; participe de \all{finden} : \all{gefunden} ; verbe conjugué en fin de subordonnée après \all{dass} ; \all{im Internet} = \all{in dem Internet} (datif, lieu où l'on se trouve).
\end{corrige}

\begin{corrige}[Exercice 11 — Rédaction guidée (type Bac)]
\textbf{Proposition de rédaction modèle :}

\begin{textebox}[Die Vorteile und Nachteile des Smartphones]
Heute haben viele Jugendliche in Kamerun ein Smartphone. Das Handy ist ein Teil ihres Lebens.

Das Smartphone hat viele Vorteile. Man kann telefonieren, Nachrichten schreiben und im Internet surfen. Es ist praktischer als der Computer, weil man es überall benutzen kann. Viele Schüler lernen auch mit dem Handy: Sie suchen Informationen für die Schule. Letzte Woche habe ich einen Text für meine Hausaufgaben im Internet gefunden.

Aber das Smartphone hat auch Nachteile. Es ist teurer als das Internetcafé, und nicht alle Familien können es kaufen. Manche Jugendliche chatten zu viel und lernen weniger. Mein Freund hat gestern drei Stunden nur gespielt!

Ich denke, dass das Smartphone ein gutes Werkzeug ist, aber man muss es klug benutzen: zum Lernen und nicht nur zum Spielen.
\end{textebox}

(Environ 110 mots ; à adapter à 80--100 mots en supprimant une phrase.)

Traduction : Aujourd'hui, beaucoup de jeunes au Cameroun ont un smartphone. Le portable fait partie de leur vie. Le smartphone a beaucoup d'avantages. On peut téléphoner, écrire des messages et surfer sur Internet. Il est plus pratique que l'ordinateur, car on peut l'utiliser partout. Beaucoup d'élèves apprennent aussi avec le portable : ils cherchent des informations pour l'école. La semaine dernière, j'ai trouvé un texte pour mes devoirs sur Internet. Mais le smartphone a aussi des inconvénients. Il est plus cher que le cybercafé, et toutes les familles ne peuvent pas l'acheter. Certains jeunes discutent trop en ligne et apprennent moins. Mon ami n'a fait que jouer hier pendant trois heures ! Je pense que le smartphone est un bon outil, mais il faut l'utiliser intelligemment : pour apprendre et pas seulement pour jouer.

Barème indicatif (sur 20) : respect du plan et du sujet 5 points ; richesse du vocabulaire thématique 4 points ; correction grammaticale (comparatifs, \all{dass}, \all{Perfekt}, place du verbe) 7 points ; orthographe et majuscules 2 points ; présentation 2 points.

Points de vigilance :
\begin{itemize}
\item Vérifier la présence d'au moins deux comparatifs (\all{praktischer als, teurer als}), d'une subordonnée en \all{dass} avec le verbe conjugué à la fin, et de deux verbes forts au \all{Perfekt} (\all{habe gefunden, hat gespielt}).
\item Majuscule à tous les noms ; Umlaute obligatoires (\all{Jugendliche, Werkzeug, praktischer}).
\item Éviter les calques du français : « pour apprendre » se traduit par \all{zum Lernen}, et non par une construction avec \all{für} ; « sur Internet » se dit \all{im Internet}.
\end{itemize}
\end{corrige}

\newpage

\coursec{Fiche bilan}

\begin{popfigure}
\begin{tikzpicture}[
  mindmap,
  grow cyclic,
  every node/.style={concept, align=center, font=\small},
  root concept/.append style={concept color=popblue, font=\bfseries\small, text=white},
  level 1/.append style={level distance=3.2cm, sibling angle=60},
  level 1 concept/.append style={font=\footnotesize}
]
\node[root concept]{Smartphones und Internetcafés}
  child[concept color=poporange]{ node[align=center]{Wortschatz\\ \all{das Smartphone}\\ \all{der Vorteil}\\ \all{chatten}} }
  child[concept color=popgreen]{ node[align=center]{Komparativ\\ \all{schneller}\\ \all{besser}\\ \all{am besten}} }
  child[concept color=poppurple]{ node[align=center]{Nebensatz\\ mit \all{dass}\\ Verb am Ende} }
  child[concept color=popgold]{ node[align=center]{Perfekt\\ der starken\\ Verben} }
  child[concept color=poppink]{ node[align=center]{Textverständnis\\ global + detailliert} }
  child[concept color=popblue!70]{ node[align=center]{Methode\\ Kommentar\\ + kurzer Text} };
\end{tikzpicture}
\legende{Carte mentale de la leçon : les six blocs à maîtriser.}
\end{popfigure}

\begin{bilan}
\textbf{1. Lexique clé (à connaître avec article et pluriel)}
\begin{itemize}\setlength{\itemsep}{1pt}
  \item \all{das Smartphone, -s} : le smartphone ; \all{das Handy, -s} : le portable
  \item \all{das Internetcafé, -s} : le cybercafé ; \all{das Internet} : l'internet
  \item \all{der Vorteil, -e} : l'avantage ; \all{der Nachteil, -e} : l'inconvénient
  \item \all{die Verbindung, -en} : la connexion ; \all{das Netz, -e} : le réseau
  \item \all{chatten} : tchatter ; \all{telefonieren} : téléphoner ; \all{surfen} : surfer
  \item \all{die Nachricht, -en} : le message ; \all{schicken / senden} : envoyer
  \item \all{billig / teuer} : bon marché / cher ; \all{schnell / langsam} : rapide / lent
\end{itemize}

\textbf{2. Grammaire à connaître par cœur}
\begin{center}
\begin{tabularx}{\linewidth}{|l|X|X|}
\hline
\rowcolor{popblueL} Point & Règle & Exemple \\
\hline
Comparatif & adjectif + \all{-er} (+ \all{als}) & \all{Das Smartphone ist schneller als der Brief.} \\
\hline
Superlatif & \all{am} + adjectif + \all{-sten} & \all{Das Internetcafé ist am billigsten.} \\
\hline
Irréguliers & \all{gut $\to$ besser $\to$ am besten} ; \all{gern $\to$ lieber $\to$ am liebsten} ; \all{viel $\to$ mehr $\to$ am meisten} & \all{Ich chatte lieber.} \\
\hline
Subordonnée \all{dass} & verbe conjugué \textbf{à la fin} & \all{Ich denke, dass das Handy praktisch ist.} \\
\hline
Perfekt (forts) & \all{haben/sein} + participe en \all{ge-...-en} & \all{gehen $\to$ ist gegangen} ; \all{lesen $\to$ hat gelesen} \\
\hline
\end{tabularx}
\end{center}

\textbf{3. Méthodes express}
\begin{itemize}\setlength{\itemsep}{1pt}
  \item Compréhension globale : lire le titre, repérer le thème, puis lire le texte une première fois sans dictionnaire.
  \item Compréhension détaillée : souligner les mots-clés (\all{Vorteil, Nachteil, aber, weil}) et répondre par des phrases complètes.
  \item Commentaire : 1) présenter le texte, 2) dégager l'idée principale, 3) relever avantages et inconvénients, 4) donner son avis avec \all{Ich meine, dass ...}.
  \item Production guidée : 4 à 6 phrases simples, verbe en position 2, un comparatif et une subordonnée en \all{dass}.
\end{itemize}
\end{bilan}

\begin{aretenir}[Checklist avant le Bac]
\begin{itemize}\setlength{\itemsep}{1pt}
  \item[$\square$] Je connais l'article et le pluriel de \all{das Smartphone, der Vorteil, der Nachteil, die Verbindung, das Internetcafé}.
  \item[$\square$] Je sais former le comparatif : adjectif + \all{-er} + \all{als}.
  \item[$\square$] Je sais former le superlatif : \all{am} + adjectif + \all{-sten}.
  \item[$\square$] Je connais par cœur \all{gut -- besser -- am besten} et \all{gern -- lieber -- am liebsten}.
  \item[$\square$] Je place le verbe conjugué à la fin dans une subordonnée en \all{dass}.
  \item[$\square$] Je mets une virgule avant \all{dass}.
  \item[$\square$] Je connais le parfait des verbes forts fréquents : \all{gehen, kommen, lesen, schreiben, sehen, nehmen}.
  \item[$\square$] Je choisis correctement l'auxiliaire \all{sein} pour les verbes de mouvement.
  \item[$\square$] Je mets la majuscule à tous les noms allemands.
  \item[$\square$] Je sais rédiger 5 phrases sur les avantages et les inconvénients du smartphone.
  \item[$\square$] Je sais exprimer mon avis : \all{Ich denke / Ich meine, dass ...}.
  \item[$\square$] Je relis ma production : place du verbe, terminaisons, orthographe.
\end{itemize}
\end{aretenir}

\findecours

\end{document}
